%% LyX 2.3.7 created this file.  For more info, see http://www.lyx.org/.
%% Do not edit unless you really know what you are doing.
\documentclass[oneside,american]{amsart}
\usepackage[T1]{fontenc}
\usepackage[latin9]{inputenc}
\usepackage{color}
\definecolor{note_fontcolor}{rgb}{0.800781, 0.800781, 0.800781}
\usepackage{babel}
\usepackage{refstyle}
\usepackage{enumitem}
\usepackage{amstext}
\usepackage{amsthm}
\usepackage{amssymb}
\usepackage{setspace}
\PassOptionsToPackage{normalem}{ulem}
\usepackage{ulem}
\usepackage[unicode=true,pdfusetitle,
 bookmarks=true,bookmarksnumbered=false,bookmarksopen=false,
 breaklinks=false,pdfborder={0 0 1},backref=false,colorlinks=false]
 {hyperref}

\makeatletter

%%%%%%%%%%%%%%%%%%%%%%%%%%%%%% LyX specific LaTeX commands.

\AtBeginDocument{\providecommand\thmref[1]{\ref{thm:#1}}}
%% Because html converters don't know tabularnewline
\providecommand{\tabularnewline}{\\}
%% The greyedout annotation environment
\newenvironment{insight}
  {\textcolor{note_fontcolor}\bgroup\ignorespaces}
  {\ignorespacesafterend\egroup}
\RS@ifundefined{subsecref}
  {\newref{subsec}{name = \RSsectxt}}
  {}
\RS@ifundefined{thmref}
  {\def\RSthmtxt{theorem~}\newref{thm}{name = \RSthmtxt}}
  {}
\RS@ifundefined{lemref}
  {\def\RSlemtxt{lemma~}\newref{lem}{name = \RSlemtxt}}
  {}


%%%%%%%%%%%%%%%%%%%%%%%%%%%%%% Textclass specific LaTeX commands.
\numberwithin{equation}{section}
\numberwithin{figure}{section}
\theoremstyle{plain}
\newtheorem{thm}{\protect\theoremname}
\theoremstyle{definition}
\newtheorem{definition}[thm]{\protect\definitionname}
\theoremstyle{definition}
\newtheorem{example}[thm]{\protect\examplename}
\theoremstyle{remark}
\newtheorem{remark}[thm]{\protect\remarkname}
\theoremstyle{remark}
\newtheorem{claim}[thm]{\protect\claimname}
\theoremstyle{definition}
\newtheorem{xca}[thm]{\protect\exercisename}
\theoremstyle{plain}
\newtheorem{corollary}[thm]{\protect\corollaryname}
\theoremstyle{plain}
\newtheorem{proposition}[thm]{\protect\propositionname}
\newlist{casenv}{enumerate}{4}
\setlist[casenv]{leftmargin=*,align=left,widest={iiii}}
\setlist[casenv,1]{label={{\itshape\ \casename} \arabic*.},ref=\arabic*}
\setlist[casenv,2]{label={{\itshape\ \casename} \roman*.},ref=\roman*}
\setlist[casenv,3]{label={{\itshape\ \casename\ \alph*.}},ref=\alph*}
\setlist[casenv,4]{label={{\itshape\ \casename} \arabic*.},ref=\arabic*}
\theoremstyle{plain}
\newtheorem{fact}[thm]{\protect\factname}
\theoremstyle{plain}
\newtheorem{lemma}[thm]{\protect\lemmaname}
\theoremstyle{definition}
\newtheorem{sol}[thm]{\protect\solutionname}
\newenvironment{lyxlist}[1]
	{\begin{list}{}
		{\settowidth{\labelwidth}{#1}
		 \setlength{\leftmargin}{\labelwidth}
		 \addtolength{\leftmargin}{\labelsep}
		 \renewcommand{\makelabel}[1]{##1\hfil}}}
	{\end{list}}
\theoremstyle{remark}
\newtheorem*{claim*}{\protect\claimname}

%%%%%%%%%%%%%%%%%%%%%%%%%%%%%% User specified LaTeX commands.
\usepackage{amsmath}
\DeclareMathOperator*{\argmax}{arg\,max}
\DeclareMathOperator*{\argmin}{arg\,min}
\usepackage{bbm}

\makeatother

\providecommand{\casename}{Case}
\providecommand{\claimname}{Claim}
\providecommand{\corollaryname}{Corollary}
\providecommand{\definitionname}{Definition}
\providecommand{\examplename}{Example}
\providecommand{\exercisename}{Exercise}
\providecommand{\factname}{Fact}
\providecommand{\lemmaname}{Lemma}
\providecommand{\propositionname}{Proposition}
\providecommand{\remarkname}{Remark}
\providecommand{\solutionname}{Solution}
\providecommand{\theoremname}{Theorem}

\begin{document}
\begin{titlepage} 
\title{Mathematical Tools for Computer Science}
\author{Jonathan Weiss}
\maketitle
\begin{spacing}{0}
\noindent \begin{center}
Please send comments to: \href{mailto:Jonathan.weiss1@mail.huji.ac.il}{Jonathan.weiss1@mail.huji.ac.il}
\par\end{center}
\end{spacing}

\end{titlepage}

\pagebreak{}

\global\long\def\ball#1{B_{\left|\left|\cdot\right|\right|_{#1}}}%
\global\long\def\C{\mathbb{C}}%
\global\long\def\bin#1#2{\binom{#1}{#2}}%
\global\long\def\p#1{\Pr\left[#1\right]}%
\global\long\def\limn#1{\lim\limits _{n\to\infty}\left(#1\right)}%
\global\long\def\one{\mathbbm{1}}%
\global\long\def\cov#1#2{COV\left[#1,#2\right]}%
\global\long\def\E#1{\mathbb{E}\left[#1\right]}%
\global\long\def\var#1{var\left[#1\right]}%
\global\long\def\R{\mathbb{R}}%
\global\long\def\N{\mathbb{N}}%
\global\long\def\then{\Rightarrow}%
\global\long\def\norm#1{\left\Vert #1\right\Vert }%
\global\long\def\Z{\mathbb{Z}}%
\tableofcontents

\pagebreak{}

\part{probability}

\section{Probability recap}
\begin{definition}
Probability space (merchav midgam)

a set $\Omega$ and a probability measure $p_{r}$. 
\end{definition}


\subsubsection*{Simplest case:}

Finite discrete probability space $\left|\Omega\right|<\infty$ individual
element $\omega\in\Omega$ we associate non negative number $p_{r}\left(\omega\right)\ge0$
s.t $\sum\limits _{\omega\in\Omega}p_{r}\left(\omega\right)=1$.
\begin{definition}
Events (meora)

We call $A\subseteq\Omega$ an event. The probability of an event
is: $p_{r}\left(A\right)=\sum\limits _{\omega\in A}p_{r}\left(\omega\right)$. 
\end{definition}


\subsubsection*{Obvious extension: }

$\left|\Omega\right|=\aleph_{0}$ for example $\Omega=\mathbb{N}$.
$\Omega=\left\{ x_{1},x_{2},\dots\right\} $ and $p_{1}.p_{2}\dots\ge0,\sum p_{i}=1$.

\subsubsection*{Last case:}

$\Omega=\left[0,1\right],0<\alpha<\beta<1,A=\left[\alpha,\beta\right]$.
$p_{r}\left(A\right)=\beta-\alpha$
\begin{definition}
Independent events (meoraot bilti tloolim)

Let $\left(\Omega,p_{r}\right)$ be a probability space\c{ }and let
$A,b\subseteq\Omega$ be two events. We say that events $A$ and $B$
are independent if $p_{r}\left(A\cap B\right)=p_{r}\left(A\right)\cdot p_{r}\left(B\right)$
\end{definition}

\begin{example}
intuition for random variables

a fair dice and a fair coin are being tossed. $\Omega=\left[6\right]\times\left\{ \mathbb{H},\mathbb{T}\right\} $,
$\forall\omega\in\Omega,p_{r}\left(\omega\right)=\frac{1}{12}$ (uniform
distribution). We treat the dice and the coin as independent. Thus
$p_{r}\left(\text{dice is even}\right)=p_{r}\left(A\right)=\frac{1}{2}$,
$p_{r}\left(\text{coin is on }H\right)=p_{r}\left(B\right)=\frac{1}{2}$.
$p_{r}\left(\text{dice is even}\cap\text{coin is on }H\right)=p_{r}\left(A\cap B\right)=p_{r}\left(A\right)\cdot p_{r}\left(B\right)=\frac{1}{4}$.
Now we can inspect all possible cases of $\left[6\right]\times\left\{ \mathbb{H},\mathbb{T}\right\} $
and find the event. $A\cap B=\left\{ \left(2,H\right),\left(4,H\right),\left(6,H\right)\right\} $
thus $p_{r}\left(A\cap B\right)=\frac{3}{\left|\Omega\right|}=\frac{3}{12}=\frac{1}{4}$
\end{example}

\begin{definition}
Conditioning (hatnaya)
\begin{definition}
conditioning on the event $A$ is defined as follows: $p_{r}\left(B|A\right)=\frac{p_{r}\left(A\cap B\right)}{p_{r}\left(A\right)}$.
In case of independence between the events it is simple to deduce
that $p_{r}\left(B|A\right)=p_{r}\left(B\right)$: 
\[
p_{r}\left(B|A\right)=\frac{p_{r}\left(A\cap B\right)}{p_{r}\left(A\right)}=\frac{p_{r}\left(A\right)\cdot p_{r}\left(B\right)}{p_{r}\left(A\right)}=p_{r}\left(B\right)
\]
.
\end{definition}

\end{definition}

\begin{remark}
We can think of conditioning as looking at event $A$ as our probability
space thus we normalize according to it. $A$ happened what is the
probability that $B$ will happen too? so we do a conjunction between
the two and normalize by A to get the new probability.
\end{remark}

\begin{definition}
Random Variable: Let $\left(\Omega,p_{r}\right)$ be a probability
space and $X:\Omega\to\mathbb{R}$ is called a real random variable.
In general it is very typical to not care about where $X$ maps $\omega\in\Omega$.
In most cases we care about the event that is tied to the random variable. 
\end{definition}

\begin{example}
Random Variable: 

$X:\Omega\to\mathbb{R}$, $\left[\alpha,\beta\right]\subseteq\mathbb{R}$,
let $A\subseteq\Omega$ be the event $\left\{ \omega\in\Omega|X\left(\omega\right)\in\left[\alpha,\beta\right]\right\} $.
Using random variable we this way allows us to inquire important information
about real life events. Each $\omega\in\Omega$ in this example can
be a person and $X$ can map from a person to it's salary. 
\end{example}

\begin{definition}
Random variable independency 

Let $X,Y:\Omega\to S$ be Random Variables (R.V) we say that $X,Y$
are independent if $\forall r,\sigma\in S$ the events $\left\{ \omega|X\left(\omega\right)=r\right\} ,\left\{ \omega|Y\left(\omega\right)=\sigma\right\} $
are independent events.
\end{definition}

\begin{definition}
Indicator 

$X:\Omega\to\left\{ 0,1\right\} $ is called an ``indicator random
variable''.
\end{definition}

When we are working with random variables we might face questions
like the following: What is the average grade of $X$
\begin{definition}
Expectation (Tochelet)

Denote $\mathbb{E}\left[x\right]=\sum\limits _{\omega\in\Omega}p_{r}\left(\omega\right)\cdot X\left(\omega\right)$
the expectation of $X$.
\end{definition}

\begin{claim}
Linearity of expectation (with no proof)

$\forall X,Y:\Omega\to\mathbb{R},\forall a,b\in\mathbb{R},\mathbb{E}\left[aX+bY\right]=a\mathbb{E}\left[X\right]+b\mathbb{E}\left[Y\right]$
\end{claim}

Aside from the average (expectation) of a R.V we might also wish to
know how \uline{concentrated} it is. how can we find how far are
all the values $p_{1}a_{1}+\dots+p_{n}a_{n}=\mathbb{E}\left[X\right]$
from the average of $X$?

\paragraph*{First try}

$\mu:=\mathbb{E}\left[X\right]=\sum p_{i}a_{i}$ , now to calculate
how far each value (with it's appropriate size factor):$\sum p_{i}\left(a_{i}-\mu\right)=\sum\left(p_{i}a_{i}\right)-1\cdot\sum p_{i}a_{i}=0$
\\
Didn'\'{t} go as wanted.

A working definition would be $Var\left[X\right]=\sum p_{i}\left(a_{i}-\mu\right)^{2}$
\begin{definition}
Variance of a random variable

$Var\left[X\right]=\sum p_{i}\left(a_{i}-\mu\right)^{2}=\mathbb{E}\left[\left(X-\mu\right)^{2}\right]=\mathbb{E}\left[\left(X-\mathbb{E}\left[X\right]\right)^{2}\right]$
Let as simplify this equation:
\begin{align*}
\mathbb{E}\left[\left(X-\mathbb{E}\left[X\right]\right)^{2}\right] & =\mathbb{E}\left[X^{2}-\underbrace{2\mathbb{E}\left[X\right]}_{scalar}\cdot X+\mathbb{E}\left[X\right]^{2}\right]=\\
\text{Linearity} & =\mathbb{E}\left[X^{2}\right]-2\cdot\mathbb{E}\left[X\right]\cdot\mathbb{E}\left[X\right]+\mathbb{E}\left[X\right]^{2}=\\
 & =\mathbb{E}\left[X^{2}\right]-2\mathbb{E}\left[X\right]^{2}+\mathbb{E}\left[X\right]^{2}=\\
 & =\mathbb{E}\left[X^{2}\right]-\mathbb{E}\left[X\right]^{2}
\end{align*}
\end{definition}

\begin{thm}[Markov inequality]
 If $X$ is a non-negative R.V and $c>1$ then $p_{r}\left(x\ge c\cdot\mathbb{E}\left[X\right]\right)\le\frac{1}{c}$
\end{thm}

\begin{thm}[Chebyshev inequality]
 If $X$ is a real R.V then $p_{r}\left(\left|X-\mathbb{E}\left[X\right]\right|\ge c\cdot\sigma\right)\le\frac{1}{c^{2}}$
where $var\left[X\right]=\sigma^{2}$ .

In general we can view Chebyshev's inequality as a way to see how
far $X$ is from It's expectation by units of it's standard deviation
($\sigma=\sqrt{Var\left[X\right]}$).

The intuition of proving that claim is to apply Markov's inequality
over $Y=\left(X-\mathbb{E}\left[X\right]\right)^{2}$ and understanding
that $p_{r}\left(\left(x-\mathbb{E}\left[X\right]\right)^{2}\ge c^{2}\sigma^{2}\right)=p_{r}\left(\left|X-\mathbb{E}\left[X\right]\right|\ge c\cdot\sigma\right)$.

\end{thm}

\begin{xca}
If you pick at random a permutation $\pi:\left[n\right]\to\left[n\right]\in S_{n}$
How many fixed points does it have?

In order to answer this question we will create a random variable
that counts the number of fixed points in a given permutation. afterward
we will search it's expectation. Let $X\left(\pi\right)=\sum\limits _{i=1}^{n}X_{i}\left(\pi\right)$
where $X_{i}=\begin{cases}
1 & \pi\text{ has a fixed point at }i\\
0 & otherwise
\end{cases}$ . Thus $X$ is a counter for fixed points in $\pi$. Let calculate
the expectation of $X$: $\mathbb{E}\left[X\right]=\mathbb{E}\left[\sum_{i}X_{i}\right]$
and from linearity of expectation $\mathbb{E}\left[X\right]=\sum\limits _{i=1}^{n}\mathbb{E}\left[X_{i}\right]$.
Now we need to find the expectation of the indicator. $\mathbb{E}\left[X_{i}\right]=1\cdot p_{r}\left(\text{there is a fixed point at }i\right)+0\cdot p_{r}\left(otherwise\right)$.
$p_{r}\left(\text{there is a fixed point at }i\right)=\frac{\left|\left\{ \pi|\pi\text{ has a fixed point at }i\right\} \right|}{n!}.$
The value of $\left|\left\{ \pi|\pi\text{ has a fixed point at }i\right\} \right|$
is $\left(n-1\right)!$ because we fixed one position and then we
still have $n-1$ elements to choose for them a position, therefore
$\left(n-1\right)!$ possible permutations. As a result $p_{r}\left(\text{there is a fixed point at }i\right)=\frac{\left(n-1\right)!}{n!}=\frac{1}{n}$.
Which in turn means that $\mathbb{E}\left[X\right]=\sum\limits _{i=1}^{n}\frac{1}{n}=1$.
The answer to our question is 1.
\end{xca}

\begin{remark}
The above exercise is known as the disoriented secretary.
\end{remark}


\section{Ramsey's Theorem and the probabilistic method}

First lets define $R\left(k,l\right)$: the minimum number of vertices
in a graph that is needed in order to make either $k$ size clique
that is coloured blue on it's edges or $l$ red clique. (clique and
anti-clique)
\begin{example}
What is the number of people needed to be invited to a party in order
to be sure that either 3 know each other or 3 do not know each other?
$R\left(3,3\right)=6$.
\end{example}

\begin{thm}[Erdos-Szekeres 35]
 $R\left(k,l\right)\le{k+l-2 \choose k-1}$. The theorem is without
proof.

For $R\left(k,k\right)\le{k+k-2 \choose k-1}=\binom{2k-2}{k-1}\le2^{2k}=n$,
it implies that in every graph with $n$ vertices there exists a clique
or an anti-clique of size $\ge\frac{1}{2}\log_{2}n$ (because $\frac{1}{2}\log_{2}n=\frac{1}{2}\log_{2}\left(2^{2k}\right)=k\log_{2}2=k$). 

\end{thm}

\begin{remark}
\label{probRVreualsZero}if $X$ is a R.V that is non negative then
$p\left[X\ge c\cdot\mathbb{E}\left[X\right]\right]\le\frac{1}{c}$.
If X takes values in $\N\cup\left\{ 0\right\} $ and $\mathbb{E}\left[X\right]=\epsilon$
then $P\left[X=0\right]\ge1-\epsilon$. (from Markov) 
\end{remark}

\begin{remark}
Chebyshev inequality: $p\left[X=0\right]\le\frac{Var\left[X\right]}{\mathbb{E}\left[X\right]^{2}}$
(proved in problem set 1).
\end{remark}

\begin{corollary}
$\forall N$ every $N$-vertex graph must have a clique or an anti-clique
on $\ge\frac{1}{2}\log_{2}N$ vertices.
\end{corollary}

\begin{thm}[Erdos '47]
 $\exists N$-vertex graph with no clique nor anti-clique $>2\log_{2}N$.
In fact this is true for almost all graphs.
\end{thm}

\begin{proof}
$\Omega_{n}=\left\{ \text{all graphs on }V=\left\{ 1,2,\dots N\right\} \right\} $,
there is a uniform distribution; $\forall G\in\Omega_{N},p_{r}\left[G\right]=\frac{1}{2^{\binom{n}{2}}}$.
\begin{remark}
the number of all possible graphs is $2^{\binom{n}{2}}$ because for
each edge we have two choices and we have $\binom{n}{2}$ edges.
\end{remark}

This probability space is named $G\left(n,\frac{1}{2}\right)$. $k$
is a parameter, $X,Y:\Omega_{n}\to\N\cup\left\{ 0\right\} $. $X\left(G\right)=\text{Number of k-cliques in }G$.
Y is the same but for anti-cliques.
\begin{proposition}
Suppose that $\mathbb{E}\left[X+Y\right]=o\left(1\right)$, where
$o\left(1\right)$ is something that goes to zero when $N\to\infty$.
If this happens then we can conclude that $p_{r}\left[X+Y=0\right]=1-o\left(1\right)$
(remember remark \ref{probRVreualsZero}). Well that means that $X=Y=0$
and that most $N$ - vertex graphs have no k-clique and no k-anti-clique.
So we want to find that $\mathbb{E}\left[X+Y\right]=o\left(1\right)$.
\begin{remark}
Please note that $\mathbb{E}\left[X+Y\right]=2\cdot\mathbb{E}\left[X\right]$
because they are the same R.V but for different colour.
\end{remark}

\end{proposition}

Lets start: We define $X=\sum\limits _{\begin{array}{c}
S\subseteq V\\
\left|S\right|=k
\end{array}}X_{S}$ where $S$ are groups of $k$ vertices from $V$, where
\[
X_{s}=\begin{cases}
1 & \text{S is a clique in G}\\
0 & \text{otherwise}
\end{cases}
\]
 now lets find the expectation:
\begin{align*}
\mathbb{E}\left[X\right] & =\sum\limits _{S}\mathbb{E}\left[X_{S}\right]=\\
 & =\binom{n}{k}p_{r}\left[\underbrace{\left\{ 1,\dots k\right\} }_{vertices}\text{ is a clique}\right]=\\
 & =\binom{n}{k}\cdot\left(\frac{1}{2}\right)^{\binom{k}{2}}
\end{align*}
notice that $\left(\frac{1}{2}\right)^{\binom{k}{2}}$ is the probability
to form an edge $\binom{k}{2}$ times (between all vertices). Lets
inspect the expression above again:
\[
=\frac{\binom{n}{k}}{\left(2\right)^{\binom{k}{2}}}
\]
We need to show that $\binom{n}{k}\cdot\left(\frac{1}{2}\right)^{\bin k2}=o\left(1\right)$.
Notice that $\binom{n}{k}<n^{k}<2^{\frac{k\left(k-1\right)}{2}}\Rightarrow\binom{N}{k}\ll2^{\binom{k}{2}}$.
Consequently the expectation strives to zero as $k$ tends to infinity.
Lets choose $k$ in respect to $N$: 
\begin{align*}
N^{k} & <2^{\frac{k\left(k-1\right)}{2}}\\
N & <2^{\frac{k-1}{2}}\\
\log_{2}N & <\frac{k-1}{2}\\
1+2\cdot\log_{2}N & <k
\end{align*}
 for $k>1+2\log_{n}N$ it holds that $\mathbb{E}\left[X+Y\right]=o\left(1\right)$
and then $p_{r}\left[X+Y=0\right]=1-o\left(1\right)$ as wanted.
\end{proof}
\begin{example}
We will inspect another view of this method: What is the probability
that $G\sim G\left(n,p\right)$ contains a $K_{4}$? ($k_{4}$is a
clique of 4).
\begin{remark}
$G\left(n,p\right);p_{r}\left[G\right]=p^{\left|E\left(G\right)\right|}\cdot\left(1-p\right)^{\binom{n}{2}-\left|E\left(G\right)\right|}$.
Explanation: $p^{\left|E\left(G\right)\right|}$ means the probability
to get the edges. $\left(1-p\right)^{\binom{n}{2}-\left|E\left(G\right)\right|}$
means to not get an edge number of possible edges in the graph minus
the number of edges we already have.
\end{remark}

\begin{thm}
$p=n^{-\frac{2}{3}}$is the threshold probability for the occurrences
of $k_{4}$ in a $G\left(n,p\right)$ graph.
\begin{enumerate}
\item for $p=o\left(n^{-\frac{2}{3}}\right)$ i.e $p\ll n^{-\frac{2}{3}}$
``almost surely'' G contains no $k_{4}$.
\item for $p=\omega\left(n^{-\frac{2}{3}}\right)$ i.e $p\gg n^{-\frac{2}{3}}$
``almost surely'' G contains $k_{4}$.
\end{enumerate}
\begin{remark}
``almost surely'' means: when we sample with p from the graph ,
and the probability strives to zero as $n\to\infty$.
\end{remark}

\end{thm}

$X:G\left(n,p\right)\to\N\cup\left\{ 0\right\} $, $X\left(G\right)=\text{num }k_{4}\text{s in }G$.
$X=\sum\limits _{\left|S\right|=4}X_{s}$ where $X_{s}=\begin{cases}
1 & \text{S is a 4-clique in G}\\
0 & o/w
\end{cases}$. 
\end{example}

\begin{casenv}
\item $p=o\left(n^{-\frac{2}{3}}\right)$ and as a result 
\begin{align*}
\mathbb{E}\left[X\right] & =\sum\mathbb{E}\left[X_{s}\right]=\binom{n}{4}\cdot p_{r}\left[\left\{ 1,2,3,4\right\} \text{ is a clique}\right]=\\
 & =\binom{n}{4}p^{\binom{4}{2}}=\binom{n}{4}p^{6}\le n^{4}p^{6}\\
 & =o\left(1\right)
\end{align*}
As a result we recall $p_{r}\left[X=0\right]=1-\E X=1-o\left(1\right)$
as wanted.
\item $p=\omega\left(n^{-\frac{2}{3}}\right)$: recall that if $X:\Omega\to\N\cup\left\{ 0,1,2,\dots\right\} $
then $p_{r}\left[X=0\right]\le\frac{var\left[X\right]}{\mathbb{E}\left[X\right]^{2}}$.
We want to show that $var\left[X\right]\ll\mathbb{E}\left[X\right]^{2}$.
$X$ is the same, $\mathbb{E}\left[X\right]$ is the same. Lets calculate
the variance: 
\begin{align*}
var\left[X\right] & =\mathbb{E}\left[X^{2}\right]-\mathbb{E}\left[X\right]^{2}=\\
 & =\sum_{\left|S\right|=4}\left(\mathbb{E}\left[X_{s}^{2}\right]-\mathbb{E}\left[X_{s}\right]^{2}\right)+2\cdot\sum_{\begin{array}{c}
S\ne T\\
\left|S\right|=\left|T\right|=4
\end{array}}\left(\underbrace{\mathbb{E}\left[X_{s}\cdot X_{T}\right]-\mathbb{E}\left[X_{s}\right]\cdot\mathbb{E}\left[X_{T}\right]}_{cov\left(X_{s},X_{t}\right)}\right)=\\
 & =\sum_{\left|S\right|=4}\left(\mathbb{E}\left[X_{s}\right]-\mathbb{E}\left[X_{s}\right]^{2}\right)+2\cdot\sum_{\begin{array}{c}
S\ne T\\
\left|S\right|=\left|T\right|=4\\
\left|S\cap T\right|\in\left\{ 2,3\right\} 
\end{array}}\left(\underbrace{\mathbb{E}\left[X_{s}\cdot X_{T}\right]-\mathbb{E}\left[X_{s}\right]\cdot\mathbb{E}\left[X_{T}\right]}_{cov\left(X_{s},X_{t}\right)}\right)
\end{align*}
First let us handle the Cov summation: Notice that there are 2 cases
to handle $\left|S\cap T\right|=2$ and $\left|S\cap T\right|=3$,
that is because the R.V are either independent when $\left|S\cap T\right|=1$
or they are disjointed. when $\left|S\cap T\right|=4$ they are the
same variable hence not in the summation.
\begin{remark}
Note that we indeed compute only where the covariance is no\'{t} zero,
but we can enlarge the value of $Var\left[X\right]$ and make our
life easier by simply computing 
\begin{align*}
\sum\limits _{\left|S\cap T\right|\in\left\{ 2,3\right\} }\mathbb{E}\left[X_{s}\cdot X_{T}\right]-\mathbb{E}\left[X_{s}\right]\cdot\mathbb{E}\left[X_{T}\right]\\
<\sum\limits _{\left|S\cap T\right|\in\left\{ 2,3\right\} }\mathbb{E}\left[X_{s}\cdot X_{t}\right]
\end{align*}
 so we will address only the values in $\sum\limits _{\left|S\cap T\right|\in\left\{ 2,3\right\} }\mathbb{E}\left[X_{s}\cdot X_{t}\right]$.
\end{remark}

\begin{casenv}
\item 
\begin{align*}
\sum\limits _{\left|S\cap T\right|=2}\mathbb{E}\left[X_{s}\cdot X_{t}\right] & =\\
\underbrace{\binom{n}{4}}_{\left(1\right)}\cdot\underbrace{\binom{4}{2}}_{\left(2\right)}\cdot\underbrace{\binom{n-4}{2}}_{\left(3\right)}\cdot\underbrace{p^{11}}_{\left(4\right)}
\end{align*}
 explanation:
\begin{enumerate}
\item number of possibilities to choose the set $S$ (choose 4 vertices
out of the n possible vertices).
\item the possibilities to choose two vertices off of $S$ to be in the
intersection $T\cap S$.
\item all other possibilities for the 2 other vertices of $T$.
\item 11 is the number of edges we needed to to be connected. 
\end{enumerate}
\item 
\[
\sum\limits _{\left|S\cap T\right|=3}\mathbb{E}\left[X_{s}\cdot X_{t}\right]=\binom{n}{4}\cdot\binom{4}{3}\cdot\binom{n-4}{1}\cdot p^{9}
\]
 The explanation is similar to case $i$.
\end{casenv}
\begin{align*}
var\left[X\right] & \le\\
 & \mathbb{E}\left[X\right]+\underbrace{\Theta\left(n^{6}p^{11}\right)+\Theta\left(n^{5}p^{9}\right)}_{\begin{array}{c}
\text{we removed the array because }\\
\text{it is encapsulated in the Thetas}
\end{array}}\\
 & <\mathbb{E}\left[X\right]^{2}=\left(n^{4}p^{6}\right)^{2}=n^{8}p^{12}
\end{align*}
 $var\left[X\right]\ll\mathbb{E}\left[X\right]^{2}$ which means that
$p_{r}\left[X=0\right]\le\frac{var\left[X\right]}{\mathbb{E}\left[X\right]^{2}}\to0$
We got what we wanted: $X\ge1$ almost surely.\\

\end{casenv}
%

\section{Balls and bins}
\begin{fact}
$1+x\le e^{x}$ 
\end{fact}

\begin{fact}
for a small $x$: $1+x\ge e^{x}-O\left(x^{2}\right)$ 
\end{fact}


\subsection{Birthday paradox}

When we throw multiple balls into multiple bins ($n$ bins) we expect
to see the first collision (a second ball into an occupied bin) in
our $\sim\sqrt{n}$ throw. $p_{r}\left[\text{no collision}\right]=1\cdot\left(1-\frac{1}{n}\right)\cdots\left(1-\frac{k-1}{n}\right)$
that is the first ball doesn't have a collision with probability 1\c{ }the
second doesn't have a collision with probability $\left(1-\frac{1}{n}\right)$
because to hit the occupied bin is $\frac{1}{n}$ and that is its
complimentary event, and so on. 
\begin{align*}
p_{r}\left[\text{no collision}\right] & =1\cdot\left(1-\frac{1}{n}\right)\cdots\left(1-\frac{k-1}{n}\right)=\\
\left(\text{fact 32.}\right) & \le1\cdot e^{-\frac{1}{n}}\cdots e^{-\frac{k-1}{n}}=\\
 & =e^{-\left(\frac{1}{n}\sum_{i=1}^{k-1}i\right)}=e^{-\frac{1}{n}\left(\frac{k\left(k-1\right)}{2}\right)}=\\
 & =e^{-\frac{k\left(k-1\right)}{2n}}
\end{align*}
 We can see from the above expression that if $k=o\left(\sqrt{n}\right)$
then $e^{-\frac{k\left(k-1\right)}{2n}}$ the above expression as
n tends to infinity get closer to 1 because $k$ is a lot smaller
than $\sqrt{n}$ and as a result $\limn{e^{\frac{-k^{2}-k}{2n}}}=e^{0}$.
On the other hand $k=\omega\left(\sqrt{n}\right)$ $e^{-\frac{k\left(k-1\right)}{2n}}=o\left(1\right)$
. So we conclude that after $\sqrt{n}$ throws it is quite certain
that there will be collisions.
\begin{example}
Suppose we pick two random subsets $A,B$ of $\left\{ 1,\dots,n\right\} $
of cardinals (size) $k,l$. How likely is it that the two sets are
disjoint? 

Without the loss of generality (WLOG): $A=\left\{ 1\dots k\right\} $
and B is chosen uniformly at random. $p_{r}\left[A\cap B=\emptyset\right]=\frac{\binom{n-k}{l}}{\binom{n}{l}}$
. ($\binom{n-k}{l}$ is because we want to chose $l$ elements that
do not collide with elements from $A$, $\binom{n}{l}$ is all possible
sets of size $l$ - our probability space). 
\begin{align*}
\p{A\cap B=\emptyset} & =\frac{n-k}{n}\cdot\frac{n-k-1}{n-1}\cdot\dots\cdot\frac{n-k-l+1}{n-l+1}=\\
 & =\left(1-\frac{k}{n}\right)\cdot\left(1-\frac{k}{n-1}\right)\cdots\left(1-\frac{k}{n-l+1}\right)\cong\\
 & \cong e^{-\left(\frac{k}{n}+\frac{k}{n-1}+\dots+\frac{k}{n-l+1}\right)}\cong e^{-\frac{kl}{n}}
\end{align*}
\begin{insight}
\begin{align*}
p_{r}\left[A\cap B=\emptyset\right] & =\frac{\binom{n-k}{l}}{\binom{n}{l}}=\frac{\frac{\left(n-k\right)!}{l!\left(n-k-l\right)!}}{\frac{n!}{l!\left(n-l\right)!}}=\\
 & =\frac{\frac{\left(n-k\right)\left(n-k-1\right)\dots\left(n-k-l\right)\left(n-k-l-1\right)\dots}{l!\left(n-k-l\right)\left(n-k-l-1\right)\dots}}{\frac{n\left(n-1\right)\dots\left(n-l\right)\left(n-l-1\right)\dots}{l!\left(n-l\right)\left(n-l-1\right)\dots}}=\\
 & =\frac{\frac{\left(n-k\right)\left(n-k-1\right)\dots\left(n-k-l+1\right)}{l!}}{\frac{n\left(n-1\right)\dots\left(n-l\right)\left(n-l+1\right)}{l!}}=\\
 & =\frac{\left(n-k\right)\left(n-k-1\right)\dots\left(n-k-l+1\right)}{n\left(n-1\right)\dots\left(n-l\right)\left(n-l+1\right)}=\\
 & =\frac{n-k}{n}\cdot\frac{n-k-1}{n-1}\cdots\frac{n-k-l+1}{n-l+1}=\\
 & =\left(1-\frac{k}{n}\right)\left(1-\frac{k}{n-1}\right)\cdots\left(1-\frac{k}{n-l+1}\right)=\\
 & \backsimeq e^{-\left(\frac{k}{n}+\frac{k}{n-1}+\dots+\frac{k}{n-l+1}\right)}
\end{align*}
 please notice that $e^{-\frac{k}{n-1}}\backsimeq e^{-\frac{k}{n}}$
hence 
\[
\backsimeq e^{-\left(\underbrace{\frac{k}{n}+\frac{k}{n}+\dots+\frac{k}{n}}_{l\text{ times}}\right)}=e^{-\frac{kl}{n}}
\]
\end{insight}
{} Now, if $kl=\omega\left(n\right)$ then asymptotic almost surely
$\left(a.a.s\right)$ the two sets intersect because $p_{r}\left[A\cap B=\emptyset\right]=e^{-\frac{kl}{n}}=o\left(1\right)$.
On the other hand if $kl=o\left(n\right)$ then $p_{r}\left[A\cap B=\emptyset\right]=e^{-\frac{kl}{n}}=1-o\left(1\right)$.
\end{example}


\subsection{The coupons collector}
\begin{claim}
It takes roughly $n\log n$ throws of $m$ balls into $n$ bins in
order to make sure that we have inserted at least one ball into every
bin. 
\end{claim}

\begin{proof}
Let $X$ be the R.V that counts the number of empty bins. $X=\sum\limits _{i=1}^{n}X_{i}$
where $X_{i}$ is an indicator R.V for the event ``the $i$-th bin
is empty''. The random variable $X_{i}$ are identically distributed
and their expectation $\E{X_{i}}=\Pr\left[X_{i}=1\right]=\left(1-\frac{1}{n}\right)^{m}$
which means that after $m$ throws a specific bin is still empty.
By linearity of expectation\"{ }$\E X=\sum\limits _{i=1}^{n}\E{X_{i}}=n\cdot\left(1-\frac{1}{n}\right)^{m}\backsimeq n\cdot e^{\frac{-m}{n}}$. 
\begin{casenv}
\item if $m=\omega\left(n\log n\right)$ then $n\cdot e^{\frac{-m}{n}}=o\left(1\right)$
\item if $m=o\left(n\log n\right)$ then $n\cdot e^{\frac{-m}{n}}=\omega\left(1\right)$ 
\end{casenv}
\end{proof}
%
\begin{proof}
Another proof: define $T_{i}$ as the time between having seen $\left(i-1\right)$
bins to the first time bin $i-th$ is visited. this means that $T=\sum T_{i}$.
$\E T=\sum_{i=1}^{n}\E{T_{i}}$. 
\begin{remark}
$T_{i}$ is a geometrically distributed random variable $T_{i}\sim Geo\left(p\right)$.
$\E{T_{i}}=\frac{1}{p}$.
\end{remark}

thus $T_{i}\sim Geo\left(\frac{n-i+1}{n}\right)$, the probability
is defined as follows $n-i+1$ bins are currently empty, so we need
any one of them over the probability space which is $n$. Now we can
compute $\E T=\sum_{i=1}^{n}\E{T_{i}}=\sum_{i=1}^{n}\frac{n}{n-i+1}=n\sum_{i=1}^{n}\frac{1}{n-i+1}$
notice that $\sum_{i=1}^{n}\frac{1}{n-i+1}=\left(\frac{1}{n}+\frac{1}{n-1}+\dots\frac{1}{1}\right)=\sum\limits _{i=1}^{n}\frac{1}{i}$
and thus $\E T=n\sum\limits _{i=1}^{n}\frac{1}{i}=n\left(\log_{e}n+O\left(1\right)\right)$. 
\end{proof}
We would like to say that after exceeding $n\log n+O\left(n\right)$
throws it's almost certain that there is no empty bin. We cannot use
Markov's inequality because it isn't tight enough: 
\[
\Pr\left[T\ge10\cdot\E T\right]\le\frac{1}{10}\iff\Pr\left[T\ge10\cdot n\log n\right]\le\frac{1}{10}
\]
We can use Chebyshev's inequality
\[
\Pr\left[\left|T-\E T\right|\text{\ensuremath{\ge c\cdot\var T}}\right]\le\frac{1}{c^{2}}
\]
Notice that $\var T=\var{\sum T_{i}}$ and that the $T_{i}$ R.Vs
are independent %
\begin{insight}
Notice that each $T_{i}$ is independent because they count number
of tries given a static event. So we check for each$T_{i}$ how many
tries should we do when $i$ slots are already filled.%
\end{insight}
. 
\begin{remark}
A geometric distributed R.V $X$ with parameter $p$ has $\var X=\frac{1-p}{p^{2}}$.
\end{remark}

Therefore 
\[
\var T=\var{\sum T_{i}}=\sum\var{T_{i}}=\sum_{i=1}^{n}\left(\frac{n}{n-i+1}\right)^{2}=n^{2}\sum_{j=1}\frac{1}{j^{2}}\le\frac{\pi^{2}}{6}
\]
 using Chebyshev's inequality: 
\[
\Pr\left[T>n\ln n+cn\right]\le\frac{10}{c^{2}}
\]

\begin{insight}
Not sure whys that chebyshev%
\end{insight}
{} If we choose $c=\sqrt{\log n}$ we get that the probability tends
to 0.

\section{Random walks}

A drunkard takes a walk, the drunkard can either take a step to the
right or to the left. There are multiple questions that might interest
us, such as: 
\begin{enumerate}
\item How far typically will the drunkard reach? $\pm O\left(\sqrt{n}\right)$. 
\item How often will the drunkard return to the origin of the random walk?
$\Theta\left(\sqrt{n}\right)$.
\item We can ask the same questions above higher dimensions. 
\end{enumerate}
We can model the problem as having $N$ RVs $X_{1},\dots,X_{n}$ which
are i.i.d (independent identically distributed). $\p{X_{i}=1}=\p{X_{i}=-1}=\frac{1}{2}$,
$X:=\sum\limits _{i=1}^{N}X_{i}$ ($X$ represents our random walk
after $N$ steps). 
\begin{proposition}
$\p{X\ge a}\le e^{-\frac{a^{2}}{2N}}$ We will use this and prove
it afterward. 
\begin{example}
$\p{X>5\sqrt{N}}\le e^{-\frac{25N}{2N}}=e^{\frac{-25}{2}}$, using
that fact we get that $X$ is quite close to $\sqrt{N}$. In fact
the probability of $X$ being bigger than $\sqrt{N}$ or smaller than
$-\sqrt{N}$ tends to zero.
\end{example}

\end{proposition}

\begin{proof}
$\forall t>0$, $\p{X\ge a}=\p{e^{tX}\ge e^{ta}}$ (the events $e^{tX}\ge e^{ta}$
is identical to $X\ge a$). 
\begin{align*}
\p{\underbrace{e^{tX}}_{R.V}\ge e^{ta}} & \overset{\left(Markov\right)}{\le}\frac{1}{\left(\frac{e^{ta}}{\E{e^{tX}}}\right)}=e^{-ta}\E{e^{tX}}=\\
 & =e^{-ta}\E{e^{t\sum_{i=1}^{N}X_{i}}}=e^{-ta}\E{\prod_{i=1}^{N}e^{tX_{i}}}=\\
\left(independence\right) & =e^{-ta}\prod_{i=1}^{N}\E{e^{tX_{i}}}
\end{align*}
 We know that $X_{i}\left(\omega\right)\in\left\{ -1,1\right\} $,
as a result $\E{e^{tX_{i}}}=\frac{1}{2}\cdot\left(e^{t}+e^{-t}\right)$.
collecting everything up to now: $\p{X\ge a}\le e^{-ta}\left[\frac{e^{t}+e^{-t}}{2}\right]^{N}$.
\begin{claim}
$\forall t\ge0$, $\frac{e^{t}+e^{-t}}{2}\le e^{\frac{t^{2}}{2}}$.
\begin{proof}
$e^{u}=\sum\limits _{j=0}^{\infty}\frac{u^{j}}{j!}\Rightarrow$ 
\begin{align*}
e^{t} & =1+\frac{t}{1!}+\frac{t^{2}}{2!}+\dots\\
e^{-t} & =1+\frac{(-t)}{1!}+\frac{(-t)^{2}}{2!}\dots
\end{align*}
 Notice that the even powers are with the same sign, and the odd powers
cancel each other. we get that
\[
\frac{e^{t}+e^{-t}}{2}=\sum\limits _{j=0}^{\infty}\frac{t^{2j}}{\left(2j\right)!}\le\sum\limits _{j=0}^{\infty}\frac{\left(t^{2}\right)^{j}}{2^{j}\cdot j!}=\sum\limits _{j=0}^{\infty}\frac{\left(\frac{t^{2}}{2}\right)^{j}}{j!}=e^{\frac{t^{2}}{2}}
\]
 
\end{proof}
now lets continue 
\[
\p{X\ge a}\le e^{-ta}\cdot\left[\frac{e^{t}+e^{-t}}{2}\right]^{N}\le e^{-ta}\cdot e^{\frac{t^{2}N}{2}}=e^{\frac{t^{2}N}{2}-ta}
\]
 choose $t=\frac{a}{N}$ and we get 
\[
\p{X\ge a}\le e^{\frac{t^{2}N}{2}-ta}=e^{\frac{\left(\frac{a}{N}\right)^{2}N}{2}-\frac{a}{N}a}=e^{-\frac{a^{2}}{2N}}
\]

as wanted.
\end{claim}

\end{proof}
%
\begin{fact}
$\sum_{i=1}^{N}\binom{N}{i}=2^{n}$
\end{fact}

\begin{fact}
Stirling Formula: $k!\approx\left(\frac{k}{e}\right)^{k}\cdot\sqrt{2\pi k}$ 
\end{fact}

Assuming that $N$ is even $X=\sum\limits _{i=1}^{N}X_{i}$ then the
probability to go back to origin is $\p{X=0}=\frac{\binom{N}{N/2}}{2^{N}}$,
where $\binom{N}{N/2}$ is all possible choosing for $N/2$ steps
right and the the other $N/2$ steps are to the left, and $2^{N}$
are all possible sequence of steps over $N$ walks. The probability
after $N$ steps to be two steps to the right from the origin namely
$\p{X=2}$ is $\frac{\binom{N}{1+N/2}}{2^{N}}$, where $1+N/2$ because
we do $1+N/2$ steps to the right and $\left(N/2\right)-1$ steps
to the left.

A useful lemma to express with better ease expressions like above
$\left(\p{X=0}=\frac{\binom{N}{N/2}}{2^{N}}\right)$:
\begin{lemma}
$\forall0<\alpha<1$, $\bin N{\alpha N}=2^{1+o\left(1\right)}\cdot2^{NH\left(\alpha\right)}$
where $H\left(\alpha\right)=-\alpha\log_{2}\alpha-\left(1-\alpha\right)\log_{2}\left(1-\alpha\right)$.
\begin{proof}
$1=\left(\alpha+\left(1-\alpha\right)\right)^{N}=\sum\limits _{j=0}^{N}\bin Nj\alpha^{j}\left(1-\alpha\right)^{N-j}$
For which $j$ the term in the sum is the biggest? We want to use
derivatives, but we can only use discrete derivation, we compare to
the closest term and and check how close it it to zero: 
\begin{align*}
\bin Nj\alpha^{j}\left(1-\alpha\right)^{N-j} & \backsimeq\bin N{j+1}\alpha^{j+1}\left(1-\alpha\right)^{N-j-1}\\
\Updownarrow\\
\left(j+1\right)\left(1-\alpha\right) & \backsimeq\left(N-j\right)\alpha\Rightarrow j\backsimeq\alpha N
\end{align*}
 consequently 
\[
\frac{1}{N+1}\le\bin N{\alpha N}\alpha^{\alpha N}\left(1-\alpha\right)^{N-\alpha N}\le1
\]
 it is clear why $\bin N{\alpha N}\alpha^{\alpha N}\left(1-\alpha\right)^{N-\alpha N}\le1$
because the largest term in $1=\left(\alpha+\left(1-\alpha\right)\right)^{N}=\sum\limits _{j=0}^{N}\bin Nj\alpha^{j}\left(1-\alpha\right)^{N-j}$
is $\bin N{\alpha N}\alpha^{\alpha N}\left(1-\alpha\right)^{N-\alpha N}$
. Furthermore 1 is the maximal value and the average of a sum is less
then the largest value ($\frac{1}{N+1}$ is the average of sum valued
1 with $N+1$ elements), hence $\frac{1}{N+1}\le\bin N{\alpha N}\alpha^{\alpha N}\left(1-\alpha\right)^{N-\alpha N}$.
Now lets make the above term a bit prettier:
\begin{align*}
\frac{1}{N+1} & \le\bin N{\alpha N}\alpha^{\alpha N}\left(1-\alpha\right)^{N-\alpha N}\le1\\
 & \Updownarrow\\
\frac{1}{N+1}\cdot\frac{1}{\left(\alpha^{\alpha}\left(1-\alpha\right)^{1-\alpha}\right)^{N}} & \le\bin N{\alpha N}\le\frac{1}{\left(\alpha^{\alpha}\left(1-\alpha\right)^{1-\alpha}\right)^{N}}\\
 & \Updownarrow\\
\frac{1}{N+1}\cdot\frac{1}{2^{\log_{2}\left[\left(\alpha^{\alpha}\left(1-\alpha\right)^{1-\alpha}\right)^{N}\right]}} & \le\bin N{\alpha N}\le\frac{1}{2^{\log_{2}\left[\left(\alpha^{\alpha}\left(1-\alpha\right)^{1-\alpha}\right)^{N}\right]}}\\
 & \Updownarrow\\
\frac{1}{N+1}\cdot2^{-N\left[\alpha\log_{2}\left(\alpha\right)+\left(1-\alpha\right)\log_{2}\left(1-\alpha\right)\right]} & \le\bin N{\alpha N}\le2^{-N\left[\alpha\log_{2}\left(\alpha\right)+\left(1-\alpha\right)\log_{2}\left(1-\alpha\right)\right]}\\
 & \Updownarrow\\
2^{N\left(H\left(\alpha\right)-o\left(1\right)\right)}=\frac{2^{NH\left(\alpha\right)}}{N+1} & \le\bin N{\alpha N}\le2^{NH\left(\alpha\right)}
\end{align*}
 which means we get that $\bin N{\alpha N}$ is between $2^{NH\left(\alpha\right)}$
and $2^{NH\left(\alpha\right)-\epsilon}$ ergo $\bin N{\alpha N}=2^{1+o\left(1\right)}\cdot2^{NH\left(\alpha\right)}$,
as wanted.
\end{proof}
\end{lemma}


\section{More bounds}
\begin{remark}
we used previously $e^{tX}$ - the moment generating function. not
that the expectation of such random variable is defined as follows:
$\E{e^{tX}}=\sum_{j=0}^{\infty}\frac{t^{j}}{j!}\E{X^{j}}$
\end{remark}


\subsection{Chernoff bound}
\begin{thm}[Chernoff Bound]
\begin{flushleft}
 Let $X_{1},X_{2},\dots,X_{n}$ be independent indicator R.Vs, $X=\sum X_{i}$
$\mu=\E X$ then
\par\end{flushleft}
\begin{enumerate}
\item $\forall\delta>0$, $\p{X\ge\left(1+\delta\right)\mu}\le\left[\frac{e^{\delta}}{\left(1+\delta\right)^{1+\delta}}\right]^{\mu}$
\item $\forall0<\delta<1$, $\p{X\le\left(1-\delta\right)\mu}\le\left[\frac{e^{-\delta}}{\left(1-\delta\right)^{1-\delta}}\right]^{\mu}$
\end{enumerate}
\begin{proof}
We are proving only part 1, the second part should be similar. 

$\forall t>0$, $\p{X\ge\underbrace{\left(1+\delta\right)\cdot\mu}_{a}}=\p{e^{tX}\ge e^{ta}}\overset{markov}{\le}\frac{\E{e^{tX}}}{e^{ta}}$
. Lets work on the expectation:
\begin{align*}
\E{e^{tX}} & =\E{e^{t\sum_{i=1}^{n}X_{i}}}=\E{\prod e^{tX_{i}}}\overset{independent}{=}\prod\E{e^{tX_{i}}}
\end{align*}
Notice that if $X_{1}v,\dots X_{n}$ are independent so are $e^{tX_{1}}\dots e^{tX_{n}}$.
Because $X_{i}$ are indicators:
\[
\E{e^{tX_{i}}}=p_{i}\cdot e^{t\cdot1}+\left(1-p_{i}\right)e^{t\cdot0}=1+\underbrace{p_{i}\left(e^{t}-1\right)}_{x}\le e^{p_{i}\left(e^{t}-1\right)}
\]
 So 
\[
\E{e^{tX}}\le e^{\left(e^{t}-1\right)\cdot\sum p_{i}}
\]
 and because $\E{X_{i}}=p_{i}$ then 
\begin{align*}
\E{e^{tX}} & \le e^{\left(e^{t}-1\right)\cdot\sum p_{i}}=e^{\left(e^{t}-1\right)\sum\E{X_{i}}}=\\
 & =e^{\left(e^{t}-1\right)\E{\sum X_{i}}}=e^{\left(e^{t}-1\right)\E X}=e^{\left(e^{t}-1\right)\mu}
\end{align*}
 So $\p{X\ge a}\le e^{\left(e^{t}-1\right)\mu}$ choose $t=\ln\left(\delta+1\right)$
and we get
\begin{align*}
\p{X\ge a} & =\p{X\ge\left(\delta+1\right)\mu}\le\frac{e^{\left(e^{t}-1\right)\mu}}{e^{t\left(1+\delta\right)\mu}}=\\
 & =\frac{e^{\left(e^{\ln\left(1+\delta\right)}-1\right)\mu}}{e^{ln\left(1+\delta\right)\cdot\left(1+\delta\right)\mu}}=\frac{e^{\left(\left(1+\delta\right)-1\right)\mu}}{e^{ln\left(1+\delta\right)\cdot\left(1+\delta\right)\mu}}=\\
 & =\frac{e^{\delta\mu}}{\left(1+\delta\right)^{\left(1+\delta\right)\mu}}=\\
 & \left(\frac{e^{\delta}}{\left(1+\delta\right)^{\left(1+\delta\right)}}\right)^{\mu}
\end{align*}
\end{proof}
\end{thm}


\subsection{a brief encounter with martingales and the Azuma inequality}

According to Nati this isn't the standard way martingales are defined
but this is equivalent. 
\begin{definition}[Martingale]
 $\Omega$ is a probability space, $X:\Omega\to\R$ is a real R.V.
$\pi_{0}=\Omega$, $\pi_{0}\prec\pi_{1}\prec\pi_{2}\prec\dots$ are
partitions of $\Omega$. 

Notice that partition $i$ is a refined partition of $i-1$. For example
if $\pi_{1}$ splits $\Omega$ into two halves than $\pi_{2}$ splits
these two halves in some manner. 
\end{definition}

Now having these partitions we define the following random variables.
$X_{i}$ will be associated with partition $\pi_{i}$. $X_{0}\left(\omega\right)=\E X$.
for each $X_{i}$ we define $X_{i}$ over the part of the partition
that an atomic event emerged from. For example assume that $\pi_{1}$
is a partition of $\Omega$ into three events: $A,B,C$ if $\omega\in A$
then $X_{1}\left(\omega\right)=$ the expectation over event $A$
only. (for $B$ or $C$ this applies the same but their own expectation).
Now that we understand how these random variables are defined:
\begin{thm}[Azuma's inequality]
 Let $\Omega$ be a probability space, $X$ a real R.V on $\Omega$,
$\pi_{0}\prec\pi_{1}\prec\pi_{2}\prec\dots$ a sequence of refining
partitions and $X_{0}=\E X$, $X_{1},X_{2},\dots$ the corresponding
R.Vs. If $\forall i$ $\left|X_{i}-X_{i+1}\right|\le d_{i}$ then
$\p{\left|X_{0}-X_{n}\right|>\Delta}\le e^{-\frac{\Delta^{2}}{2\sum_{i=1}^{N}\left(d_{i}^{2}\right)}}$.
(Where $d_{i}$ is defined by us)
\end{thm}

Some intuition: We can think of each partition a different level of
a tree where the leaves are the atomic events. When we calculate $X_{i}$
for level $i$ we look from which leaf came $\omega$ and to what
node from level $i$ that leaf is a descendant. Then we look at the
expectation over all the leaves of that node and that is the value
of $X_{i}\left(\omega\right)$.


\section{The central limit theorem}
\begin{thm}[The central limit Theorem]
 Let $X_{1},X_{2},\dots$ be i.i.d R.V with expectation $\mu$ and
variance $\sigma^{2}$, then 
\[
\frac{\sum_{i=1}^{n}\left(x_{i}\right)-n\mu}{\sigma\sqrt{n}}\overset{d}{\to}\mathcal{N}\left(0,1\right)
\]
 That is $\forall a<b\in\R$ $\limn{\p{a<\frac{\sum_{i=1}^{n}\left(x_{i}\right)-n\mu}{\sigma\sqrt{n}}<b}}=\frac{1}{\sqrt{2\pi}}\int\limits _{a}^{b}e^{-\frac{x^{2}}{2}}dx$. 
\end{thm}

\begin{insight}
First of all note that $\E{X_{i}}=\mu$ for any $i$ and $\var{X_{i}}=\sigma^{2}$.
now when removing from each $X_{i}$ the value of the mean we get
that each $X_{i}$ have a mean 0 and now we can look at it like random
walk with $\sigma$ steps. so we will divide by $\sigma$and now we
have steps of size around $1$. so we divide by $\sqrt{n}$ which
is the expectation of random walk.%
\end{insight}


\part{Linear algebra}

\subsection{recap}
\begin{definition}
A vector space over field $\mathbb{F}$ is a set $V$ with two operations:
Vector addition and scalar multiplication.
\end{definition}

\begin{definition}
$T:V\to W$ ($W,V$ are linear spaces) is a linear map ``if it respects''
the space's linear structure. That is $T\left(u+v\right)=T\left(u\right)+t\left(v\right)$
and $\forall\alpha$, $T\left(\alpha u\right)=\alpha T\left(u\right)$.
\end{definition}

\begin{definition}
Let $V$ be a f.d (finite dimension) vector space $\Rightarrow$ there
are $v_{1}\dots v_{d}\in V$ s.t. $\forall v\in V\exists\alpha_{1}\dots\alpha_{d}$
s.t. $v=\sum\alpha_{i}v_{i}$. The smallest such $d$ is called $dim\left(V\right)$,
and a set $v_{1}\dots v_{d}$ -- a basis.
\end{definition}

\begin{claim}
Let us fix a basis. Then every $v\in V$ is uniquely determined by
the scalars $\alpha_{1}\dots\alpha_{d}$ 
\begin{align*}
v & \leftrightarrow\left(\alpha_{1},\dots\alpha_{n}\right)\\
v_{1} & \leftrightarrow\left(1,0,\dots0\right)
\end{align*}
\end{claim}

\begin{claim}
Every linear mapping $T:V\to W$ ($W$ has a standard basis) is represented
by a matrix $A_{k\times d}$ ($dim\left(W\right)=k$.
\end{claim}

\[
T\left(v\right)=T\left(\alpha_{1},\dots,\alpha_{n}\right)=A\alpha
\]
 %
\begin{insight}
what? $A\alpha$?

What are infinite dim vector spaces? Consider the space of all continuous
functions $\left[a,b\right]\to\R$.%
\end{insight}

\begin{definition}
A vector $v$ is an \textbf{eigenvector} of linear map $T$ with \textbf{eigenvalue}
$\lambda$ of $T$ if $T\left(v\right)=\lambda\cdot v$.
\end{definition}


\subsubsection{\textquotedblleft Rough statements\textquotedblright{} }
\begin{itemize}
\item Linear maps are comprised of rotations and stretching.
\item ``Size of a vector'': The most important example: $L_{2}$ norm.
The size of the vector $\left|\left|\left(x_{1},\dots x_{n}\right)\right|\right|_{2}=\sqrt{\sum x_{i}^{2}}$
.
\item Rotation: a linear map that preserves the $L_{2}$ norm, $\forall v,\left|\left|T\left(v\right)\right|\right|_{2}=\left|\left|v\right|\right|_{2}$.
\item When do we say that $T$ stretches the vector $v$ by $\lambda$?
$T\left(v\right)=\lambda v$. The more standard term is: $v$ is an
eigenvector of $T$ with eigenvalue of $\lambda$.
\end{itemize}
Let's make things more concrete.

\section{Metric spaces}
\begin{definition}
A metric space is a pair $\left(X,d\right)$ where $X$ is a set and
$d$ is a distance function $d:X\times X\to\R$ such that:
\begin{enumerate}
\item $d\left(x,y\right)=d\left(y,x\right)$
\item $\forall x,y\in X$ $d\left(x,y\right)\ge0$ and equality to zero
iff $x=y$. 
\item triangle inequality -- $\forall x,y,z\in X$ , $d\left(x,y\right)+d\left(y,z\right)\ge d\left(x,z\right)$
\end{enumerate}
\end{definition}


\section{normed spaces}

``From size to distance'' (in linear spaces) $\left|\left|u\right|\right|$
``the size of the vector'' we can define as $d\left(u,v\right):=\left|\left|u-v\right|\right|$%
\begin{insight}
like in real numbers, the size is it's abs value, and the distance
is the subtract of two numberss inside abs value%
\end{insight}
.
\begin{definition}
Let $V$ be a real linear space, $\left|\left|\cdot\right|\right|:V\to R$
is a norm if for every $u,v\in V$ and $\alpha\in\R$ it has the following
properties:
\begin{enumerate}
\item $\forall v\in V,\left|\left|v\right|\right|\ge0,\left|\left|v\right|\right|=0\iff v=0$
\item $\forall\alpha\in\R,v\in V,$ $\left|\left|\alpha v\right|\right|=\left|\alpha\right|\cdot\left|\left|v\right|\right|$
\item triangle inequality -- $\forall u,v\in V,\left|\left|u+v\right|\right|\le\left|\left|u\right|\right|+\left|\left|v\right|\right|$
\end{enumerate}
\end{definition}


\subsection{Important examples of norms:}

$L_{p}$ norms, where $1\le p\le\infty$ are defined as follows: $\left|\left|\left(x_{1},\dots,x_{n}\right)\right|\right|_{p}:=\left(\sum\left|x_{i}\right|^{p}\right)^{\frac{1}{p}}$

The most important norms are for $p=1,2,\infty$
\begin{enumerate}
\item $\left|\left|\left(x_{1},\dots,x_{n}\right)\right|\right|_{\infty}:=\max\left\{ \left|x_{i}\right|\right\} $
\item $\left|\left|\left(x_{1},\dots,x_{n}\right)\right|\right|_{1}:=\sum\left|x_{i}\right|$
(Manhattan distance)
\item $\left|\left|\left(x_{1},\dots,x_{n}\right)\right|\right|_{2}:=\left(\sum\left|x_{2}\right|^{2}\right)^{\frac{1}{2}}$
\end{enumerate}
The fact that $L_{p}$ is a norm is not trivial. (The triangle inequality
requires proof)

\subsection{Balls}

Associated with every norm $\left|\left|\cdot\right|\right|$ is its
\uline{unit ball} on a space $V$: $B_{\left|\left|\cdot\right|\right|}:=\left\{ x\in V:\left|\left|x\right|\right|\le1\right\} $.
\begin{itemize}
\item The unit ball of $L_{2}$ is a round ball. 
\item The unit ball for $L_{\infty}$ is a cube: $\left|\left|\left(x_{1},\dots x_{n}\right)\right|\right|_{\infty}\le1\iff\forall i,-1\le x_{i}\le1$.
\item The unit ball of $L_{1}$ in 3 dims is an octahedron ($\left|x\right|+\left|y\right|+\left|z\right|\le1$
in $\R^{3}$ ).
\end{itemize}

\subsubsection{inequalities}

how many vertices, edges and inequalities do we have?
\begin{itemize}
\item cube: 
\begin{align*}
-1 & \le x_{1}\le1\\
-1 & \le x_{2}\le1\\
 & \vdots\\
-1 & \le x_{n}\le1
\end{align*}
 So we have $2^{n}$ vertices.
\item Cross polytope: (the $L_{1}$): $\pm e_{1},\pm e_{2},\dots$ In $\R^{3}$:
\begin{align*}
x+y+z & \le1\\
-x+y+z & \le1\\
x-y+z & \le1\\
 & \vdots
\end{align*}
 we have 8 inequalities in $\R^{3}$ and thus $8$ facets.
\end{itemize}
\begin{fact}
$\forall$ norm $\left|\left|\cdot\right|\right|$, its unit ball
$\ball{}$ is :
\begin{enumerate}
\item centrally symmetric. ($x\in B\Rightarrow-x\in B$ $\left|\left|-x\right|\right|=\left|-1\right|\cdot\left|\left|x\right|\right|\le1$
)
\item convex %
\begin{insight}
if $x_{1}$ and $x_{2}$ are in $B$ then the interval defined by
$\left\{ \lambda x_{1}+\left(1-\lambda\right)x_{2}:\lambda\in\left[0,1\right]\right\} \subseteq B$.%
\end{insight}
\item full dimensional set in $\R^{n}$ %
\begin{insight}
That is if $x\in B$ then $-x\in B$.%
\end{insight}
\end{enumerate}
\end{fact}

\begin{proposition}
if $B\subseteq\R^{n}$ and $B$ has the properties: 
\begin{enumerate}
\item centrally symmetric. ($x\in B\Rightarrow-x\in B$ $\left|\left|-x\right|\right|=\left|-1\right|\cdot\left|\left|x\right|\right|\le1$
)
\item convex %
\begin{insight}
if $x_{1}$ and $x_{2}$ are in $B$ then the interval defined by
$\left\{ \lambda x_{1}+\left(1-\lambda\right)x_{2}:\lambda\in\left[0,1\right]\right\} \subseteq B$.%
\end{insight}
\item full dimensional set in $\R^{n}$ %
\begin{insight}
That is if $x\in B$ then $-x\in B$.%
\end{insight}
{} 
\end{enumerate}
then $\exists\left|\left|\cdot\right|\right|$ s.t $B=B_{\left|\left|\cdot\right|\right|}$.

\end{proposition}

\begin{definition}[isometry]
 Let $\left(X,\left|\left|\cdot\right|\right|\right)$ be a normed
space, a map $T:X\to X$ is an isometry if $\forall x\left|\left|Tx\right|\right|=\left|\left|x\right|\right|$
.
\end{definition}

Another key property of $L_{2}$ is that this norm is derived from
an \uline{inner product}. $\left\langle x,y\right\rangle =\sum x_{i}y_{i}$.
and $\left|\left|x\right|\right|=$$\left\langle x,x\right\rangle ^{\frac{1}{2}}$.

We associate with every $x\in\R^{n}$ the map $y\to\left\langle x,y\right\rangle $
. the map: $\left\langle x,\cdot\right\rangle \left(y\right)$
\begin{definition}[dual norm]
 Let $\left|\left|\cdot\right|\right|$ be a norm on $\R^{n}$, it's
dual norm is $\left|\left|x\right|\right|^{*}=\sup\limits _{y\in\R^{n}}\frac{\left|\left\langle x,y\right\rangle \right|}{\left|\left|y\right|\right|}=\sup\limits _{\left|\left|y\right|\right|=1}\left|\left\langle x,y\right\rangle \right|$%
\begin{insight}
We can write $\max$ instead of $\sup$%
\end{insight}
{} %
\begin{insight}
it's equal because $y\cdot\alpha$ in$\frac{\left|\left\langle x,\alpha y\right\rangle \right|}{\left|\left|\alpha y\right|\right|}=\frac{\left|\left\langle x,y\right\rangle \right|}{\left|\left|y\right|\right|}$
so the expression doesn\'{t} change for any $\alpha$ so we can multiplie
any $y$ by some number to get $\left|\left|y\right|\right|=1$ and
the norm stays the same%
\end{insight}
.
\end{definition}

\begin{xca}
What is the dual norm of $L_{\infty}$?
\begin{sol}
$x\in\R^{n}$ ,
\[
\left|\left|x\right|\right|_{\infty}^{*}=\max_{\left|\left|y\right|\right|_{\infty}=1}\left|\left\langle x,y\right\rangle \right|=\left|\left|x\right|\right|_{1}
\]

\begin{insight}
to start it up lets use some real vector to get a hunch : $x=\left(7,-3,5,9\right)$
so we want $y$ with these signs to the coordinates. $y=\left(+,-,+,+\right)$
so the best $y$ is such that $\left|\left|y\right|\right|_{\infty}=1$
is $y=\left(1,-1,1,1\right)$. %
\begin{insight}
$L_{1}$ and $L_{\infty}$ are the dual norms of each other.%
\end{insight}
\end{insight}

\end{sol}

\end{xca}

\begin{example}
What is the dual norm of $L_{2}$? $x\in\R^{n}$. what is $\left|\left|x\right|\right|_{2}^{*}$?
\[
\left|\left|x\right|\right|_{2}^{*}=\max_{\left|\left|y\right|\right|_{2}=1}\left|\left\langle x,y\right\rangle \right|
\]
 By Cauchy Schwartz the best choice is $y=\frac{x}{\left|\left|x\right|\right|_{2}}$
\begin{insight}
$\left\langle x,y\right\rangle =\left|\left|x\right|\right|\cdot\left|\left|y\right|\right|\cdot cos\left(\alpha\right)$
where $\alpha$ is the angle between the two, thus we want $\cos\alpha=1$
when $x=y$ else this is smaller then it.%
\end{insight}
\[
\left|\left|x\right|\right|_{2}^{*}=\max_{\left|\left|y\right|\right|_{2}=1}\left|\left\langle x,y\right\rangle \right|=\left\langle x,\frac{x}{\left|\left|x\right|\right|_{2}}\right\rangle =\frac{\left(\left|\left|x\right|\right|_{2}\right)^{2}}{\left|\left|x\right|\right|_{2}}=\left|\left|x\right|\right|_{2}
\]
\end{example}

\begin{claim}
If $1\le p,q\le\infty$, $\frac{1}{q}+\frac{1}{p}=1\Rightarrow\left|\left|x\right|\right|_{p}^{*}=\left|\left|x\right|\right|_{q}$.
\end{claim}

\begin{insight}
Find $x$ such that $Ax\thickapprox b$. no such $x$ (too many constraints)
Find $x$ such that $\left|\left|Ax-b\right|\right|$ is minimized%
\end{insight}

\begin{insight}

\subsection{some reminders from previous class(?)}

Norm on real vector space $V$ is a map $\left|\left|\cdot\right|\right|:V\to\R$
if
\begin{enumerate}
\item $\forall x\in V$ $\left|\left|x\right|\right|\ge0$ and there\textquoteright s
equality if and only if $x=0$
\item Linearity: $\forall v\in V$$\lambda\in\R,\left|\left|\lambda v\right|\right|=\left|\lambda\right|\cdot\left|\left|v\right|\right|$
\item triangle inequality: $\left|\left|x+y\right|\right|\le\left|\left|x\right|\right|+\left|\left|y\right|\right|$
\end{enumerate}
%

\subsubsection{The geometric perspective:}

corresponding to $\left|\left|\cdot\right|\right|$ is it's unit ball
$B_{\left|\left|\cdot\right|\right|}=\left\{ x\in V|\left|\left|x\right|\right|\le1\right\} $
\begin{fact}
B is
\begin{enumerate}
\item convex %
\begin{insight}
(coming from traingle inequallity in norms)%
\end{insight}
\item centrally symmetric
\item full dimensional subset of $V$. Means ``it contains a Euclidean''
\begin{insight}
any vector who is not zero has positive norm%
\end{insight}
\end{enumerate}
\end{fact}

on the other hand: every set $B\subseteq V$ with propeties $\left(1\right),\left(2\right),\left(3\right)$
is the unit ball of some norm.
\begin{example}
Very important examples for norms are: $l_{p}$%
\begin{insight}
\"{s}mall $l_{p}$ is when we work with vectors (for sequnce spaces)
and $L_{p}$ is for function spaces%
\end{insight}
{} norm on $\R^{n}$ $\left|\left|x\right|\right|_{p}:=\left(\sum\left|x_{i}\right|^{p}\right)^{\frac{1}{p}}$
in perticular $p=1,2,3,\dots\infty$
\end{example}

\end{insight}

\begin{claim}
For $\infty\ge p\ge1$, $l_{p}$ is a norm.

We prove only the triangle inequality because the other values are
obvious. We need to show: $x,y\in\R^{n}$ $\left|\left|x+y\right|\right|_{p}\le\left|\left|x\right|\right|_{p}+\left|\left|y\right|\right|_{p}$.
We start with Young's inequality
\begin{proposition}
Let $p,q>0,\frac{1}{p}+\frac{1}{q}=1$ , $a,b>0$ then $a\cdot b\le\frac{1}{p}a^{p}+\frac{1}{q}b^{q}$.
\begin{proof}
recall that ``$log$'' is a concave function. %
\begin{insight}
take any two points $x,y$ on the curve of the function, they are
below the curve of the log ($\forall x,y$ $1\ge\lambda\ge0$, $\lambda\log\left(x\right)+\left(1-\lambda\right)\log y\le\log\left(\lambda x+\left(1-\lambda\right)y\right)$)%
\end{insight}
{} then 
\begin{align*}
\log ab & =\log\left(a\right)+\log b=\frac{1}{p}\cdot p\log\left(a\right)+\frac{1}{q}\cdot q\log b\\
 & =\frac{1}{p}\log\left(a^{p}\right)+\frac{1}{q}\log\left(b^{q}\right)\le\text{\ensuremath{\log\left(\frac{1}{p}a^{p}+\frac{1}{q}b^{q}\right)}}
\end{align*}
\end{proof}
\begin{proposition}[Holder's inequality]
$\forall x,y\in\R^{n}$ $\left|\left\langle x,y\right\rangle \right|\le\left|\left|x\right|\right|_{p}\left|\left|y\right|\right|_{q}$
if $\frac{1}{p}+\frac{1}{q}=1$ $p,q>0$.
\begin{proof}
WLOG (without loss of generality) $x,y\ge0$. $\sum x_{j}y_{j}\overset{?}{\le}\left(\sum\left|x_{i}\right|^{p}\right)^{\frac{1}{p}}\cdot\left(\sum\left|x_{i}\right|^{q}\right)^{\frac{1}{q}}$.
\begin{align*}
\frac{\left\langle x,y\right\rangle }{\left|\left|x\right|\right|_{p}\left|\left|y\right|\right|_{q}} & =\sum\limits _{j}\frac{x_{j}}{\left|\left|x\right|\right|_{p}}\frac{y_{j}}{\left|\left|x\right|\right|_{q}}\le\sum\left[\frac{1}{p}\left(\frac{x_{j}}{\left|\left|x\right|\right|_{p}}\right)^{p}+\frac{1}{q}\left(\frac{y_{j}}{\left|\left|y\right|\right|_{y}}\right)^{q}\right]=\\
 & =\frac{1}{p}\frac{\sum x_{j}^{p}}{\left|\left|x\right|\right|_{p}^{p}}+\frac{1}{q}\frac{\sum y_{j}^{q}}{\left|\left|y\right|\right|_{q}^{q}}=\frac{1}{p}+\frac{1}{q}=1
\end{align*}
 we can multiply the denominator and get what we wanted.
\end{proof}
\end{proposition}

\end{proposition}

\begin{proof}
What we want to show $\forall x,y\in\R^{n}$ $\left(\sum\left|x_{j}+y_{j}\right|^{p}\right)^{\frac{1}{p}}\le\left(\sum\left|x_{j}\right|^{p}\right)^{\frac{1}{p}}+\left(\sum\left|y_{j}\right|^{p}\right)^{\frac{1}{p}}$.
$l_{p}$ is a norm: again WLOG $x,y\ge0$ then 
\begin{align*}
\left|\left|x+y\right|\right|_{p}^{p} & =\sum_{j}\left(x_{j}+y_{j}\right)^{p}=\sum\left(x_{j}+y_{j}\right)\left(x_{j}+y_{j}\right)^{p-1}=\\
 & =\sum x_{j}\left(x_{j}+y_{j}\right)^{p-1}+y_{j}\left(x_{j}+y_{j}\right)^{p-1}=\\
 & =\left[\sum x_{j}\underbrace{\left(x_{j}+y_{j}\right)^{p-1}}_{v_{j}}\right]+\left[\sum y_{j}\underbrace{\left(x_{j}+y_{j}\right)^{p-1}}_{vj}\right]
\end{align*}
 (we will talk about $v_{j}$ in a moment) so %
\begin{insight}
we use holder inequallity with $\left|\left\langle u,v\right\rangle \right|\le\left|\left|u\right|\right|_{p}\left|\left|v\right|\right|_{q}$%
\end{insight}
{} with holders: 
\begin{align*}
 & \le\left|\left|x\right|\right|_{p}\left|\left|x+y\right|\right|_{p}^{p-1}+\left|\left|y\right|\right|_{p}\cdot\left|\left|x+y\right|\right|_{p}^{p-1}=\\
 & \left|\left|x+y\right|\right|_{p}^{p-1}\left(\left|\left|x\right|\right|_{p}+\left|\left|y\right|\right|_{p}\right)
\end{align*}
divide by $\left|\left|x+y\right|\right|_{p}^{p-1}$ and we get what
we wanted.\\
Lets get back to the missing part: If $v_{j}=\left(x_{j}+y_{j}\right)^{p-1}\Rightarrow\left|\left|v\right|\right|_{q}=\left|\left|x+y\right|\right|_{p}^{p-1}$
explanation: $\frac{1}{q}+\frac{1}{p}=1\iff q=\frac{p}{p-1}$ 
\begin{align*}
\left|\left|v\right|\right|_{q} & =\left(\sum v_{j}^{q}\right)^{\frac{1}{q}}=\left(\sum\left(x_{j}+y_{j}\right)^{q\left(p-1\right)}\right)^{\frac{1}{q}}=\left(\sum\left(x_{j}+y_{j}\right)^{p}\right)^{\frac{1}{q}}=\\
 & =\left(\sum\left(x_{j}+y_{j}\right)^{p}\right)^{\frac{p-1}{p}}=\left(\left(\sum\left(x_{j}+y_{j}\right)^{p}\right)^{\frac{1}{p}}\right)^{p-1}=\left|\left|x+y\right|\right|_{p}^{p-1}
\end{align*}
 Now using this $v_{j}$ we can simply use holder's inequality to
get $\le\left|\left|x\right|\right|_{p}\left|\left|x+y\right|\right|_{p}^{p-1}+\left|\left|y\right|\right|_{p}\cdot\left|\left|x+y\right|\right|_{p}^{p-1}=$
and finish our proof.
\end{proof}
\end{claim}

\begin{definition}[operator norm]
 $\left(U,\left|\left|\cdot\right|\right|_{U}\right)\overset{\psi}{\to}\left(V,\left|\left|\cdot\right|\right|_{v}\right)$
. 
\[
\left|\left|\psi\right|\right|:=\sup_{u\in U}\frac{\left|\left|\psi u\right|\right|_{V}}{\left|\left|u\right|\right|_{U}}
\]
 If $\psi$ is linear then 
\[
\left|\left|\psi\right|\right|:=\sup_{u\in U}\frac{\left|\left|\psi u\right|\right|_{V}}{\left|\left|u\right|\right|_{U}}=\sup_{\begin{array}{c}
u\in U\\
\left|\left|u\right|\right|_{U}=1
\end{array}}\left|\left|\psi u\right|\right|_{V}
\]
 
\end{definition}

an important special case: $V=\R,\left|\left|\cdot\right|\right|_{V}=\left|\cdot\right|$.
to every $x\in\R^{n}$ corresponds a linear map $\R^{n}\to\R,y\in\R^{n}\to\left\langle x,y\right\rangle $,
``We think of x as a linear functional''. So the norm of the operator
$\left\langle x,\cdot\right\rangle $: $\left|\left|x\right|\right|^{*}:=\max\limits _{\left|\left|y\right|\right|=1}\left|\left\langle x,y\right\rangle \right|$.%
\begin{insight}
Therefore if $\left|\left|\cdot\right|\right|$ is a norm on $\R^{n}$
this induces a new norm called the dual of $\left|\left|\cdot\right|\right|$
denoted $\left|\left|x\right|\right|^{*}:=\max\limits _{\left|\left|y\right|\right|=1}\left|\left\langle x,y\right\rangle \right|$.%
\end{insight}
\begin{insight}
to understand the notations: $\left|\left|x\right|\right|^{*}:\left|\left|\left\langle x,\cdot\right\rangle \right|\right|\max\limits _{\left|\left|y\right|\right|=1}\left|\left\langle x,y\right\rangle \right|$%
\end{insight}
{} 
\begin{definition}
The unit ball $B^{*}$ of the dual norm is called the polar body of
$B$ and is defined as follows:

\[
B^{*}=\left\{ y\in\R^{n}|\forall x\in B\left\langle x,y\right\rangle \le1\right\} 
\]
\end{definition}

\begin{example}
Question: Let $B$ be the unit ball of some norm $\left|\left|\cdot\right|\right|$.
What can we say about the unit ball of the dual norm?
\begin{itemize}
\item $l_{\infty}$ Ball is $\left[-1,1\right]^{n}$ $2^{n}$ vertices $2n$
facets.
\item $l_{1}$ Ball is the cross polytope $\sum\left|x_{i}\right|\le1$.
Has $2n$ vertices (the standard basis $\pm e_{i}$), and $2^{n}$
facets $\epsilon\in\left\{ -1,1\right\} ^{n}$ $\sum\epsilon_{i}\cdot x_{i}\le1$.
\end{itemize}
What is the polar $B^{*}$ of the unit cube $B=\left[-1,1\right]^{n}$?
\[
B^{*}=\left\{ y\in\R^{n}|\begin{array}{c}
-1\le x_{1}\le1,-1\le x_{2}\le1\dots\\
\sum x_{i}y_{i}\le1
\end{array}\right\} 
\]
 So given $y$ the ``most difficult'' $x$ to satisfy %
\begin{insight}
$x$ that maximizes this expression%
\end{insight}
. $x_{i}=\text{sign}\left(y_{i}\right)$ in which case $\sum x_{i}y_{i}=\sum\left|y_{i}\right|$.
\end{example}


\section{\textquotedblleft What do linear maps do\textquotedblright ?}

$A$ is a matrix

\subsection{Stretch}

If $Av=\lambda v$ holds we say that $v$ is an eigenvector of $A$
with eigenvalue $\lambda$. \\
$\left(A-\lambda I\right)v=0$ i.e. $v$ is in the kernel of $\left(A-\lambda I\right)$.
In particular $\left(A-\lambda I\right)$ is \textbf{singular}, hence
$\det\left(A-\lambda I\right)=0$, and because $\det\left(A-\lambda I\right)$
is polynomial in $\lambda$. ``The characteristic polynomial $p_{A}$
of $A$'' $\Rightarrow$ Every eigenvalue is a root of $p_{A}$.
\begin{proposition}
If A is a real $n\times n$ matrix and if $u_{1}\dots u_{n}$ are
eigenvectors with \uline{distinct} eigenvalues then $u_{1}\dots u_{n}$
are linearly independent.

\begin{insight}
This is useful in case $u_{i}$ forms a basis. $x=\sum\alpha_{i}u_{i}$
then $Ax=\sum\alpha_{i}\lambda_{i}u_{i}$ then we are in a good situation.
``We have good understanding of the linear operator''%
\end{insight}
\begin{proof}
Assume to the contrary that they are linearly dependent and let $k$
be the smallest size of a linearly dependent subset. 
\[
\sum\limits _{i=1}^{k}b_{i}u_{i}=0\Rightarrow A\left(\sum_{i=1}^{k}b_{i}u_{i}\right)=\left(\sum_{i=1}^{k}b_{i}\lambda_{i}u_{i}\right)=0
\]
furthermore we know that $\sum_{i=1}^{k}\left(b_{i}u_{i}\right)=0$
(from our assumption) so
\[
0=\sum_{i=1}^{k}\left(b_{i}u_{i}\right)=\lambda_{k}\sum_{i=1}^{k}\left(b_{i}u_{i}\right)=\sum_{i=1}^{k}\left(b_{i}\lambda_{k}u_{i}\right)
\]
 Let us subtract the first from the second:
\begin{align*}
0 & =\sum_{i=1}^{k}b_{i}\left(\lambda_{i}-\lambda_{k}\right)u_{i}=\\
 & =b_{1}\left(\lambda_{1}-\lambda_{k}\right)u_{1}+\dots+b_{k-1}\left(\lambda_{k-1}-\lambda_{k}\right)u_{k-1}+b_{k}\left(\lambda_{k}-\lambda_{k}\right)u_{k}=\\
 & =b_{1}\left(\lambda_{1}-\lambda_{k}\right)u_{1}+\dots+b_{k-1}\left(\lambda_{k-1}-\lambda_{k}\right)u_{k-1}+0=\\
 & =\sum_{i=1}^{k-1}b_{i}\left(\lambda_{i}-\lambda_{k}\right)u_{i}
\end{align*}
 We found a shorter non trivial linear dependency, which is a contradiction. 
\end{proof}
How should we deal with general linear transformation? we should find
eigenvectors and if they are a basis we can use it to understand the
matrix. but what do we do when this isn't the case? What can go wrong? 
\end{proposition}

\begin{example}
example of a matrix where the eigenvectors do not form a basis:
\end{example}

\[
A=\left(\begin{array}{cccc}
0 & 1 & \dots & 0\\
0 & \ddots & \ddots & \vdots\\
\vdots & \ddots & \ddots & 1\\
0 & \dots & 0 & 0
\end{array}\right)
\]
 
\[
\det\left(\lambda I-A\right)=\left|\begin{array}{cccc}
\lambda & -1 & \dots & 0\\
0 & \ddots & \ddots & \vdots\\
\vdots & \ddots & \ddots & -1\\
0 & \dots & 0 & \lambda
\end{array}\right|=\lambda^{n}=0
\]
$\lambda=0$ is the only eigenvalue.
\[
A\left(\begin{array}{c}
x_{1}\\
\vdots\\
x_{n-1}\\
x_{n}
\end{array}\right)=\left(\begin{array}{c}
x_{2}\\
\vdots\\
x_{n}\\
0
\end{array}\right)=0\cdot x
\]
 $x_{2}=x_{3}=\dots=x_{n}=0$ the set of eigenvectors corresponding
to $\lambda=0$ is $\left\langle e_{1}\right\rangle $. $A^{n}=0$
(When we reapply $A$ it maps to zero after enough times) How to deal
with it, at least in practical situations?
\begin{definition}
Say that a matrix $B$ $n\times n$ is ``problem free'' if it has
distinct eigenvalues.
\end{definition}

\begin{fact}
``Most'' matrices are problem free. %
\begin{insight}
thats why we don\'{t} need to deal with the ones that don't%
\end{insight}
\end{fact}


\subsection{Rotations}

Linear maps which are isometries. It suffices that $l_{2}$ norms
are preserved $\left|\left|x-y\right|\right|_{2}^{2}=\left|\left|x\right|\right|_{2}^{2}+\left|\left|y\right|\right|_{2}^{2}-2\left\langle x,y\right\rangle $
\begin{insight}
the only norm for which such things exist is $l_{2}$.%
\end{insight}

\begin{definition}
A real matrix $A$ is called orthogonal if $A^{T}A=I$.

\begin{insight}
(This means that they are inverses too)%
\end{insight}

\end{definition}

\begin{definition}
A complex matrix $U$ is called unitary if $UU^{*}=I$ %
\begin{insight}
where $U^{*}=\overline{U^{T}}$?%
\end{insight}
\end{definition}

\begin{claim}
$\left\langle x,Ay\right\rangle =\left\langle A^{T}x,y\right\rangle $
\begin{proof}
in general $\left\langle u,v\right\rangle =\sum u_{i}v_{i}$, so
\begin{align*}
\left\langle x,Ay\right\rangle  & =\sum x_{i}\left(Ay\right)_{i}=\sum_{i}x_{i}\sum_{j}a_{ij}y_{j}=\sum_{i,j}a_{ij}x_{i}y_{j}=\\
\end{align*}
 The $a_{ij}$ of $A^{T}$.
\[
=\sum_{i,j}a_{ji}^{T}x_{i}y_{j}=\sum_{j}y_{j}\sum_{i}a_{ji}^{T}x_{i}=\left\langle A^{T}x,y\right\rangle 
\]
\begin{insight}
$a_{ji}^{T}$ is like take the original matrix\c{ }transpose it and
then look at $a_{ji}$%
\end{insight}
\end{proof}
\end{claim}

\begin{thm}
The linear map $\R^{n}\to\R^{n}$ corresponding to a matrix $A$ is
an isometry if and only if $A$ is orthogonal.
\begin{proof}
Suppose that $A$ is orthogonal. We want to show that $\forall x$
$\left|\left|Ax\right|\right|_{2}=\left|\left|x\right|\right|_{2}$.
need to show :
\[
\left\langle Ax,Ax\right\rangle \overset{?}{=}\left\langle x,x\right\rangle 
\]
we know that $\left\langle Ax,By\right\rangle =\left\langle B^{\top}Ax,y\right\rangle $
so

\[
\left\langle Ax,Ax\right\rangle =\left\langle A^{\top}Ax,x\right\rangle =\left\langle Ix,x\right\rangle =\left\langle x,x\right\rangle 
\]
 as wanted.
\end{proof}
\end{thm}

\begin{claim}
An orthogonal matrix $A$ is one that satisfies $AA^{T}=I$($\iff A^{T}A=I$)
\begin{proof}
the proof is for $AA^{T}=I\iff A^{T}A=I$: $A$ has a right inverse,
(namely $A^{T}$) $\Rightarrow$ it has a left inverse, say $B$ $BA=I$
\[
A^{T}=IA^{T}=BAA^{T}=BI=B
\]
 and we have what we wanted.
\end{proof}
\end{claim}

.What does $\left(HH^{T}\right)_{i,j}$ equals to? $\left(HH^{T}\right)_{i,j}=$
the inner product of the $i$th row of $H$ and it's $j$th row. $H^{T}H=$
the inner product of the $i$th column of $H$ and the $j$th column.
\begin{proposition}
A real square matrix $A$ is orthogonal iff $x\to Ax$ is an isometry,
i.e. $\forall x\left|\left|Ax\right|\right|_{2}=\left|\left|x\right|\right|$.
\begin{proof}
If $A$ is orthogonal: $\left|\left|Ax\right|\right|^{2}=\left\langle Ax,Ax\right\rangle =\left\langle A^{T}Ax,x\right\rangle =\left|\left|x\right|\right|^{2}$.
The other side ($A$ is an isometry)
\[
\left|\left|x+y\right|\right|^{2}=\left|\left|x\right|\right|^{2}+\left|\left|y\right|\right|^{2}+2\left\langle x,y\right\rangle 
\]
 we also know that $\left|\left|x+y\right|\right|^{2}=\left|\left|A\left(x+y\right)\right|\right|^{2}=\left|\left|Ax+Ay\right|\right|^{2}=\left|\left|Ax\right|\right|^{2}+\left|\left|Ay\right|\right|^{2}+2\left\langle Ax,Ay\right\rangle $
we know that $\left|\left|Ax\right|\right|^{2}=\left|\left|x\right|\right|^{2},\left|\left|Ay\right|\right|^{2}=\left|\left|y\right|\right|^{2}$.
So 
\begin{align*}
\left|\left|Ax\right|\right|^{2}+\left|\left|Ay\right|\right|^{2}+2\left\langle Ax,Ay\right\rangle  & =\left|\left|x+y\right|\right|^{2}=\left|\left|x\right|\right|^{2}+\left|\left|y\right|\right|^{2}+2\left\langle x,y\right\rangle \\
2\left\langle Ax,Ay\right\rangle  & =2\left\langle x,y\right\rangle \\
\left\langle Ax,Ay\right\rangle  & =\left\langle x,y\right\rangle 
\end{align*}
\end{proof}
We saw: if $A$ is an isometry $\forall x,y$. $\left\langle Ax,Ay\right\rangle =\left\langle x,y\right\rangle $
\[
\left\langle A^{T}Ax,y\right\rangle =\left\langle x,y\right\rangle 
\]
$fix$ x $\forall y$ $\left\langle A^{T}Ax-x,y\right\rangle =0\Rightarrow$$A^{T}A=I$.%
\begin{insight}
because for all $y$ it happens that $A^{T}Ax=x$ which is $\left\langle A^{T}Ax-x,y\right\rangle =\left\langle x-x,y\right\rangle =0$
this wasn\'{t} the case if $A^{T}Ax\ne x$.%
\end{insight}

\end{proposition}


\section{Singular value decomposition}

\begin{insight}
spectral: overall name that speaks on eigenvalues and eigenvectors%
\end{insight}

\begin{thm}[The spectral theorem for real symmetric matrices]
 Every real symmetry matrix $A$ can be expressed as: $A=UDU^{T}$
where $U$ is orthogonal, $D$ is diagonal on $D's$ diagonal are
the eigenvalues of $A$. (Without proof) %
\begin{insight}
Note that $U^{T}AU=D$%
\end{insight}
\end{thm}

let us remember again the operator norm:
\begin{definition}
\begin{insight}
Again (note that you defined it back there\.{)}%
\end{insight}
{} Let there be a real matrix $A$ (not necessarily square), it's operator
norm is $\left|\left|A\right|\right|_{op}=\max\limits _{x}\frac{\left|\left|Ax\right|\right|}{\norm x}=\max\limits _{\left|\left|x\right|\right|=1}\left|\left|Ax\right|\right|$.
\end{definition}

\begin{thm}[SVD -- Singular Value Decomposition]
 Every real matrix $A_{m\times n}$ can be expressed as 
\[
A_{m\times n}=U_{m\times m}D_{m\times n}V_{n\times n}^{T}
\]
 where $U$ and $V$ are orthogonal, and $D$ is ``diagonal''. 
\[
D=\left(\begin{matrix}\sigma_{1} & \dots & 0 & 0\\
0 & \sigma_{2} & \dots & 0\\
\vdots &  & \ddots & \vdots\\
0 & \dots &  & \sigma_{n}\\
0 & 0 & \dots & 0\\
\vdots &  &  & \vdots\\
0 & 0 & \dots & 0
\end{matrix}\right),m\ge n
\]
where $\sigma_{1}\ge\sigma_{2}\ge\dots\ge0$ $'A'$s singular values.
\end{thm}

Note that for $n\ge m$ $D=\left(\begin{matrix}\sigma_{1} & \dots & 0 & 0 & 0 & \dots & 0\\
0 & \sigma_{2} & \dots & 0 & 0 &  & 0\\
\vdots &  & \ddots & \vdots & 0 &  & 0\\
0 & \dots &  & \sigma_{n} & 0 & \dots & 0
\end{matrix}\right),m\ge n$

\begin{proof}
The proof is by induction on $m+n$. %
\begin{insight}
Base: $m=n=1$ thus $A$ is a $1\times1$ matrix. so let us use 
\[
I_{1}AI_{1}
\]
 $A$ is diagonal and it's eigenvalue is it's own value. Step%
\end{insight}
\begin{insight}
We use induction on $m+n$. Main idea: We want $U^{T}AV=D$, we will
do it in stages, not all at once. First stage : We seek orthogonal
matrices $U_{1},V_{1}$ such that 
\[
U_{1}^{T}AV_{1}=\left(\begin{array}{c|ccc}
\text{real number} & 0 & \dots & 0\\
\hline 0\\
\vdots &  & B\\
0
\end{array}\right)
\]
We want to apply induction to repeat this process and then the process
should do it for us.. We take the decomposition of the orthogonal
to those that take the key roles of $SVD$ in $B$. and by induction
we are done. $\tilde{U}B\tilde{V}^{T}=\tilde{D}$. so we take 
\[
\left(\begin{array}{c|ccc}
1 & 0 & \dots & 0\\
\hline 0\\
\vdots &  & \tilde{U}\\
0
\end{array}\right)
\]
 to get the wanted $U_{2}$ and $V_{2}$.%
\end{insight}

We will find $U_{1},V_{1}$ orthogonal s.t. 
\[
U_{1}^{T}AV_{1}=\left(\begin{array}{c|ccc}
\left|\left|A\right|\right|_{op} & 0 & \dots & 0\\
\hline 0\\
\vdots &  & B\\
0
\end{array}\right)
\]
First lets figure out where $\left|\left|A\right|\right|_{op}$ comes
from. $\exists v_{1},u_{1}$ , $\left|\left|v_{1}\right|\right|=\left|\left|u_{1}\right|\right|=1$
s.t. $Av_{1}=\underbrace{\left|\left|A\right|\right|_{op}}_{\sigma_{1}}u_{1}$
(From the definition of the operator norm%
\begin{insight}
Actually Nati didn'\'{t} show this, 
\[
u_{1}=\frac{Ax}{\left|\left|Ax\right|\right|},v_{1}=x
\]
 and $x$ is the vector that gains $\left|\left|A\right|\right|_{op}$
so $\left|\left|v_{1}\right|\right|=1\iff v=\frac{x}{\left|\left|x\right|\right|}$
and we get $u_{1},v_{1}$ as wanted.%
\end{insight}
). $V_{1}=\left(v_{1}|V_{2}\right)$ (we extend $v_{1}$ into an orthogonal
matrix)%
\begin{insight}
(meaning that $\left(v_{1}|V_{2}\right)$ is an orthogonal matrix
where it's first row is $v_{1}$)%
\end{insight}
{} and $U_{1}=\left(u_{1}|U_{2}\right)$. So 
\begin{align*}
U_{1}^{T}AV_{1} & =\left(\begin{array}{c}
u_{1}^{T}\\
\hline U_{2}^{T}
\end{array}\right)A\left(\begin{array}{c|c}
v_{1} & V_{2}\end{array}\right)=\\
 & =\left(\begin{array}{c}
u_{1}^{T}\\
\hline U_{2}^{T}
\end{array}\right)\left(A\begin{array}{c|c}
v_{1} & AV_{2}\end{array}\right)=\\
 & =\left(\begin{array}{c|c}
u_{1}^{T}Av_{1} & u_{1}^{T}AV_{2}\\
\hline U_{2}^{T}Av_{1} & U_{2}^{T}AV_{2}
\end{array}\right)
\end{align*}
 so $u_{1}^{T}Av_{1}=\left|\left|A\right|\right|_{op}$, $U_{2}^{T}Av_{1}=0$
because $U_{2}^{T}\cdot\left(\left|\left|A\right|\right|_{op}u_{1}\right)=0$
because $U_{2}$ is orthogonal %
\begin{insight}
$u_{1}$ is some basis vector from the orthogonal matrix which means
that $\left\langle u_{1},u_{i}\right\rangle =0$%
\end{insight}
. Currently: we found orthogonal matrices $U_{1},V_{1}$ s.t 
\[
U_{1}^{T}AV=\left(\begin{array}{ccc|ccc}
 & \sigma_{1} &  & w_{1}\\
\hline  & 0 & \\
 & \vdots &  & B\\
 & 0 & 
\end{array}\right)
\]
 and we want to show that $w_{1}=0$. The plan: Show that if $w_{1}\ne0$
then we can find a vector that $A$ extends by a vector $>\sigma_{1}$
\begin{insight}
(in contradiction)%
\end{insight}
. Observe that if a matrix is multiplied by an orthogonal matrix either
from left or from the right, the operator norm does not change.
\[
\left(\begin{array}{ccc|ccc}
 & \sigma_{1} &  & w_{1}\\
\hline  & 0 & \\
 & \vdots &  & B\\
 & 0 & 
\end{array}\right)\cdot\left(\begin{array}{c}
\sigma_{1}\\
\hline \\
w_{1}^{^{T}}\\
\\
\end{array}\right)=\left(\begin{array}{c}
\sigma_{1}^{2}+\left|\left|w_{1}\right|\right|^{2}\\
\\
\vdots\\
\\
\end{array}\right)
\]
 This extends by a factor $\ge\sqrt{\sigma_{1}^{2}+\left|\left|w_{1}\right|\right|^{2}}\ge\sigma_{1}$.
What we conclude from that: 
\[
\sigma_{1}=\left|\left|A\right|\right|_{op}=\left|\left|U_{1}^{T}AV_{1}\right|\right|_{op}\ge\frac{\left|\left|\left(\begin{array}{c}
\sigma_{1}^{2}+\left|\left|w_{1}\right|\right|^{2}\\
\\
\vdots\\
\\
\end{array}\right)\right|\right|}{\sqrt{\sigma_{1}^{2}+\left|\left|w_{1}\right|\right|^{2}}}\ge\sqrt{\sigma_{1}^{2}+\left|\left|w_{1}\right|\right|^{2}}\underbrace{>}_{\text{*}}\sigma_{1}
\]
\end{proof}
\begin{itemize}
\item From the assumption that $w_{1}\ne0$ so we got a contradiction.
\end{itemize}
Hence $w_{1}=0$ and we got what we wanted. now we can apply this
idea using induction. Finishing by induction: We found $U_{1},V_{1}$
orthogonal s.t 
\[
U_{1}^{T}AV_{1}=\left(\begin{array}{ccc|ccc}
 & \sigma_{1} &  & 0 & \dots & 0\\
\hline  & 0 & \\
 & \vdots &  &  & B\\
 & 0 & 
\end{array}\right)
\]
 and $B=\tilde{U}\tilde{D}\tilde{V}^{T}$. So 
\begin{align*}
A & =U_{1}\left(\begin{array}{ccc|ccc}
 & 1 &  & 0 & \dots & 0\\
\hline  & 0 & \\
 & \vdots &  &  & \tilde{U}\\
 & 0 & 
\end{array}\right)\left(\begin{array}{ccc|ccc}
 & \sigma_{1} &  & 0 & \dots & 0\\
\hline  & 0 & \\
 & \vdots &  &  & \tilde{D}\\
 & 0 & 
\end{array}\right)\left(\begin{array}{ccc|ccc}
 & 1 &  & 0 & \dots & 0\\
\hline  & 0 & \\
 & \vdots &  &  & \tilde{V}\\
 & 0 & 
\end{array}\right)V_{1}^{T}\\
U_{1}^{T}AV & =\left(\begin{array}{ccc|ccc}
 & 1 &  & 0 & \dots & 0\\
\hline  & 0 & \\
 & \vdots &  &  & \tilde{U}\\
 & 0 & 
\end{array}\right)\left(\begin{array}{ccc|ccc}
 & \sigma_{1} &  & 0 & \dots & 0\\
\hline  & 0 & \\
 & \vdots &  &  & \tilde{D}\\
 & 0 & 
\end{array}\right)\left(\begin{array}{ccc|ccc}
 & 1 &  & 0 & \dots & 0\\
\hline  & 0 & \\
 & \vdots &  &  & \tilde{V}\\
 & 0 & 
\end{array}\right)
\end{align*}
\begin{insight}
That we get in this process 
\[
\left(\begin{array}{ccc|ccc}
 & \sigma_{1} &  & w_{1}\\
\hline  & 0 & \\
 & \vdots &  & B\\
 & 0 & 
\end{array}\right)\to\left(\begin{array}{ccc|c|cc}
 & \sigma_{1} &  & w_{1}\\
\hline  & 0 &  & \sigma_{2}\\
\hline  & \vdots &  &  & C\\
 & 0 &  & 
\end{array}\right)
\]
and $\sigma_{1}\ne\sigma_{2}$ because in the next step we find $\left|\left|B\right|\right|_{op}$
and not $\left|\left|A\right|\right|_{op}$ again. So we get $\left|\left|A\right|\right|_{op},\left|\left|B\right|\right|_{op},\left|\left|C\right|\right|_{op}\dots$%
\end{insight}

Suppose that indeed $A=UDV^{T}$ as in the theorem, how to interpret
the numbers $\sigma_{1}\ge\sigma_{2}\ge\dots\ge0$? 
\begin{align*}
AA^{T} & =UD\underbrace{V^{T}V}_{I}D^{T}U^{T}=\\
 & =UDD^{T}U^{T}
\end{align*}
\[
DD^{T}=\left(\begin{matrix}\sigma_{1}^{2} & \dots & 0 & 0\\
0 & \sigma_{2}^{2} & \dots & 0\\
\vdots &  & \ddots & \vdots\\
0 & \dots &  & \sigma_{n}^{2}
\end{matrix}\right)_{m\times m}
\]
 Based on the spectral theorem for real symmetric matrices (and note
that $AA^{T}$ is r.s%
\begin{insight}
real symmetric%
\end{insight}
) we conclude that $\sigma_{1}^{2},\dots,\sigma_{n}^{2},\underbrace{0,\dots,0}_{m-n}$
are the eigenvalues of the (Positive semi-definite) real matrix $AA^{T}$.

We take another look at the proof of the SVD theorem to make sure
it is understood by everyone.
\begin{thm}[SVD]
 Every real $m\times n$ matrix $A$ can be expressed as $A=UDV^{T}$
where $U_{m\times m}$ is an orthogonal matrix, $V_{n\times n}$ is
orthogonal and $D$ is ``diagonal'' and on the diagonal of $D$
$\sigma_{1}\ge\dots\ge\sigma_{\min\left\{ m,n\right\} }\ge0$
\end{thm}

When we look at $U^{T}AV$ we get a ``diagonal'' matrix.

When we prove this statement we want to find $U_{1}AV_{1}^{T}$ such
that they give us 
\[
U_{1}^{T}AV_{1}=\left(\begin{array}{ccc|ccc}
 & \sigma_{1} &  & 0 & \dots & 0\\
\hline  & 0 & \\
 & \vdots &  &  & B\\
 & 0 & 
\end{array}\right)
\]
 and by induction we can get for $B$ orthogonal matrices that get
as what we need. so how do we do it? first we find$u_{1},v_{1}$ such
that $Av_{1}=\sigma u_{1}$and $\left|\left|v_{1}\right|\right|=1=\left|\left|u_{1}\right|\right|$.
Now if $U_{1}$ and $V_{1}$ are defined as in the proof we get
\[
U_{1}^{T}AV_{1}=\left(\begin{array}{ccc|ccc}
 & \sigma_{1} &  &  & w\\
\hline  & 0 & \\
 & \vdots &  &  & B\\
 & 0 & 
\end{array}\right)
\]
How large can we make $\sigma$? it can be as large as $\left|\left|A\right|\right|_{op}$
- the biggest such $\sigma$. and we get 
\[
U_{1}^{T}AV_{1}=\left(\begin{array}{ccc|ccc}
 & \left|\left|A\right|\right|_{op} &  &  & w\\
\hline  & 0 & \\
 & \vdots &  &  & B\\
 & 0 & 
\end{array}\right)
\]
 our claim is that for $\sigma=\left|\left|A\right|\right|_{op}$
then $w$ is 0. else we get $\left|\left|A\right|\right|_{op}>\sigma$
which is impossible. So we lets choose some specific vector to get
a contradiction 
\[
\left(\begin{array}{c}
\left|\left|A\right|\right|_{op}\\
\hline w^{T}
\end{array}\right)
\]
Which leads to:
\[
\sigma_{1}=\left|\left|A\right|\right|_{op}=\left|\left|U_{1}^{T}AV_{1}\right|\right|_{op}\ge\frac{\left|\left|\left(\begin{array}{c}
\sigma_{1}^{2}+\left|\left|w_{1}\right|\right|^{2}\\
\\
\vdots\\
\\
\end{array}\right)\right|\right|}{\sqrt{\sigma_{1}^{2}+\left|\left|w_{1}\right|\right|^{2}}}\ge\sqrt{\sigma_{1}^{2}+\left|\left|w_{1}\right|\right|^{2}}\underbrace{>}_{\text{*}}\sigma_{1}
\]
 so $w=0$. and we're done.

Question: Why does the sub matrix $B$ has a smaller operator norm?Working
with $B$ is the same as working with matrix $A$ but with a row filled
with zeros, thus it has a smaller or equal operator norm.

Note $A=UDV^{T}\iff A=\sum\limits _{i=1}^{\min\left\{ m,n\right\} }\sigma_{i}u_{i}\otimes v_{i}$
meaning the outer product of $u_{i}$and $v_{i}$ where $u_{i}$ is
the ith column of $U$ and $v_{i}$ is the ith column %
\begin{insight}
column because we have there $V^{T}$%
\end{insight}
of $V$. 
\[
UDV^{T}=\left(\begin{matrix}| &  & |\\
u_{1} & \dots & u_{m}\\
| &  & |
\end{matrix}\right)D\left(\begin{matrix}- & v_{1} & -\\
 & \vdots\\
- & v_{n} & -
\end{matrix}\right)
\]


\subsection{another operator norm}
\begin{xca}
How to evaluate how ``close'' is matrix $B$ to matrix $A$?
\begin{sol}
two answers:
\begin{enumerate}
\item $\left|\left|A-B\right|\right|_{op}$
\item $\left|\left|A-B\right|\right|_{2}=\left|\left|A-b\right|\right|_{\text{Frobenius}}=\sqrt{\sum\left(a_{ij}-b_{ij}\right)^{2}}$
\end{enumerate}
\end{sol}

\end{xca}

$\left|\left|A\right|\right|_{frobenius}=\sqrt{\sum a_{ij}^{2}}$
note that $\left|\left|A\right|\right|_{fr}^{2}=trace\left(AA^{T}\right)$.
because:
\begin{align*}
\left|\left|A\right|\right|_{fr}^{2} & =trace\left(AA^{T}\right)=\sum_{i}\left(AA^{T}\right)_{ii}\\
\end{align*}
 
\[
\left(AA^{T}\right)_{ii}=\left|\left|a_{i}\right|\right|^{2}
\]
 where $a_{i}$ is the ith row of $A$ so 
\[
\left|\left|a_{i}\right|\right|^{2}=\sum_{j}a_{ij}^{2}
\]
 %
\begin{insight}
\[
trace\left(AA^{T}\right)=\sum_{i,j}a_{ij}^{2}
\]
\end{insight}

\begin{claim}
Let $A$ be a real matrix, $U$ an orthogonal matrix, note that 
\[
\left|\left|A\right|\right|_{fr}=\left|\left|AU\right|\right|_{fr}
\]
\begin{proof}
$ $ 
\[
\left|\left|AU\right|\right|_{fr}^{2}=tr\left(\left(AU\right)\left(AU\right)^{T}\right)=tr\left(AUU^{T}A^{T}\right)=tr\left(AIA^{T}\right)=tr\left(AA^{T}\right)=\left|\left|A\right|\right|_{fr}^{2}
\]
\end{proof}
\end{claim}

\begin{thm}
Let $A_{m\times n}$ be a real matrix and let $k<m,n$. For every
$B_{m\times n}$ matrix of rank $\le k$ there holds $\left|\left|A-B\right|\right|_{op}\ge\sigma_{k+1}$
where $\sigma_{1}\ge\dots\ge0$ are the singular values of $A$. Equality
holds for$B=UD^{\left(k\right)}V^{T}$, where $A=UDV^{T}$ and 
\[
D^{\left(k\right)}=\left(\begin{array}{cccccc}
\sigma_{1} & 0 & \dots &  & \dots & 0\\
0 & \ddots &  &  &  & \vdots\\
\vdots &  & \sigma_{k}\\
 &  &  & 0\\
 &  &  &  & \ddots & \vdots\\
\vdots &  &  &  &  & 0\\
0 &  &  &  &  & 0
\end{array}\right)
\]
 
\[
B=\sum_{i=1}^{k}\sigma_{i}u_{i}\otimes v_{i}
\]
(We can\'{t} do better then than $\left|\left|A-B\right|\right|_{op}\ge\sigma_{k+1}$)
\end{thm}

Note this: $A-B_{k}=U\left(\begin{array}{ccccc}
0\\
 & \ddots\\
 &  & 0\\
 &  &  & \sigma_{k+1}\\
 &  &  &  & \ddots
\end{array}\right)V^{T}\Rightarrow$$\left|\left|A-B\right|\right|_{op}=\sigma_{k+1}$
\begin{proof}
\begin{insight}
all Lecture page 31%
\end{insight}
{} What do we need to show? for $C$ matrix such that $\text{rank}\left(C\right)\le k\Rightarrow\left|\left|A-C\right|\right|_{op}\ge\sigma_{k+1}$.
\begin{insight}
to show this we need to show some vector $z$ such that $\left|\left|\left(A-C\right)z\right|\right|\ge\sigma_{k+1}$%
\end{insight}
Lets start: $\exists z,\left|\left|\left(A-C\right)z\right|\right|_{2}\ge\sigma_{k+1}$
where $\left|\left|z\right|\right|_{2}=1$. How can we find such a
vector? Plan: We will select $z$ in such a way that $Cz=0$ and then
$\left|\left|\left(A-C\right)z\right|\right|=\left|\left|Az-Cz\right|\right|=\left|\left|Az\right|\right|$
- lets do just that: Since $\text{rank}\left(C\right)\le k$ the kernel
of $C$ intersects (not in $\bar{0}$) every linear subspace $W$
of $\dim W\ge k+1$. Let $W=span\left\{ v_{1}\dots v_{k+1}\right\} $
(the columns of $V$).
\[
z=\sum_{i=1}^{k+1}\alpha_{i}v_{i};Cz=0;\left|\left|z\right|\right|_{2}=1
\]
 
\begin{align*}
\left(A-C\right)z & =Az-Cz=Az=\sum_{i}\sigma_{i}\left(u_{i}\otimes v_{i}\right)z=\\
 & =\sum\sigma_{i}u_{i}\left\langle v_{i},z\right\rangle =\sum^{m}\sigma_{i}u_{i}\left\langle v_{i},z\right\rangle \\
\left(1\right) & =\sum_{i=1}^{k+1}\sigma_{i}u_{i}\left\langle v_{i},z\right\rangle =\sum_{i=1}^{k+1}\sigma_{i}\left\langle v_{i},z\right\rangle u_{i}=
\end{align*}
\begin{enumerate}
\item everything beyond $k+1$ equals zero. because $v_{1}\dots v_{m}$are
orthonormal then $\left\langle v_{i},z\right\rangle =\left\langle v_{i},\sum_{j}^{k}\alpha_{j}v_{j}\right\rangle $
equals 0 for every $j\ne i$.
\end{enumerate}
Because the $u_{i}$'s are orthogonal:
\begin{align*}
\left|\left|\left(A-C\right)z\right|\right|^{2} & =\sum_{i=1}^{k+1}\sigma^{2}\left\langle v_{i},z\right\rangle ^{2}\ge\sigma_{k+1}^{2}\underbrace{\sum\limits _{i=1}^{k+1}\left\langle v_{i},z\right\rangle ^{2}}_{(*)}\\
 & =\sigma_{k+1}^{2}\cdot1
\end{align*}

\begin{itemize}
\item because it's an orthogonal basis.%
\begin{insight}
To my understanding because it\'{s} an orthonormal basis%
\end{insight}
\end{itemize}
\end{proof}
%
\begin{definition}
A Real symmetric matrix is said to be positive semi-definite (PSD)
if all it's eigenvalues are non negative.
\end{definition}

\begin{definition}
A Real symmetric matrix is said to be positive definite (PD) if all
it's eigenvalues are positive.
\end{definition}

\begin{thm}
The following are equivalent (TFAE) regarding a real symmetric matrix
$A$
\begin{enumerate}
\item $A$ is PSD.
\item $A=MM^{T}$ for some real $M$.
\item $\forall x,xAx^{T}\ge0$
\end{enumerate}
\end{thm}

\begin{proof}
$ $
\end{proof}
\begin{itemize}
\item[$1\Rightarrow2$:]  recall the spectral theorem for real symmetric matrices: it is possible
to express $A=U\Lambda U^{T}$ where $U$ is orthogonal $\Lambda$
is diagonal, the diagonal terms in $\Lambda$ are $A$'s eigenvalues
and the row of $U$ are the corresponding eigenvectors. Let $M=U\sqrt{\Lambda}$
($\sqrt{\Lambda}$ = The diagonal matrix whose $i$ths entry is $\sqrt{\lambda_{i}}$.)
$MM^{T}=U\sqrt{\Lambda}\sqrt{\Lambda}U^{T}=U\Lambda U^{T}=A$. As
wanted.
\item[$2\Rightarrow3$:]  Since $A=MM^{T}$ $xAx^{T}=xMM^{T}x=\left|\left|xM\right|\right|_{2}^{2}\ge0$,
As wanted.
\item[$3\Rightarrow1$:]  Let $v$ be an eigenvector with eigenvalue $\lambda$. We know that
$vAv^{T}\ge0$ but on the other hand 
\[
0\le vAv^{T}=\lambda\left\langle v,v\right\rangle =\lambda\left|\left|v\right|\right|_{2}^{2}
\]
\begin{insight}
Notice that the rows are revresed: (full explanation for the equation)
\[
0\le vAv^{T}=v\left(\lambda v^{T}\right)=\lambda\left\langle v,v\right\rangle =\lambda\left|\left|v\right|\right|_{2}^{2}
\]
\end{insight}
{} so $\lambda\ge0$ because $\left|\left|v\right|\right|_{2}^{2}\ge0$.
\end{itemize}
\begin{corollary}
The collection of all $n\times n$ PSD matrices form a \textbf{cone}
$PSD_{n}$. Namely if $A,B\in PSD_{n}\Rightarrow A+B\in PSD_{n}$
and if $t>0\Rightarrow tA\in PSD_{n}$.%
\begin{insight}
We can use the 3 property we proved to show that.%
\end{insight}
\end{corollary}


\subsection{The variational definition of eigenvalues of real symmetric matrices}
\begin{thm}[Rayleigh-Ritz]
 Let $A$ be real symmetric matrix and let $\lambda_{1}\ge\lambda_{2}\ge\dots\ge\lambda_{n}$
be it's eigenvalues. Then 
\[
\lambda_{1}=\max\frac{xAx^{T}}{\left|\left|x\right|\right|_{2}^{2}}=\max\limits _{\left|\left|x\right|\right|_{2}=1}\left(xAx^{T}\right)
\]
$\lambda_{2}=\max\limits _{x\perp x_{1}}\frac{xAx^{T}}{\left|\left|x\right|\right|_{2}^{2}}$,
$\lambda_{3}=\max\limits _{x\perp x_{1}\perp x_{2}}\frac{xAx^{T}}{\left|\left|x\right|\right|_{2}^{2}}$,
and so on. 
\end{thm}

To prove that we will use some analysis:
\begin{remark}
Recall that $\left(\frac{f}{g}\right)^{'}=0\iff fg'=f'g$ we'll have
$f=\left(xAx^{T}\right)$ and $g=\left|\left|x\right|\right|^{2}$
and we want to get the maximum of that.
\end{remark}

\begin{remark}
$\left(xA\right)_{\nu}=\sum\limits _{s}a_{s\nu}x_{s}$ and $xAx^{T}=\sum_{\alpha,\beta}x_{\alpha}a_{\alpha,\beta}x_{\beta}$.
\end{remark}

We will show that if $x$ maximizes $\frac{xAx^{T}}{\left|\left|x\right|\right|_{2}^{2}}$
then $x$ is an eigenvector. First lets derive $\left|\left|x\right|\right|_{2}^{2}$:
$\frac{\partial}{\partial x_{i}}\left|\left|x\right|\right|_{2}^{2}=2x_{i}$.
Now from the previous remark:
\[
\frac{\partial}{\partial x_{i}}xAx^{T}=\frac{\partial}{\partial x_{i}}\left(\sum\limits _{\alpha,\beta}a_{\alpha,\beta}x_{\alpha}x_{\beta}\right)=2\left(Ax\right)_{i}
\]
 twice is because the row can either come from $a_{ij}x_{i}x_{j}$
or from $a_{si}x_{s}x_{i}$. for $x$ to maximize the Rayleigh quotient
$\frac{xAx^{T}}{\left|\left|x\right|\right|^{2}}$ a necessary condition
is that $\forall i$ $\underbrace{\left(Ax\right)_{i}}_{f'}\cdot\underbrace{\left|\left|x\right|\right|^{2}}_{g}=\underbrace{\left(xAx^{t}\right)}_{f}\underbrace{x_{i}}_{g'}$
notice that $\left|\left|x\right|\right|^{2},\left(xAx^{T}\right)$
are numbers, and 
\[
\underbrace{\left(Ax\right)_{i}}_{f'}\cdot\underbrace{\left|\left|x\right|\right|^{2}}_{g}=\underbrace{\left(xAx^{t}\right)}_{f}\underbrace{x_{i}}_{g'}\iff\left(Ax\right)\left|\left|x\right|\right|^{2}=\left(xAx^{T}\right)x
\]
so we've seen that $\left(Ax\right)\underbrace{\left|\left|x\right|\right|^{2}}_{number}=\underbrace{\left(xAx^{T}\right)}_{number}x$.
We've seen that if $x$ maximizes $\frac{xAx^{T}}{\left|\left|x\right|\right|_{2}^{2}}$
then $x$ is an eigenvector.

\begin{proof}[proof of Rayleigh - Ritz:]
 (We prove only for $\lambda_{1}$, the next steps ($\lambda_{2}\dots$)
need to use Lagrange multipliers. %
\begin{insight}
The concept is that he shows here that if x maxsimizes this (f is
the nominator and $g$ is the denominator of what we want) so if $x$
maximises it then it is an eigenvalue!! So if we maximise this thing
we should find the biggest lambda! (maximal x is always eigenvector
so we use maximal eigenvalue to get the maximum of $\frac{xAx^{T}}{\left|\left|x\right|\right|_{2}^{2}}$)%
\end{insight}
) Using the spectral theorem: $A=V\Lambda V^{T}\Rightarrow$$xAx^{T}=xV\Lambda V^{T}x^{T}$
We think of it as $y\Lambda y^{T}=\sum\lambda_{i}y_{i}^{2}$ where
$y=xV^{T}$ So 
\[
=\underbrace{\sum\lambda_{i}\left(Vx\right)_{i}^{2}}_{\text{positive items}}\le\left(\max\lambda_{i}\right)\cdot\sum\left(Vx\right)_{i}^{2}=\lambda_{1}\sum\left(Vx\right)_{i}^{2}=\lambda_{1}\left|\left|Vx\right|\right|^{2}=\lambda_{1}\left|\left|x\right|\right|^{2}
\]
 if $Ax=\lambda_{1}x$ then $\frac{xAx^{T}}{\left|\left|x\right|\right|^{2}}=\lambda_{1}$%
\begin{insight}
for $\lambda_{2}$: The same thing but if $x$ is orthogonal to the
previous $x$ then the equations zeros out.%
\end{insight}
\end{proof}
%
\begin{thm}[Rayliegh-Ritz 2]
 Let $A$ be a real symmetric matrix with eigenvalues $\lambda_{1}\ge\lambda_{2}\ge\dots\ge\lambda_{n}$
. Then$\lambda_{n}=\min\frac{xAx^{T}}{\left|\left|x\right|\right|^{2}};$$\lambda_{n-1}=\min\limits _{x\perp x_{n}}\frac{xAx^{T}}{\left|\left|x\right|\right|^{2}}$.
(With no proof)
\end{thm}

How do we derive the second theorem from the first? we look at the
minus of the same thing.
\begin{thm}[Courant-Fischer]
 $A_{n\times n}$ real symmetric, with eigenvalues $\lambda_{1}\ge\lambda_{2}\ge\dots\ge\lambda_{n}$.
Then $\forall i$ 
\[
\lambda_{i+1}=\min_{\begin{array}{c}
F\\
\dim F=i
\end{array}}\max_{x\perp F}\left\{ \frac{xAx^{T}}{\left|\left|x\right|\right|^{2}}\right\} 
\]
 and $\max\limits _{\dim F=j}\min\limits _{x\perp F}\frac{xAx^{T}}{\left|\left|x\right|\right|^{2}}=\lambda_{n-j}$.
\begin{insight}
Why is there an optimum for all these subspaces? (There's infinite
amount of subspaces) Because of compactness.%
\end{insight}
\end{thm}

\begin{thm}[The Interlacing theorem]
Let $A$ be a real symmetric matrix with eigenvalues $\lambda_{1}\ge\dots\ge\lambda_{n}$.
Let $B$ be the matrix that results by erasing the i-th row and the
i-th column of $A$. %
\begin{insight}
B is real symmetric too. and has $n-1$ eigenvalues.%
\end{insight}
{} Let $\mu_{1}\ge\dots\ge\mu_{n-1}$ be the eigenvalues of $B$. So
$\lambda_{1}\ge\mu_{1}\ge\lambda_{2}\ge\mu_{2}\ge\dots\ge\lambda_{n-1}\ge\mu_{n-1}\ge\lambda_{n}$.
\begin{insight}
This is derived from Rayleigh -- Ritz. We need to connect $A$ to
$B$. so the way to think about this is that we can pretendt that
$B$ was never born and we keep working with $A$. we insist that
the $i'th$ coordinate of $x$ is 0. Doing so we get what we get $B$.%
\end{insight}
\end{thm}


\section{Perron -- Frobenius Theorems}
\begin{thm}[Perron]
 Let $A$ be a square positive matrix %
\begin{insight}
positive means its real.%
\end{insight}
(i.e $\forall i,j$ $a_{i,j}>0$) Then $A$ has a positive eigenvalue
$\rho>0$ s.t. for every other eigenvalue $\mu$ (in general complex)
$\rho>\left|\mu\right|$. The eigenvector corresponding to $\rho$
is positive, no other eigenvector is positive.
\end{thm}

The proof's stages:
\begin{enumerate}
\item $\rho>0$ corresponding eigenvector $>0$.
\item $\forall\mu\ne\rho,\,\,\rho>\left|\mu\right|$.
\item $\rho$ is a simple eigenvalue. (Has an algebraic multiplicity 1 --
Note that we don't actually show algebraic multiplicity 1) 
\item no other positive eigenvector (We don't prove this)
\end{enumerate}
\begin{insight}
we want a theorem for non negative matrices, rather then strictly
positive matrix. to show that $A^{k}$ is positive we switch every
non zero number to 1. it still holds that if this new matrix is positive
so is $A^{k}$. now that $A$ is zeros or 1s we know that $A$ is
a the adjacency matrix of a directed graph. $Au=\lambda u$ then $A^{k}u=\lambda^{k}u$.
If we knew persons theorem and we know that $A^{k}>0$ we can derive
the same conclusions for Perron's Theorem to a non negative matrix.
$A^{k}\ge0$.%
\end{insight}

\begin{proof}[proof of Perron's Theorem]
\begin{insight}
We want to think about it as an optimization problem.%
\end{insight}
Let us observe: if $v\ge0,v\ne0\Rightarrow Av>0$. Let $S\subseteq\mathbb{R}^{n}$
be defined as: 
\[
S=\left\{ u\in\R^{n}:\left|\left|u\right|\right|_{2}=1,u\ge0\right\} 
\]
\begin{insight}
this is a compact set%
\end{insight}

We define $\forall u\in S$, 
\[
L\left(u\right):=\max\left\{ b>0:\begin{array}{c}
\left(Au\right)_{i}\ge bu_{i},\\
\forall i\,\,u_{i}\ne0
\end{array}\right\} =\min_{\begin{array}{c}
i\\
u_{i}\ne0
\end{array}}\frac{\left(Au\right)_{i}}{u_{i}}
\]
 %
\begin{insight}
The coordinate that holds be to become bigger is the $i$ that makes
this$\frac{\left(Au\right)_{i}}{u_{i}}$ the minimum%
\end{insight}

It can be shown that $L\left(u\right)$ is compact%
\begin{insight}
The edge of a set is a point that contains around it points in the
set and points outside of the set.

A close set is a set such it contains it's edge.

Compact set:
\begin{enumerate}
\item Let there be a metric space $X$. A compact set $Y\subseteq X$ is
called compact if for any series $y_{n}\in Y$ There is a converging
sub series of $y_{n}$ such that it converge into a point in $Y$.
\item Every closed set $Z\subseteq Y$ is compact.
\item a compact set is bound.
\item The image of a compact set by a function on that set is also compact.
\item Any non empty compact set contains a maximum and a minimum.
\item A close and bound set in $\R^{n}$ is compact.
\end{enumerate}
\end{insight}
, hence $\rho:=\max\limits _{u\in S}L\left(u\right)$. Let $v\in S$
be a vector realizing $Av\ge\rho v$%
\begin{insight}
from the L$\left(u\right)$ statement%
\end{insight}
. %
\begin{insight}
we want to show that this is an eigenvalue. but we saw that we can
get a bigger value by contradiction.%
\end{insight}
We want to show $Av=\rho v$, if it is -- were done, else we assume
by contradiction that $Av-\rho v>0$, it follows that %
\begin{insight}
non negative and non zero vector multiplied by a positive matrix is
positive.%
\end{insight}
$A\left(Av-\rho v\right)>0$. Because that is a positive vector $\exists\epsilon>0$
s.t. $A\left(Av-\rho v\right)>\epsilon Av$, Let $c>0$ s.t. $cAv$
is a unit vector, actually $cAv\in S$ so%
\begin{insight}
from this: $Av-\rho v>0$ we expand the term:
\begin{align*}
AAv-Apv & >\epsilon Av\\
cAAv-Acpv & >\epsilon cAv\\
AcAv & >\epsilon cAv+Acpv\\
A\left(cAv\right) & >\left(p+\epsilon\right)cAv
\end{align*}
 Notice that $cAv$ is a vector in $S$ for which we defined $\rho=\max\limits _{u\in S}L\left(u\right)$
which will lead us to a contradiction. HELL YES!%
\end{insight}
\[
A\left(cAv\right)>\left(\rho+\epsilon\right)cAv
\]
{} This is contrary to the maximality of $\rho$.
\begin{claim}
If $\mu\ne\rho$ is an eigenvalue of $A$ then $\rho>\left|\mu\right|$. 
\begin{proof}
Let $Ay=\mu y\iff\forall i\,\,\underbrace{\left(Ay\right)_{i}}_{=\sum\limits _{j}a_{i,j}y_{j}}=\mu y_{i}$
So $\left|\sum\limits _{j}a_{ij}y_{j}\right|=\left|\mu\right|\left|y_{i}\right|$
then from triangle inequality 
\[
\left|\mu\right|\cdot\left|y_{i}\right|=\left|\sum a_{ij}y_{j}\right|\le\sum\left|a_{i,j}\right|\left|y_{j}\right|=\sum a_{i,j}\left|y_{j}\right|
\]
. Denote $z_{i}=\left|y_{i}\right|$ so $\sum a_{i,j}z_{i}\ge\left|\mu\right|\cdot z_{i}$
meaning:
\[
Az\ge\left|\mu\right|z
\]
 From this it follows that $\rho\ge\left|\mu\right|$. %
\begin{insight}
$\rho$ is defined as the largest number such that $Av\ge\rho v$
so $\mu$ cannot be bigger then $\rho$.%
\end{insight}
\begin{insight}
z should be normalized and then $z\in S$ and as a result $\rho$
is the maximal eigenvalue.%
\end{insight}
\end{proof}
\end{claim}

\begin{claim}
$\rho$ has a geometric multiplicity i.e. the linear space of eigenvectors
with eigenvalue $\rho$ is $1-dimensional$
\begin{proof}
suppose that $Au=\rho u$. We can write $u=\phi+i\psi$ and each part
will be an eigenvector separately. This means we can just assume that
$u$ is a real vector. Generally $u$ may be complex, but then the
vectors $Re\left(u\right),Im\left(u\right)$ are also eigenvectors
with eigenvalue $\rho$. WLOG $u$ is real.%
\begin{insight}
We should apply this for both vectors - to the imaginary vector too
and we get the same contradiction.%
\end{insight}

Let us assume by contradiction that $Av=\rho v,Au=\rho u$. $\exists\epsilon$
(might be negative too) s.t. $v+\epsilon u\ge0$ with at least one
zero coordinate. because $v,u$ are eigenvectors of $\rho$: $A\left(v+\epsilon u\right)=\rho\left(v+\epsilon u\right)$.
(This is a contradiction because $A\left(v+\epsilon u\right)$ can't
have an entry that is zero because a positive matrix multiplied by
a non negative vector creates a positive vector -- can't have a cell
with zero in it like the vector $\rho\left(v+\epsilon u\right)$ and
from this we conclude that there can't be another vector $u$ such
that $Au=\rho u$.)

(To my understanding we did not show in this proof that the algebraic
multiplicity is 1)%
\begin{insight}
There cant be an entry with zero and thus the equation is a contradiction
- not sure whats going on. not sure what happened here%
\end{insight}

\end{proof}
\end{claim}

 Additional item in Perron\'{s} Theorem: $A$ positive square matrix
with $\rho$ as before...
\[
\limn{\frac{A}{\rho}}^{n}=\frac{x_{r}\otimes x_{\ell}}{\left\langle x_{r},x_{\ell}\right\rangle }
\]
 ($\otimes$ means outer-product) Where $x_{r},x_{\ell}$, are the
right and left eigenvectors of $A$ corresponding to $\rho$.%
\begin{insight}
We divide each eigenvalue by $\rho^{n}$%
\end{insight}

\end{proof}
%
To pass from Perron to Frobenius: if $A^{^{k}}>0$ for some $k\in\N$
then $A$ satisfies the conclusions of Persons theorem.

Questions: Given $A\ge0$ when does $\exists k$ s.t. $A^{k}>0$ ?
$\iff$ Given a directed graph $G$ under what conditions does $\exists k\in\N$
s.t. $\forall x,y\in V$ $\exists$ directed $x\to y$ of length $k$
.
\begin{thm}
(\textcolor{red}{This wasn't stated in class explicitly}): (Perron
- Frobenius) If $A$ is non negative and there exists a $k\in\N$
such that $A^{k}>0$ then Perron's statement holds for this matrix
too.
\end{thm}

\begin{definition}
We say that directed graph $G=\left(V,E\right)$ has period $p$ if
it is possible to split $V=V_{1}\cup V_{2}\dots\cup V_{p}$ s.t if
$x\to y\in E$ iff $\exists i$ such that $x\in V_{i},y\in V_{i+1\mod p}$.
\end{definition}

\begin{claim}
This is the case provided
\begin{enumerate}
\item $G$ is strongly connected.
\item Aperiodic. %
\begin{insight}
If it's periodic we know that $k$ moves us to some group.%
\end{insight}
\end{enumerate}
\end{claim}


\section{finite Markov chains}

This is a system ($M$) that can reside in any one of $n$ states.
It moves at every time step from state to state, and if it in presently
resides in state $i$ then with probability $p_{i,j}$ it will be
in state $j$ at the next time unit. $P=\left(p_{i,j}\right)$ is
called the transition matrix of the chain. $P$ is a stochastic matrix,
i.e. $\forall i,j$ $p_{i,j}\ge0$ each row sums to one ($\forall i\,\,\sum_{j}p_{ij}=1$).
\begin{definition}
We say that Markov chain is ergodic if the directed graph is strongly
connected and aperiodic.
\end{definition}


\subsubsection{informal chatter}

$P_{i,j}^{t}=$ probability that starting from state $i$ we will
be at state $j$ after $t$ units of time.

Let us restrict ourselves to deal with ergodic Markov chains.

Since the underlying directed graph of $M$ (markov chain) is strongly
connected and aperiodic the Perron -- Frobenius theorem applies.
(This matrix contains one positive eigenvector, other eigenvalues
are strictly smaller then the corresponding eigenvalue) $P\cdot\one=\one\Rightarrow\one$
is an eigenvector corresponding to eigenvalue $1$. By $P.F$ (Perron
-- Frobenius) $\one$ and the eigenvalue $1$ are the eigenvector
and eigenvalue pair given by the theorem $\Rightarrow$ all other
eigenvalues $\lambda$ of $p$ satisfies $\left|\lambda\right|<1$.

Let $u$ be the probability vector(where $u_{i}$ = is the probability
that the system starts at state $i$, and$u\ge0$,$\sum u_{i}=1$):
so 
\[
uP\one=u\one=\left\langle u,\one\right\rangle =1
\]
$\left(uP\right)_{\alpha}=\sum_{i}u_{i}p_{i\alpha}=$probability that
the chain is in state $\alpha$ at the next time unit.

$\left(uP^{2}\right)_{\beta}=\left(\left(uP\right)\cdot P\right)_{\beta}=$
prob that the chain will be in state $\beta$ in two units from now.

What do we want to know using Markov chains?
\begin{enumerate}
\item what is the long-run behavior of the chain.
\item How soon does it get there.
\end{enumerate}
%
From $P.F$ we know that there is a positive \uline{left} eigenvector
$\pi$ corresponding to eigenvalue $1$, $\pi P=\pi$. Let us normalize
$\pi$ so that $\left\langle \pi,\one\right\rangle =1$.
\begin{thm}[from Perron -- frobenious]
 $\limn{\left(\frac{A}{p}\right)^{n}}=\frac{x_{r}\otimes x_{\ell}}{\left\langle x_{r},x_{\ell}\right\rangle }$
where $\rho$ is the eigenvalue from perron--frobenious.$x_{r},x_{\ell}$
are the right and left eigenvectors of $\rho$.
\end{thm}

\begin{proof}
in the case of Markov Chains $A=P,\rho=1$. $x_{r}=\one,x_{\ell}=\pi$$\Rightarrow$
$\limn{\frac{P^{n}}{\rho^{n}}}=\one\otimes\pi=\begin{pmatrix}- & \pi & -\\
- & \pi & -\\
 & \vdots\\
- & \pi & -
\end{pmatrix}$. If $\limn{P^{n}}$ exists let us denote it by $Q$, thus $QP=Q$
$\Rightarrow$ Every row of $Q$ is a left eigenvector of $P$ with
eigenvalue $1$. %
\begin{insight}
(This gives us the matrix above)%
\end{insight}
. Why does the limit exist? The set of all $m\times n$ stochastic
matrices is compact %
\begin{insight}
This a simplex closed and bounded%
\end{insight}
$\Rightarrow$ every infinite sequence of stochastic matrices has
a convergent sub-sequence. As we saw only $Q$ can be a limit point
of matrices in the sequence. %
\begin{insight}
All stochastic matrices is a compact set and we have a sequence so
it has a convergent sequence. Every stochastic matrix will converge
and as a result we know that $Q$ is the only limit? something like
that.%
\end{insight}
\end{proof}
%

\subsection{Random Walks on Graphs}

$G=\left(V,E\right)$ an undirected graph. $States=V$. Transitions
matrix defined is normalized in row $i$ by $\frac{1}{\deg\left(v_{i}\right)}$.
\begin{proposition}
If $G$ is a connected non-bipartite undirected graph, then the limit
distribution of the random walk on $G$ is $\left(\cdots\frac{d\left(x\right)}{2\left|E\right|}\cdots\right)=\pi$.
$\pi P\overset{}{=}\pi$.
\end{proposition}


\subsubsection{Expander graphs}

Expander graphs walks on them, etc.

(informal definition) These are graphs with no bottlenecks. $G=\left(V,E\right)$
d-regular graph. $h\left(G\right):=\min\limits _{\begin{array}{c}
\emptyset\ne S\subseteq V\\
\left|S\right|\le\frac{\left|V\right|}{2}
\end{array}}\frac{e\left(S,\bar{S}\right)}{\left|S\right|}$ %
\begin{insight}
i think the function $e$ counts the number of edges between two sets
of nodes.

Why do we normalize by the size of $S$?%
\end{insight}

We want $d-regular$ graph with a large edge expansion.
\begin{definition}
(TAKEN FROM A PREVIOUS YEAR) Let $\delta\ge0$, a graph $G=\left(V,E\right)$
is $\delta$-edge expanding if for every $S$ of size at most $\frac{\left|V\right|}{2}$
it holds that $e\left(S,\overline{S}\right)\ge\delta\cdot\left|S\right|$.
\end{definition}

\begin{definition}
The edge expansion of a graph $h\left(G\right)$ is the minimal expansion
of a vertices set. I.e. $h\left(G\right)=\min\limits _{\begin{array}{c}
\emptyset\ne S\subset V\\
\left|S\right|\le\frac{\left|V\right|}{2}
\end{array}}\frac{e\left(S,\overline{S}\right)}{\left|S\right|}$
\end{definition}

\uline{Question:} Given $d\ge3$ does there exists $\delta>0$
and infinitely many $d-reg$ graphs $G=G_{n}$ %
\begin{insight}
familiy of $G_{n}$ for $n=d+1\to\infty$%
\end{insight}
with $h\left(G_{n}\right)\ge\delta$?

\begin{insight}
G=Gn - a family for the rest of these claims%
\end{insight}

\begin{thm}
\begin{insight}
\end{insight}
Yes, actually \uline{most} $d-reg$ graphs %
\begin{insight}
family%
\end{insight}
{} have this property.
\end{thm}

For your personal information: To decide whether $h\left(G\right)>\alpha$
is a $co-NP-HARD$ problem.
\begin{thm}
Let $G$ %
\begin{insight}
take a graph - not a family%
\end{insight}
{} be a $d-regular$ graph and let $d=\lambda_{1}\ge\lambda_{2}\ge\dots\ge\lambda_{n}$
be the eigenvalue of adjacency matrix. so %
\begin{insight}
$\Lambda=\max\lambda_{2},-\lambda_{n}$. This is what he wrote then
erased from the board.%
\end{insight}
{} %
\begin{insight}
Note that there must be a negative eigenvalues because the trace of
a graph's matrix is $0$ (on the diagonal is zeros)%
\end{insight}
{} $\frac{d-\lambda_{2}}{2}\le h\left(G\right)\le\sqrt{d^{2}-\lambda_{2}^{2}}$.
\end{thm}

Constructing expander graphs is qualitatively equivalent to constructing
$d-regular$ graphs with $\lambda_{2}$ bounded away from $d$. %
\begin{insight}
building this kind of graph familiy such that $h\left(G_{n}\right)$
is far from zero. equivalent to build a familty of $d-reg$ graphs
where $\lambda_{2}$ is bounded away from $d$. where bounded away
means that there is some small area around $d$ where lambda cannot
reach.%
\end{insight}

\begin{definition}
If $\sigma,\tau$ are two prob distributions on a finite set $S$,
we define their distance 
\[
d_{TV}\left(\sigma,\tau\right)=\sum_{x\in S}\left|\sigma\left(x\right)-\tau\left(x\right)\right|
\]

Please Note that in the recitation we defined it as follows:
\[
d_{TV}\left(\sigma,\tau\right)=\frac{1}{2}\sum_{x\in S}\left|\sigma\left(x\right)-\tau\left(x\right)\right|=\frac{1}{2}\norm{\sigma-\tau}_{1}
\]
\end{definition}

Let $G$ be a regular graph, Let $\sigma_{0}=\left(1,0,\dots,0\right)$,
$\sigma_{t}=$ distribution of the random walker at time $t$.

$Q:$ given $\epsilon>0$, for which time $t$ does it hold that $d_{TV}\left(\sigma_{t},uniform\right)<\epsilon$?
Let $\mathcal{G}=\left(G_{n}\right)$ be a family of $d-regular$
graphs, then the following are equivalent
\begin{enumerate}
\item $h\left(G_{n}\right)$ is bounded away from zero
\item $d-\lambda_{2}$ is bounded away from zero.
\item mixing time is $O\left(\log n\right)$
\end{enumerate}
Let $A$ be a positive matrix $\rho=$ spectral radius of $A$. $\limn{\frac{A}{\rho}}^{n}=\frac{x_{R}\otimes x_{L}}{\left\langle X_{R},X_{L}\right\rangle }$
where $x_{R},x_{L}$ are the right and left eigenvectors of $\rho$.
We will prove the case of $A$ a stochastic matrix.
\begin{proof}
$ $
\begin{lemma}
Let $G$ be a strongly connected aperiodic directed graph and let
$A$ be it's adjacency matrix. Then $\exists k$ s.t $A^{k}>0$. (Without
proof).
\end{lemma}

We will prove the case of a positive stochastic matrix $\then$ same
for the stochastic ergodic case.

Let $A$ be a positive stochastic matrix and let $\alpha>0$ be the
smallest entry in $A$. Let $v$ be a probability vector $v\ge0$,
$\left\langle v,\one\right\rangle =1$. Let $M,m$ be the largest
and smallest coordinates of $v$. Let $M',m'$ be the max/min of $Av$.
\begin{align*}
M' & \le\alpha m+\left(1-\alpha\right)M\\
m' & \ge\alpha M+\left(1-\alpha\right)m
\end{align*}
\begin{insight}
We are multiplying a row in $A$ by a vector. We engineer a vector
that sums to one and gets the biggest number. what do we do? alpha
will do the least damage when it is multiplied by $m$, now after
we do $\alpha m$ we add $\left(1-\alpha\right)M$ and we ignore all
other values, %
\end{insight}
{} so 
\[
M'-m'\le\left(1-2\alpha\right)\left(M-m\right)
\]
and $A^{n}v=A^{n-1}\left(Av\right)=A^{n-2}A\left(Av\right)$ $\then$
$M'^{\left(n\right)}-m'^{\left(n\right)}\le\left(1-2\alpha\right)^{n}\left(M-m\right)$
\[
\limn{M'^{\left(n\right)}-m'^{\left(n\right)}}=0
\]
 then 
\[
\limn{A^{n}v}=c\cdot\one
\]
 for some $c>0$. $v=e_{i}\then$ The col of $\limn{A^{n}}$ is a
constant vector. Therefore $\limn{A^{n}}=\begin{pmatrix}- & u & -\\
- & u & -\\
 & \vdots\\
- & u & -
\end{pmatrix}$ and we've seen that if there such limit it equals to $=\begin{pmatrix}- & \pi & -\\
- & \pi & -\\
 & \vdots\\
- & \pi & -
\end{pmatrix}$ %
\begin{insight}
(from 106)%
\end{insight}
. $c=\left\langle u,v\right\rangle =\left\langle \pi,v\right\rangle $
for each $e_{i}$, $c$ is different.
\end{proof}
%
Edge extension of $G$: $h\left(G\right):=\min\limits _{\begin{array}{c}
\emptyset\ne S\subseteq V\\
\left|S\right|\le\frac{\left|V\right|}{2}
\end{array}}\frac{e\left(S,\overline{S}\right)}{\left|S\right|}$. $G$ is $d-regular$ with eigenvalues $d=\lambda_{1}\ge\lambda_{2}\ge\dots\ge\lambda_{n}$.
\begin{insight}
The graph is connected iff $\lambda_{1}>\lambda_{2}$%
\end{insight}

\begin{thm}
\begin{insight}
page 42 in ``AllLectures''%
\end{insight}
Let $G$ be a $d-regular$ connected non-bipartite graph with eigenvalues
$d=\lambda_{1}\ge\lambda_{2}\ge\dots\ge\lambda_{n}$ then $\frac{d-\lambda_{2}}{2}\le h\left(G\right)\le\sqrt{d^{2}-\lambda_{2}}$.
\begin{proof}
we are only proving the bottom bound: $\lambda_{2}=\max\limits _{x\perp\one}\frac{xAx^{T}}{\norm x^{2}}$
\begin{insight}
$x\perp\one$ is beause the first eigenvector is $\one$ . and we
use the Rayleigh rits where $\lambda_{2}=\max\limits _{x\perp x_{1}}\frac{xAx^{T}}{\norm x^{2}}$%
\end{insight}
.We look at some partition of $V$: $S,\overline{S}$. What do we
know about the problem that we can use? We know that the graph is
$d-reg$ and we know the partition $S,\overline{S}$. So we will define
$x$ as follows:
\[
x_{i}=\begin{cases}
\left|\overline{S}\right| & i\in S\\
-\left|S\right| & i\in\overline{S}
\end{cases}
\]
this vector $x$ holds that $x\perp\one$ ($\left\langle x,\one\right\rangle =\sum x_{i}\cdot1=\left|\overline{S}\right|\cdot\left(-\left|S\right|\right)+\left|S\right|\cdot\left|\overline{S}\right|=0$)%
\begin{insight}
This is because we have $\left|\overline{S}\right|$ times $\left(-\left|S\right|\right)$
in entries in $x$ and $\left|S\right|$ times $\left|\overline{S}\right|$
entries in x.%
\end{insight}
. Now we need to calculate $\max\limits _{x\perp\one}\frac{xAx^{T}}{\norm x^{2}}$.
First: 
\begin{align*}
\norm x^{2} & =\sum_{i=1}^{n}x_{i}^{2}=\\
 & =\left|S\right|\left|\overline{S}\right|^{2}+\left|S\right|^{2}\left|\overline{S}\right|=\left|S\right|\left|\overline{S}\right|\left(\left|\overline{S}\right|+\left|S\right|\right)=n\left|S\right|\left|\overline{S}\right|
\end{align*}
 
\begin{align*}
xAx^{T} & =\sum a_{ij}x_{i}x_{j}=
\end{align*}
 this is like summing where $a_{ij}=1$ (where $i$ is a neighbour
of $j$) so 
\[
=\sum\limits _{i\sim j}x_{i}x_{j}=\underbrace{2\cdot e\left(S\right)\cdot\left|\overline{S}\right|^{2}}_{\left(1\right)}+\underbrace{2\cdot e\left(\overline{S}\right)\cdot\left|S\right|^{2}}_{\left(2\right)}-\underbrace{2e\left(S,\overline{S}\right)\left|S\right|\left|\overline{S}\right|}_{\left(3\right)}
\]
\end{proof}
\begin{enumerate}
\item $i,j$ both are vertices in $S$ hence $x_{i}\cdot x_{j}=\left|\overline{S}\right|\cdot\left|\overline{S}\right|$
and we count them $2\cdot e\left(S\right)$ times. 
\item same explanation as above but for $i,j$ vertices in $\overline{S}$.
\item $i\in S$, $j\in\overline{S}$, $x_{i}\cdot x_{j}=\left|\overline{S}\right|\cdot\left(-\left|\overline{S}\right|\right)$.
and we count them $2\cdot e\left(S,\overline{S}\right)$ times.
\end{enumerate}
Another important observation is $d\cdot\left|S\right|=2\cdot e\left(S\right)+e\left(S,\overline{S}\right)$
because this is a d-regular graph - we count for each vertex the number
of times it is touched by the the edges in $S$ plus the edges that
connect to$S$ from $\overline{S}$. (from the other side $d\left|\overline{S}\right|=2e\left(\overline{S}\right)+e\left(S,\overline{S}\right)$)
So %
\begin{insight}
\begin{align*}
xAx^{T} & =2\cdot e\left(S\right)\cdot\left|\overline{S}\right|^{2}+2\cdot e\left(\overline{S}\right)\cdot\left|S\right|^{2}-2e\left(S,\overline{S}\right)\left|S\right|\left|\overline{S}\right|=\\
* & =2\left(\frac{1}{2}\left(d\left|S\right|-e\left(S,\overline{S}\right)\right)\left|\overline{S}\right|^{2}+\frac{1}{2}\left(d\left|\overline{S}\right|-e\left(S,\overline{S}\right)\right)\left|S\right|-e\left(S,\overline{S}\right)\left|\overline{S}\right|\cdot\left|S\right|\right)=\\
* & =d\left|S\right|\left|\overline{S}\right|\left(\left|\overline{S}\right|+\left|S\right|\right)-e\left(S,\overline{S}\right)\left(\left|\overline{S}\right|^{2}+2\left|\overline{S}\right|\left|S\right|+\left|S\right|^{2}\right)=\\
 & =d\left|S\right|\left|\overline{S}\right|n-e\left(S,\overline{S}\right)n^{2}
\end{align*}
( {*} )were taken from allLectures.%
\end{insight}
{} 
\[
xAx^{T}=\dots=d\left|S\right|\left|\overline{S}\right|n-e\left(S,\overline{S}\right)n^{2}
\]

so $\lambda_{2}\ge\frac{xAx^{T}}{\norm x^{2}}$ lets assign what we
calculated: $\lambda_{2}\ge\frac{d\left|S\right|\left|\overline{S}\right|n-e\left(S,\overline{S}\right)n^{2}}{n\cdot\left|S\right|\left|\overline{S}\right|}$
since $\left|\overline{S}\right|\ge\frac{1}{2}n$ 
\[
\lambda_{2}\ge\frac{d\left|S\right|\left|\overline{S}\right|n-e\left(S,\overline{S}\right)n^{2}}{n\cdot\left|S\right|\left|\overline{S}\right|}\ge d-2\frac{e\left(S,\overline{S}\right)}{\left|S\right|}
\]
organize this a bit more to get: 
\[
\frac{e\left(S,\overline{S}\right)}{\left|S\right|}\ge\frac{d-\lambda_{2}}{2}
\]
We saw that $\forall S\,\;\left|S\right|\le\frac{n}{2}$ then $\frac{e\left(S,\overline{S}\right)}{\left|S\right|}\ge\frac{d-\lambda_{2}}{2}$.
This means that this is true for the minimum too $h\left(G\right)=\min\limits _{\begin{array}{c}
\emptyset\ne S\subseteq V\\
\left|S\right|\le\frac{\left|V\right|}{2}
\end{array}}\frac{e\left(S,\overline{S}\right)}{\left|S\right|}\ge\frac{d-\lambda_{2}}{2}$. 
\end{thm}

\begin{thm}
Fix $d$, let $G=G_{n}$ be an infinite family of $d-regular$ graphs.
then the following are equivalent:
\begin{enumerate}
\item $h\left(G_{n}\right)$ is bounded away from zero. %
\begin{insight}
\begin{lyxlist}{00.00.0000}
\item [{what}] is the meaning of bounded away from zero?
\end{lyxlist}
\end{insight}
\item The spectral gap $\lambda_{1}-\lambda_{2}$ is bounded away from zero. 
\item The random walk on $G_{n}$ convergence in logarithmic time.
\end{enumerate}
\begin{insight}
In terms of apllication this is a huge thing. in all things ofenumaration
- checking different models this means that these things a re practical.
this means we don't have to wait.%
\end{insight}
\begin{proof}
\textcolor{red}{I Had a big problem taking notes in this part, I'll
try and fix this proof as soon as possible} ($2\then3$, What was
mostly intuition and not formal)$A$ is the matrix of the graph (a
stochastic matrix) so $A=\frac{1}{d}A_{G}$. We know that 
\[
\frac{A}{d}^{N}\to\text{\ensuremath{\begin{pmatrix}\\
 & \frac{1}{n}\\
\\
\end{pmatrix}}}
\]
Now assume we have some distribution, and we want to understand $\left(\frac{A_{G}}{d}\right)^{N}v\to\frac{1}{n}\one$.
We know nothing about $v$. We will express $v$ with the basis of
eigenvectors. $v=c_{1}\cdot\one+\sum c_{i}v_{i}$. where $\frac{A_{G}}{d}v_{i}=\lambda_{i}v_{i}$.
so what is $c_{1}$? $c_{1}=\left\langle v,\one\right\rangle $ and
because $v$ is a probability vector $c_{1}=\frac{1}{\sqrt{n}}$ .
thus $\frac{d}{\sqrt{n}}$ normalizes what happen when we apply $A^{N}$
to $v$? $A^{N}v=c_{1}\cdot A\one+\sum c_{i}Av_{i}$something that
goes exponentially fast to zero. And we stopped here.%
\begin{insight}
\begin{proof}
We are proving that $2\then3$.
\[
\left(\frac{A}{d}\right)^{n}v\to\frac{1}{n}\one
\]
we want to write $v$ as $v=c_{i}\frac{1}{\sqrt{n}}\one+\sum c_{i}v_{i}$.
where $c_{i}=\left\langle v,v_{i}\right\rangle $. Because $v$ is
a prob vector $c_{1}=\frac{1}{\sqrt{n}}..$ when $A^{n}v$ we get
that $\one$ stays the same, and all the other vectors other then
$\one$ tends to zero.
\end{proof}
Question given $d$ how large can the spectral be? $\iff$how small
can $\lambda_{2}$ be?%
\end{insight}
\end{proof}
\end{thm}

Question: given $d$ how large can the spectral gap be? $\iff$ how
small can $\lambda_{2}$ be?
\begin{thm}[Alon Boppana bound]
 %
\begin{insight}
Pg 109 in Ayelet's sum%
\end{insight}
$\lambda_{2}\ge2\sqrt{d-1}-o_{n}\left(1\right)$ .

a hint was given but not an actual proof.
\end{thm}

\begin{thm}
$\lambda_{2}\ge\sqrt{d}$
\begin{proof}
Let $A$ be the adjacency matrix of $G$%
\begin{insight}
G is d regular and have eigenvalues and stuff%
\end{insight}
. Then $tr\left(A^{2}\right)=\sum\lambda_{i}^{2}$, $\lambda_{1}=d$
so the first in the sum is $\lambda_{1}^{2}=d^{2}$ so $\sum_{i=1}\lambda_{i}=d^{2}+\sum_{i=2}\lambda_{i}^{2}$
. We know that $\lambda_{2}\ge\lambda_{3}\ge\lambda_{4}\dots\ge\lambda_{n}$
thus $\sum_{i=1}\lambda_{i}=d^{2}+\left(n-1\right)\cdot\lambda_{2}^{2}$
on the other hand 
\[
tr\left(A^{2}\right)=n\cdot d
\]
because $A^{2}$ in the $i,j$ entry counts how many paths of size
2 there are between $i$to $j$ ($d$ paths). so there are $n$ times
$d$ paths. so
\[
tr\left(A^{2}\right)=nd\le d^{2}+\left(n-1\right)\lambda_{2}^{2}
\]
 we conclude: $d-o\left(d\right)=\frac{nd-d^{2}}{n-1}\le\lambda_{2}^{2}$
\end{proof}
\end{thm}


\part{Optimization}

\section{basics}

An optimization problem is defined by a domain $\Omega$ and $f:\Omega\to\R$.
$\max\limits _{x\in\Omega}f\left(x\right)$.

\begin{insight}
$\Omega$ can be all the spanning trees in a graph.%
\end{insight}

\begin{definition}
Linear programming $\Omega\subseteq\R^{n}$ is defined by a finite
list of linear equations and inequalities, and $f$ is linear: $x\to\left\langle c,x\right\rangle $.
\begin{insight}
We can look at bit omega as a space that we cut from $\R^{n}$ with
planes and halfspaces. forexample in $\R^{2}$ we cut the space with
linear lines.The resulting omega will be some ``polygon''.%
\end{insight}
\begin{insight}
What does it mean to optimize a linear function? we want to maximize
$\left\langle c,x\right\rangle =\gamma$ we actually looking at the
last lines that still stays in the ``polygon''.%
\end{insight}
\end{definition}

$H$ set of the form $H=\left\{ x\in\R^{n}|\left\langle a,x\right\rangle =\alpha\right\} $
is called a hyperplane. (in $\R^{2}$ it is a line, in $\R^{3}$ it
is a plane). == affine subspace of dimension $n-1$.

corresponding to $H$ are two open half spaces $H^{+}=\left\{ x\in\R^{n}\big|\left\langle a,x\right\rangle >\alpha\right\} $,
$H^{-}=\left\{ x\in\R^{n}\big|\left\langle a,x\right\rangle <\alpha\right\} $
.
\begin{itemize}
\item The $\Omega$s that will interest us, are, therefore intersections
of a finite number of half-spaces. Such sets are called polyhedra
(plural of polyhedron)
\end{itemize}
A bounded polyhedron is called a polytope (peon in Hebrew).

LP is, then, the problem of maximizing a linear function over a polytope.

\subsection{feeding cows}

you can feed a cow

\[
A=\begin{array}{ccccc}
 & \text{protein} & \text{fat} & \text{vitimin A} & \dots\\
\text{grass} & 7\\
\text{grains}\\
\vdots
\end{array}
\]

\begin{itemize}
\item $0\le x_{i}\to$ How many units of each food type $i$ do we give
the cow.
\item $c_{i}=$ cost of one unit of food type $i$.
\item $b_{j}$ is the minimum amount of component $j$ (a food nutrition
like protein) that the cow needs
\end{itemize}
The farm want to find a non negative vector $x$ indicating from each
type of food what to give to the cow such that: $\min\limits _{x}\left\langle c,x\right\rangle $
(get $x$ for the least possible price) and $Ax\ge b$ (the cow receives
enough of every nutrient).

\section{examples}

\subsection{separating points}

\subsubsection{linear separation}

Given $P\subseteq\R^{n}$ finite, $N\subseteq\R^{n}$ finite.(Assume
we have two groups of points and we denote $P$ as the points with
label $+$ and $N$ the points with the negative label)

We have a few questions regarding these groups:
\begin{enumerate}
\item Is there a hyperplane that separates $P$ from $N$? $\exists H\quad s.t\;P\subseteq H^{+},N\subseteq H^{-}$
?
\begin{itemize}
\item Does there exist $\left(c_{1}\dots c_{n}\right)=c$ such that $\forall j\left\langle c,v_{j}\right\rangle >0$,
$\forall i\left\langle c,u_{i}\right\rangle <0$ where $u_{i}\in P$
and $v_{j}\in N$.
\end{itemize}
\item If such hyperplane exists how large can it's margin be?(The distance
from the separating line and the closest point)
\begin{itemize}
\item This is similar to $1$ but we want to find $\epsilon$ subjected
to $\forall i\left\langle c,v_{i}\right\rangle \ge\epsilon$ and $\forall j\left\langle c,u_{j}\right\rangle \le-\epsilon$ 
\end{itemize}
\item Feasibility: given a system of linear, equations and inequalities,
is there an $x$ that satisfies all of them?
\begin{itemize}
\item (feasibility and optimizations are highly the same question 1 is a
feasibility question. )%
\begin{insight}
When we want to optimize we first need to know if it is even possible
to answer the question.%
\end{insight}
{} 
\end{itemize}
\end{enumerate}
\begin{insight}

\subsubsection{answer to question1}

do there exist $\left(c_{1}\dots c_{n}\right)=c$ such that $\left\langle c,v_{j}\right\rangle >0$
$\forall j\,\,\,\left\langle c,u_{i}\right\rangle <0\forall i$. $u_{i}\in P$,$v_{j}\in N$.
we want $H=\left\{ x\big|\sum c_{i}x_{i}=1\right\} $. So we want
\[
\max_{\begin{array}{c}
\forall i\left\langle c,v_{i}\right\rangle \ge\epsilon\\
\forall j\left\langle c,u_{j}\right\rangle \le-\epsilon
\end{array}}\epsilon
\]
\end{insight}


\subsubsection{Polynomial separation}

Assume we have points from two groups But they sit in a space that
we cannot split the groups with a straight line, rather we need a
polynomial like $y=ax^{2}+bx+c$. This is not a linear function. how
can we use linear programming to get such a function

The points in the group $+$ are $\left(x_{i},y_{i}\right)$ for $i=1\dots n$.
The points in group $-$ will be denoted as $\left(u_{i},v_{i}\right)$
for $i=1\dots m$. The pluses should be below the polynomial we choose.
Writing this formally we would want the following:
\begin{itemize}
\item $\forall i=1\dots n\,\;y_{i}\le ax_{i}^{2}+bx_{i}+c$
\item $\forall i=1\dots m\,\;v_{i}\ge au_{i}^{2}+bu_{i}+c$
\end{itemize}
We can also ask for the biggest margin -- we want to maximize some
distance $\epsilon$ subject to these constraints. Stating this formally:
\begin{itemize}
\item $\forall i=1\dots n\,\;y_{i}+\epsilon\le ax_{i}^{2}+bx_{i}+c$
\item $\forall i=1\dots m\,\;v_{i}-\epsilon\ge au_{i}^{2}+bu_{i}+c$
\end{itemize}
You should notice that these constraints are still linear. That is
$a,b,c$ are still linear, and $u_{i}$ and $x_{i}$ are simply constant
numbers. Hence each constraint is linear in the variables. Now pushing
these constraints into a linear programming solver should output some
$a,b,c$ coefficients to get a polynomial of the form $y=ax^{2}+bx+c$.

\begin{insight}
A question was why use only $\epsilon$ in the y-axis. A better way
is to show that the distance is an actual distance and not only on
the y-axis as we show, but that is a lot more complicated.%
\end{insight}


\subsection{Humongous circle }

\begin{insight}
get the drawwing about 9 mins to the lecture%
\end{insight}
remember that a polyhedron is defined by a list of lines. For every
$i=1\dots n$ we have $y=a_{i}x+b_{i}$, consider a polygon that lies
above $\ell_{1}\dots\ell_{k}$ and below $\ell_{k+1}\dots\ell_{n}$.
What is the largest circle we can draw inside this polygon? Linear
programming can help us solve this question too. \\
Our variables are going to be: $\underbrace{\left(u,v\right)}_{\text{circle center}},\underbrace{R}_{\text{radius}}$.
Our objective function is $\max R$. The distance between the point
$\left(u,v\right)$ and the line $\ell=y=ax+b$ is $\frac{\left|v-au-b\right|}{\sqrt{1+a^{2}}}$.
We want to find for each line $\ell$ that its distance to $\left(u,v\right)$
is bigger or equal then $R$ (otherwise the circle isn't bounded by
the polygon). Let us define the problem with the tools of linear programming:
$\max R$ is subjected to $\forall i=1\dots n$ 
\[
\frac{v-a_{i}u-b_{i}}{\sqrt{1+a_{i}^{2}}}\le R
\]

\[
\frac{v-a_{i}u-b_{i}}{\sqrt{1+a_{i}^{2}}}\ge-R
\]
Notice that we dropped the absolute value and changed it into two
inequalities -- why? Because absolute value isn't linear. In this
problem too the constraints are still linear. ($u,v,R$ are the variables
and the rest is just numbers from the lines definitions)

\subsection{Interpolation}

input: $\left(x_{1},y_{1}\right),\left(x_{2},y_{2}\right),\dots,\left(x_{n},y_{n}\right)$.
We want to find a line $y=ax+b$ such that $\sum\limits _{i=1}^{n}\left(ax_{i}+b-y_{i}\right)^{2}$
(least squares method) is minimal. This is not linear! we cannot solve
this using linear programming :(

least squares is sensitive to outliers -- we square the error, and
introduce the least squares method to this sensitivity. One way to
mitigate the issue of outliers is to use this formula: $\sum\limits _{i=1}^{n}\left|ax_{i}+b-y_{i}\right|$.
this isn't yet linear too.

What we will do instead is to minimize $e_{1}+e_{2}+\dots+e_{n}$
($e_{i}$ should be equal to $\left|ax_{i}+b-y_{i}\right|$ -- Michael
will convince us that this will happen) subject to:

$\forall i=1\dots n\quad ax_{i}+b-y_{i}\le e_{i}$, $ax_{i}+b-y_{i}\ge-e_{i}$,
and $e_{1}\dots e_{n}\ge0$. In any feasible solution $a,b,e_{1},\dots,e_{n}$:
$\sum_{i=1}^{n}e_{i}\ge\sum_{i=1}^{n}\left|ax_{i}+b-y_{i}\right|$.
\begin{insight}
basically he minimize it such that $e_{i}$ is a variable that by
minimization becomes the value of the norm 1. pretty cool! if we find
$e_{i}$ they must equal because the constraints force them to do
that!%
\end{insight}
{} But in optimal solution we have equality, because if we get $e_{i}>\left|ax_{i}+b-y_{i}\right|$
then just take the solution $e_{i}'=\left|ax_{i}+b-y_{i}\right|$
and find a solution that is smaller, and still follows the constraints.

\subsection{Integer linear programming (ILP)}

\begin{insight}
This is similar to linear programming but we can only have an integer
solutions %
\end{insight}
An integer linear program is: $\max c^{T}x$ subject to $Ax\le b$
and $x\in\Z^{n}$. Linear program can be solved in linear time. Integer
linear programing is $NP-HARD$.

\subsubsection{reducing $3-SAT$ to ILP:}

Suppose we have $\bigwedge_{i=1}^{m}\left(\ell_{i}^{1}\lor\ell_{i}^{2}\lor\ell_{i}^{3}\right)$
each $\ell_{i}^{j}$ is either $x_{i}^{j}$ or its negation $\overline{x_{i}^{j}}$.
We will write a linear program that is feasible if and only if $3-SAT$
is feasible: $\forall i=1\dots m$ $Z_{i}^{1}+Z_{i}^{2}+Z_{i}^{3}\ge1$
$\forall i,j$ such that $\ell_{i}^{j}=x_{i}^{j}$ then $Z_{i}^{j}=x_{i}^{j}$
.$\forall i,j$ such that $\ell_{i}^{j}=\overline{x_{i}^{j}}$ then
$Z_{i}^{j}=1-x_{i}^{j}$.\.{ }$\forall_{i,j}x_{i}^{j}\in\left\{ 0,1\right\} $
\begin{insight}
For each literal we introduce a new x variable.%
\end{insight}
. This definition shows us that if we want to solve $3-SAT$ we can
solve the linear programming or the other way around - they are equivalent.
$x_{i}^{j}$ is true if and only if $x_{i}^{j}$ equals 1 in the linear
programming solution.

\begin{insight}
x has two indices because there%
\end{insight}


\subsubsection{Maximum weight matching:}

A set of $n$ workers, $m$ tasks (each worker can do only one task
and each task can be done by one worker) to each task a benefit $c_{i,j}\in\R$
(worker $i$ task $j$).want: find a matching between workers and
tasks that maximizes total benefit. 
\[
\max\sum_{i,j}x_{ij}c_{ij}
\]
 subjected to
\begin{itemize}
\item $\forall i=1\dots n\quad x_{i1}+x_{i2}+\dots+x_{im}\le1$ each worker
have only one task.
\item $\forall j=1\dots m\quad x_{1j}+x_{2j}+\dots+x_{nj}\le1$ each task
has only one worker.
\item $x_{i,j}\in\left\{ 0,1\right\} $ the matching must be true or false.
\end{itemize}
The actual problem fits to $ILP$ like a glove, but we can try and
solve this using regular linear programming where $x_{i,j}\in\left[0,1\right]$
rather than $x_{i,j}\in\left\{ 0,1\right\} $.

\section{Formalizing linear programming}
\begin{definition}
A linear program is in \uline{canonical} form $\max c^{T}x$ subject
to $Ax\le b$
\end{definition}

\begin{definition}
The \uline{standard} form for linear programming: $\max c^{T}x$
subject to $AX=b$ and $x\ge0$. %
\begin{insight}
the advantage of this form is that it's easier to prove things about
it. we will show how to move from canonical form to standard form.%
\end{insight}
\end{definition}

how do we move from canonical form to standard form? (This can be
done efficiently)

First we transform all inequalities to equalities:
\begin{itemize}
\item $\forall i,$ replace $\left\langle a_{i},x\right\rangle \le b_{i}$
($a_{i}$ is the $i$th row in $A$) with $\left\langle a_{i},x\right\rangle +s_{i}=b_{i}$.
\item For every $x_{i}$, replace it with $y_{i}-z_{i}$. ($y_{i},z_{i}$
are two new variables we introduce).%
\begin{insight}
the difference between two numbers can be any number we want.%
\end{insight}
{} %
\begin{insight}
We need the variables to be non negative, so what we did is that every
real number can be expressed between two negative numbers. meaning
now we only have non negative variables. nice.%
\end{insight}
\end{itemize}
\begin{insight}
so now we maximise as follows: $\max c^{T}\left(y-z\right)$ subjected
to $Ax=\left(b-s\right)$%
\end{insight}

From now on to the end of the lecture we talk on $LP$ problems in
standard form.

\subsubsection{basic feasible solution}

We have $Ax=b$, we will assume that $A\in M_{m\times n}\left(\R\right)$
where $m\le n$ and $rank\left(A\right)=m$ %
\begin{insight}
because otherwise there is no solution right?%
\end{insight}
. 
\begin{definition}
The set of feasible solutions is $P=\left\{ x\in\R^{n}:Ax=b,x\ge0\right\} $
\end{definition}

\begin{definition}
For $x\in P$ the set of feasible solutions, $support\left(x\right)=\left\{ i\in\left[n\right]:x_{i}\ne0\right\} $. 
\end{definition}

For a subset of $B\subseteq\left\{ 1\dots n\right\} $ $A_{B}$ is
the sub-matrix of $A$ with columns indexed by $B$. example of $A=\begin{pmatrix}1 & 2 & 3\\
4 & 5 & 6
\end{pmatrix}$ indexed by $1,3$: $A_{\left\{ 1,3\right\} }=\begin{pmatrix}1 & 3\\
4 & 6
\end{pmatrix}$.
\begin{definition}
$x\in P$ ($P$ is the set of feasible solution) is a \uline{basic
feasible solution} (BFS) if $\exists B\subseteq\left\{ 1\dots n\right\} $
such that:
\begin{enumerate}
\item $B\supseteq support\left(X\right)$
\item $A_{B}$ is non-singular. (invertible) %
\begin{insight}
in particular $A_{B}$ must be square to be invertible. (so $A_{B}$
must be of size $m\times m$ if youre not sure why check how $A_{B}$
is defined - it won't be square if $B$ is not of size m.)%
\end{insight}
\end{enumerate}
\end{definition}

\begin{lemma}
$x\in P$ is basic iff the columns of $A_{support\left(x\right)}$
are linearly independent.
\begin{proof}
$ $
\end{proof}
\begin{itemize}
\item[$\then$] Assume that $x$ is basic: then $\exists B\supseteq support\left(x\right)$
such that $A_{B}$ is non-singular, $A_{B}$ contains $A_{support\left(x\right)}$
as a sub matrix so it's columns are linearly independent.
\item[$\Leftarrow$] Assume that $A_{support\left(x\right)}$ has linearly independent
columns. By assumption $rank\left(A\right)=m$ so we can extend $support\left(x\right)$
to a set $B\supseteq support\left(x\right)$ such that $A_{B}$ is
non-singular.
\end{itemize}
\end{lemma}

\begin{thm}
\label{thm:consider-the-:LP126}consider the $LP$: $\max c^{T}x$
subject to $Ax=b,x\ge0$.
\begin{enumerate}
\item If there is a feasible solution and the objective function is bounded
from above, then there is an optimal solution.
\item If an optimal solution exists then there is an optimal solution that
is also basic.
\end{enumerate}
\end{thm}


\subsubsection{Algorithm for LP:}

Suppose that $B\subseteq\left\{ 1\dots n\right\} $ such that $A_{B}$
is non-singular. Then there exists a unique $x_{B}=A_{B}^{-1}b$.
Every $BFS$ has this form because if $x$ is a $BFS$ then $\exists B\supseteq supp\left(x\right)$
such that $A_{B}$ non-singular. $Ax=b\then\sum\limits _{i\in supp\left(x\right)}a_{i}x_{i}=b$
($a_{i}$ is the $i$th column of $m$) $\then$ $\underbrace{x}_{x_{B}}=\underbrace{A_{B}^{-1}b}_{\in\R^{m}}$.
$x_{B}=\begin{cases}
x_{i} & i\in B\\
0 & i\notin B
\end{cases}$. The algorithm: given matrix $A$ look at all non-singular matrices
and compute $x_{B}=A_{B}^{-1}b$. (This is really not efficient)
\begin{proof}[proof for \thmref{consider-the-:LP126}]
 Suppose $x_{0}\in P$. We will find a BFS $x\in P$ such that $c^{T}x\ge c^{T}x_{0}$
(Note that we seek maximal solutions). Let $S=\left\{ v\in P:c^{T}v\ge c^{T}x_{0}\right\} $
This is not an empty set because $x_{0}$ is in that set. Let $x\in S$
have the smallest number of non-zero entries (smallest possible support)
We will show that $A_{supp\left(x\right)}$ has linearly independent
columns (and so $x$ is a BFS). Suppose, for a contradiction that
$A_{supp\left(x\right)}$ does not have linearly independent columns
$\then$ $\exists0\ne y\in\R^{n}$ such that $Ay=0$, and $supp\left(y\right)\subseteq supp\left(x\right)$
(We take a non zero independent column from $A_{supp\left(x\right)}$
and extend it with zeros)%
\begin{insight}
The support of $y$ is contained because $y$ is defined such that
$A_{supp\left(x\right)}$ is linearily independant and $y$ is the
product of that%
\end{insight}
.We can say that $\forall t\in\R$, $A\left(x+ty\right)=Ax=b$. We
can use $t$ to change one of $x's$ values to zero (because $supp\left(y\right)\subseteq supp\left(x\right)$).
So for small enough $t$ we get $A\left(x+ty\right)=b$ and also $x+ty\ge0$
(because $supp\left(y\right)\subseteq supp\left(x\right)$). \\
Suppose that $c^{T}y\ge0$ and also $\exists j$ (explanation why
is provided by the end of the proof), such that $y_{j}<0$ which means
$j$ is in the support of $x$ this implies that $x_{j}>0$( we are
in standard form). If this happens let $t>0$ be the largest value
such that $x+ty\ge0$ so $support\left(x+ty\right)\subsetneq supp\left(x\right)$
(there is at least one more zero in the entries of $x+ty$). 
\[
c^{T}\left(x+ty\right)=\underbrace{c^{T}x}_{\ge c^{T}x_{0}}+t\underbrace{c^{T}y}_{\ge0}\ge c^{T}x_{0}
\]
This is a contradiction (we found a vector $x+ty$ which has a smaller
support than $x$ -- contradiction). We are good as long as we can
assume $c^{T}x$ and that $y_{j}<0$ how can we assume that? \\
We can always assume that $c^{T}y$ is non negative. If $c^{T}y<0$
then replace $y$ with $-y$. Now we may assume that $c^{T}y\ge0$.
We need to show why $y_{j}<0$ for some $j$:
\begin{itemize}
\item If $c^{T}y\ge0$ and $y$ doesn't have negative coordinates then $\forall t\ge0$
$x+ty\ge0$. So if $c^{T}y>0$ then $\lim\limits _{t\to\infty}c^{T}\left(x+ty\right)=\infty$
which is a contradiction to the bounding of the objective. We know
need to address $c^{T}y=0$:
\item If $c^{T}y=0$ and $y\ge0$($y$ is not trivial) then take $-y\le0$
($-y$ has a negative coordinate because $y$ is not trivial), $c^{T}\left(-y\right)\ge0$
and we are back to the first assumption which gave us what we wanted.
\end{itemize}
\end{proof}
%
We showed that Its enough to look at basic feasible solutions.

\begin{insight}

\subsection{polyhedrons}

$\max c^{T}x$ $Ax=b$ $x\ge0$.

What we should do is inspect a vertex and then inspect it's neighbours
if we cannot improve by hoping -> were done.%
\end{insight}

\begin{insight}
We are showing again some of the properties we've seen with Michael.

Linear programming is the problem of maximizing a linear function
over a domain that is defined by finitely many linear equations and
inequalities. (Polyhedron and if it is bounded it is known as a polytope).
\begin{definition}
Canonical form of an LP:
\begin{align*}
\max\left\langle c,x\right\rangle \\
 & Ax\le b
\end{align*}
\end{definition}

If we want to reduce a linear program to the canonical form we need
to come up with something that will remove equations: 
\[
\left\langle a,x\right\rangle =\alpha\iff\left\langle a,x\right\rangle \ge\alpha\text{ and }\left\langle a,x\right\rangle \le\alpha
\]

\begin{definition}
Standard form of an LP: 
\begin{align*}
\max\left\langle c,x\right\rangle \\
 & Ax=b\\
 & x\ge0
\end{align*}
\end{definition}

We only look at standard form.%
\end{insight}

Some reminders:
\begin{definition}
Canonical form of an LP:
\begin{align*}
\max\left\langle c,x\right\rangle \\
 & Ax\le b
\end{align*}

Standard form of an LP: 
\begin{align*}
\max\left\langle c,x\right\rangle \\
 & Ax=b\\
 & x\ge0
\end{align*}
If we want to reduce a linear program to the canonical form we need
to come up with something that will remove equations: 
\[
\left\langle a,x\right\rangle =\alpha\iff\left\langle a,x\right\rangle \ge\alpha\text{ and }\left\langle a,x\right\rangle \le\alpha
\]
\end{definition}

\begin{definition}
$x$ is a feasible solution if it satisfies $Ax=b$ and $x\ge0$
\end{definition}

\begin{definition}
$support\left(x\right)=\left\{ i\big|x_{i}>0\right\} $%
\begin{insight}
This is equivalent because because $x\ge0$ in the standard form%
\end{insight}
\end{definition}

\begin{definition}
$x$ is a basic feasible solution (BFS) if the set of columns $\left\{ a_{i}\big|i\in supp\left(x\right)\right\} $
is a linearly independent set.
\end{definition}

\begin{thm}
If the system $Ax=b,x\ge0$ is feasible (there is a non negative $x$
such that $Ax=b$) and if $\left\langle c,x\right\rangle $ is bounded
from above $\then$ There exists an optimal solution to the LP 
\begin{align*}
\max\left\langle c,x\right\rangle \\
 & Ax=b\\
 & x\ge0
\end{align*}
 that is a basic feasible solution.
\end{thm}

\begin{insight}
If you haven't watched the lecture, help yourself to minute 13:19
of lecture dated at 141.1.2020 ``At least you have a tree, you can
hang yourself from it.'' biggest joke of the century.%
\end{insight}


\subsubsection{Talking about simplex algorithms}

Some intuition:

In this table the rows represent equivalent things:

\begin{tabular}{|c|c|}
\hline 
Geometry & Linear algebra\tabularnewline
\hline 
\hline 
vertex of the polytope & BFS\tabularnewline
\hline 
supporting hyperplane & defining inequality\tabularnewline
\hline 
an edge $=$ two neighboring vertices & two neighboring bases\tabularnewline
\hline 
\end{tabular}

\begin{insight}
$\iff$ A linear function defined on a polytope takes it's maximum
at a vertex. (if $x$ is basic then $x$ is a vertex in the polytope)

moving from neighbour to another neighbour in linear algebra we have
some basis, we take another vector -- add it to to the basis and
take some other column vector out of the basis of the matrix. (like
in the geometric intuition)%
\end{insight}

If a solution is locally optimal then we are in an optimum. (We work
in a convex world so if we reach some local optimum we reached a global
one). A linear function defined on a polytope takes it's maximum at
a vertex. (if $x$ is basic then $x$ is a vertex in the polytope).
So if we choose a vertex and we can't improve by moving to some of
our neighboring vertices we got an optimum. We can assume that moving
from neighbor to another neighbor in linear algebra is done as follows:
Choose some base of column from $A$ which are independent. Form a
basis of the column space from these vectors. Add a vector (which
breaks the independence) and remove some vector to gain independent
basis back.

\begin{insight}
Many versions of the simplex were shown to take exponential time at
worst case, we still dont know how good they are.%
\end{insight}

A lot of simplex versions do not reach a solution in polynomial time.
Associated with every polytope there is a graph. If the graphs diameter
is big the simplex algorithm might take a long time reaching a solution. 

In practice we can assume that LP is a solved problem and use it.

\subsubsection{coming back to feasible questions}

We don't always want to maximize. When we ask questions such as 
\begin{align*}
\max\left\langle c,x\right\rangle \\
 & Ax\le b\\
 & x\ge0
\end{align*}
We might as well ask: is $Ax\le b,x\ge0$ even feasible? without maximizing
we can search: 
\begin{align*}
Ax & \le b\\
x & \ge0\\
\left\langle c,x\right\rangle  & \ge\text{some number}
\end{align*}
We seek feasibility of $Ax\le b$, and $\left\langle c,x\right\rangle $
with some number. %
\begin{insight}
if we can do it - we enlarge that by some factor and try again. if
we are bounded we can go down.%
\end{insight}

\begin{insight}
The above part is a way to write this in the sum: When we ask question's
such as 
\begin{align*}
\max\left\langle c,x\right\rangle \\
 & Ax\le b\\
 & x\ge0
\end{align*}
 is $Ax\le b,x\ge0$ even feasible? without maximizing we can do a
simply search: 
\begin{align*}
Ax & \le b\\
x & \ge0\\
\left\langle c,x\right\rangle  & \ge\text{some number}
\end{align*}
 we can do a binary search to find it. we seek feasibility of $Ax\le b$,
and $\left\langle c,x\right\rangle $ with some number. if we can
do it - we enlarge that by some factor and try again. if we are bounded
we can go down.%
\end{insight}


\section{Linear programming duality}

building up some intuition:
\begin{enumerate}
\item Any linear combinations of linear equations and non-negative combinations
of linear inequalities is valid.%
\begin{insight}
\begin{align*}
3x+5y-z & =3\\
2x-3y+7z & =8\\
 & \Downarrow\\
5x+2y+6z & =11
\end{align*}
\end{insight}
\item This is \uline{all} you can derive from the above data. you can
square, you can take log, you can try many things but there is nothing
more you can derive.
\end{enumerate}
%
Let's say we have 
\begin{align*}
\max\left\langle c,x\right\rangle \\
 & Ax\le b\\
 & x\ge0
\end{align*}
Lets rephrase this. What is the largest $T$ for which the system
$Ax\le b,x\ge0,\left\langle c,x\right\rangle \ge T$ is feasible.

Assuming the intuition from above, this means that if for some given
$T$ 
\begin{align*}
Ax & \le b\\
x & \ge0\\
\left\langle c,x\right\rangle  & \ge T
\end{align*}
 is infeasible, then we should be able to derive $0>1$ using only
allowed linear combinations.
\begin{thm}[Strong LP Duality]
\label{thm:DualLP}If the following LP is bounded 
\begin{align*}
\max\left\langle c,x\right\rangle \\
 & Ax\le b\\
 & x\ge0
\end{align*}
 Then the Dual LP (DLP) 
\begin{align*}
\min\left\langle y,b\right\rangle \\
 & yA\ge c\\
 & y\ge0
\end{align*}
\begin{insight}
$b$ and $c$ changed rows, larger is replaced with smaller , the
only thing that remains is the non negativity.%
\end{insight}
is also bounded and the moreover the optima are equal. Moreover every
$x,y$ feasible for LP and DLP respectively satisfy that $\left\langle c,x\right\rangle \le\left\langle b,y\right\rangle $.
\end{thm}

Let $x$ be a feasible solution for LP. Let $y$ be a feasible solution
for DLP. Let us calculate $yAx$. \label{weak-LP-duality-proof}

$Ax\le b,y\ge0$ so $yAx\le\left\langle y,b\right\rangle $. %
\begin{insight}
multiply from the left by $y^{T}$: $y^{T}Ax\le y^{T}b=\left\langle y,b\right\rangle $%
\end{insight}
on the other side $yA\ge c,x\ge0$ so $\left\langle c,x\right\rangle \le yAx$.
we get 
\[
\left\langle c,x\right\rangle \le yAx\le\left\langle y,b\right\rangle \then\left\langle c,x\right\rangle \le\left\langle y,b\right\rangle 
\]


\subsection{Feeding a dual cow}

We are back to the cow feeding problem

cow feed table:

\[
A=\begin{array}{ccccc}
 & \text{protein} & \text{fat} & \text{vitimin A} & \dots\\
\text{grass} & 7 & \cdot & \cdot & \cdot\\
\text{grains} & \cdot & \cdot & \cdot & \cdot\\
\vdots & \vdots &  & \vdots & \ddots
\end{array}
\]
 We will work with $A^{T}$ :
\begin{itemize}
\item The LP:
\begin{align*}
xA & \ge b\\
x & \ge0\\
\min & \left\langle c,x\right\rangle 
\end{align*}
 the wishes of the farmer - minimal price a farmer wants to pay on
feeding his cows.
\item DLP 
\begin{align*}
Ay & \le c\\
y & \ge0\\
\max & \left\langle b,y\right\rangle 
\end{align*}
Optimization problem of the producer of pills. The producer wants
to generate artificial food for cows and to check how much to charge
their costumers. $y$ is the vector prices per pill. $\max\left\langle b,y\right\rangle $
captures the price the farmer needs to pay if he switches to pills.
$Ay\le c$ express the requirement that this way of feeding is price-competitive
with the traditional ways of feeding cows.
\end{itemize}
\begin{insight}
the problem of the dualist%
\end{insight}
{} 

\thmref{DualLP} states that we can solve each problem by looking
at the dual problem, which means that the pill maker should be able
to solve the farmers problem and get the maximal price he can try
and squeeze out of farmers.
\begin{lemma}[Farkas]
 Let $A_{m\times n}$ be a real matrix and let $b\in\R^{m}$. Then
the system $Ax=b$, $x\ge0$ has no solution (i.e infeasible) iff
there exists a vector $y\in\R^{m}$ such that $yA\ge0$ and $\left\langle b,y\right\rangle <0$.
\end{lemma}

\begin{proof}
One part of the proof is very easy: If such $y\in\R^{m}$ exists,
let us multiply by it on the left: $yAx=\left\langle b,y\right\rangle <0$
and we know $yAx=\left\langle yA,x\right\rangle \ge0$ because $yA\ge0$
and $x\ge0$. %
\begin{insight}
Page 54 AllLectures%
\end{insight}
{} We will be back here later. The full proof will be in another part
of the summary (as was done in class)
\end{proof}
%
Farkas' Lemme follows from separation theorem among convex sets. Two
convex and closed sets that are disjoint - there is a hyperplane that
separates the points between them.

Why is the Farkas Lemma a separation theorem? $C=\left\{ Ax\big|x\ge0\right\} =$
the cone generated by the columns of $A$. The assumption says $b\notin C$.
\begin{lemma}[Projection Lemma]
 \label{lem:Projection}Let $C\subseteq\R^{m}$ be convex and closed.
\begin{enumerate}
\item $\forall b\in\R^{m}$ there is a unique $z\in C$ that is closest
to $b$. (i.e. $\min\limits _{z\in c}\norm{b-z}$) $z$ is called
the projection of $b$ on $c$.
\item $z$ is characterized by the following condition: For every $x\in C$
there holds $\left\langle x-z,b-z\right\rangle \le0$.
\end{enumerate}
\end{lemma}

\begin{proof}
$ $
\begin{enumerate}
\item Existence follows by compactness. Uniqueness $\then$ convexity
\item First want to show: if $\forall x\in C$ $\left\langle x-z,b-z\right\rangle \le0$
$\then$$z$ is the minimizer.
\begin{align*}
\norm{x-b}_{2}^{2} & =\norm{\left(x-z\right)+\left(z-b\right)}_{2}^{2}=\\
 & =\norm{x-z}^{2}+\norm{z-b}^{2}+2\left\langle x-z,z-b\right\rangle \ge\\
 & \ge\norm{z-b}^{2}+2\underbrace{\left\langle x-z,z-b\right\rangle }_{\ge0}
\end{align*}
 If $z$ is such that $\left\langle x-z,b-z\right\rangle \le0$ then
$z$ is the minimizer. (We compare $\norm{x-b}$ with $\norm{z-b}$
and we get $\norm{x-b}\ge\norm{z-b}+nonNegativeItem$. This means
that $z$ is the unique minimizer). The proof stops here without being
completed.
\end{enumerate}
\end{proof}
Duality. ``We can do with duality a lot of wonderful things''. ``Everything
that is implied by a finite system of linear equations and inequalities
can be derived in the standard way: Any linear combinations of equations
and any non negative linear combination of inequalities''.
\begin{itemize}
\item LP: 
\begin{align*}
\max\left\langle c,x\right\rangle \\
 & Ax\le b\\
 & x\ge0
\end{align*}
\item and $DLP$:
\begin{align*}
\min\left\langle b,y\right\rangle \\
 & yA\ge c\\
 & y\ge0
\end{align*}
\end{itemize}
Another way to represent $DLP$ and $LP$: In the upcoming theorem
we will discuss about this:
\begin{itemize}
\item $LP_{2}$:
\begin{align*}
\max\left\langle c,x\right\rangle \\
 & Ax=b\\
 & x\ge0
\end{align*}
\item $DLP_{2}$:
\begin{align*}
\min\left\langle b,y\right\rangle \\
 & yA\ge c
\end{align*}
There is no requirement for $y\ge0$%
\begin{insight}
Why is the symmetry broken in the second representation? We asked
in the second representation whether $\left\langle c,x\right\rangle =t$
is feasible. by Farkas lemma it is unsatisfiable if it can be refuted
in a specific way.%
\end{insight}
\begin{insight}
if feasibilty for the LP fails, then we can combine them however we
want?%
\end{insight}
\end{itemize}
\begin{thm}
For every feasible $y$ (for DLP) and any feasible $x$ (for LP) there
holds $\left\langle c,x\right\rangle \le\left\langle b,y\right\rangle $.
If the LP is feasible and bounded then the optima of the LP and the
DLP are equal.
\end{thm}

\begin{insight}
If we wish to show these theorem in prespective of NP, NP intersects
co-NP.%
\end{insight}
\begin{insight}

\subsubsection{Farkas$\protect\then$LP duality}
\begin{lemma}[Farkas]
 The system $Ax=b,x\ge0$ has no solution (infeasible) if and only
if $\exists y$ s.t. $yA\ge0$, $\left\langle y,b\right\rangle \le0$.
\end{lemma}

Because $yAx=\left\langle y,b\right\rangle <0$ and $yA\ge0$ and
$Ax\ge0$.%
\begin{insight}
\begin{proof}[Strong LP Duality \thmref{DualLP}]
 This is an incomplete proof,

Rather than solving $LP_{2}$ we deal with the question: is the system
\begin{align*}
\left\langle c,x\right\rangle  & =\lambda\\
 & Ax=b\\
 & x\ge0
\end{align*}
$\begin{pmatrix}A\\
c
\end{pmatrix}x=\begin{pmatrix}b\\
\lambda
\end{pmatrix}$ feasible? By Farkas it's infeasible if and only if there exists a
$y$ such that $\tilde{y}\begin{pmatrix}A\\
c
\end{pmatrix}\ge0$, $\left\langle \tilde{y},\begin{pmatrix}b\\
\lambda
\end{pmatrix}\right\rangle <0$ where $\tilde{y}=\left(y,-1\right)$. %
\begin{insight}
the last coordinate of $\tilde{y}$ is non negative since otherwise
the promal linear program is infeasible. We can assume that the last
coordinate is $-1$ otherwise we can simply devide $\tilde{y}$ with
$-\tilde{y}_{m+1}$ and denote $y=\left(\tilde{y_{1}}\dots\tilde{y_{m}}\right)$.
In this case $\left\langle \tilde{y},\begin{pmatrix}b\\
\lambda
\end{pmatrix}\right\rangle =\left\langle y,b\right\rangle -\lambda$ and the the inequality $\left\langle \tilde{y},\begin{pmatrix}b\\
\lambda
\end{pmatrix}\right\rangle <0$ can be written as $\left\langle y,b\right\rangle <\lambda$.%
\end{insight}
\[
\tilde{y}\begin{pmatrix}A\\
c
\end{pmatrix}\ge0\iff\underline{yA-c\ge0}
\]
{} 
\[
\underline{yb-\lambda<}0\iff\left\langle \tilde{y},\begin{pmatrix}b\\
\lambda
\end{pmatrix}\right\rangle <0
\]
\begin{insight}
Notice that we got that it is infeasible when $y^{T}A\ge c$ and $yb<\lambda=\left\langle c,x\right\rangle $%
\end{insight}
. %
\begin{insight}
there is such a $y$, multipling that $y$ with a negative doesn't
break our inner product negativity.%
\end{insight}
{} %
\begin{insight}
What we proved is that farkas lemma implies the strong duality. We
now showed that the critical $\lambda$ that moves the $lp$ to not
be possible is bigger then $yb$.%
\end{insight}
\begin{insight}
I think this helps because we showed that there is no solution when
we surpass the dual.%
\end{insight}
{} We saw that Farkas implies strong duality.%
\begin{insight}
because $\lambda$ is greater then the maximum of the LP if and only
if there is a vector $y$ such that $y^{T}A\ge c$ and $\left\langle y,b\right\rangle <b$.
If we relax the constraint $\left\langle y,b\right\rangle <\lambda$
to be $\left\langle y,b\right\rangle \le\lambda$ then there will
be a single critical value of $\lambda$ for which both systems have
a solution. This critical $\lambda$ is the maximum of the original
LP as well as the minimum of the dual linear program.%
\end{insight}
\end{proof}
\end{insight}
\end{insight}

\begin{lemma}[Farkas]
 The system $Ax=b,x\ge0$ has no solution (infeasible) if and only
if $\exists y$ s.t. $yA\ge0$, $\left\langle y,b\right\rangle \le0$.
\end{lemma}

\begin{proof}
First: If such $y\in\R^{m}$ exists, let us multiply by it on the
left: $yAx=\left\langle b,y\right\rangle <0$ and we know $yAx=\left\langle yA,x\right\rangle \ge0$
because $yA\ge0$ and $x\ge0$. %
\begin{insight}
Page 54 AllLectures%
\end{insight}
{} Second part:
\begin{remark}
recall infeasibility of 
\begin{align*}
\max\left\langle c,x\right\rangle \\
 & Ax=b\\
 & x\ge0
\end{align*}
means $b\notin C$ where $C$ is the cone generated by the columns
of $A$. What the lemma says is that there is a hyperplane through
the origin that separates $b$ from $C$. $H=\left\{ x\big|\left\langle x,y\right\rangle =0\right\} $.
$\left\langle y,b\right\rangle <0\iff b\in H^{-}$, $yA\ge0\iff\forall i\left\langle y,a_{i}\right\rangle \ge0$
$\iff$ $\forall ia_{i}\in H^{+}$. The hyperplane $H$ strongly separates
$b$ from the cone $C$. (we are using the separation theorem which
will be shown next)
\end{remark}

\end{proof}
%

\subsection{Separation theorem}

More generally, We have separation lemma:
\begin{lemma}[Separation lemma]
 Let $C\subseteq\R^{m}$ be closed and convex and $b\notin C$, Then
there is a hyperplane separating $b$ from $C$.
\end{lemma}

\begin{proof}
If $z\in C$ is \uline{the} closest point to $b$, then we can
take $H=\left\{ x\big|\left\langle x-z,b-z\right\rangle =0\right\} $.
\begin{insight}
$x-z$ is orthogonal to $b-z$ so $H$ is a family of lines. 
\[
\left\langle x,b-z\right\rangle =\left\langle z,b-z\right\rangle 
\]
\end{insight}
\begin{insight}
This hyperplane touch $C$ in point $z$.%
\end{insight}
{} It holds that $\left\langle b-z,b-z\right\rangle =\norm{b-z}^{2}>0$
since $b\notin C$ (which means $z\ne b$). On the other hand for
$x\in C$ it holds that $\left\langle x-z,b-z\right\rangle \le0$
(from the projection lemma \ref{lem:Projection} ). Which means $b\in H^{+}$
and $C\subseteq H^{-}$.
\end{proof}
\begin{insight}
Note that $\left\langle x-z,z-b\right\rangle $ is the opposite $\left\langle x-z,z-b\right\rangle $
(minus one multiplication)%
\end{insight}
\begin{insight}
$\norm{x-z}\ne0$ because $x\ne z$.%
\end{insight}

\begin{insight}
Projection $\then$ separation: If we let $x=b$, $\left\langle x-z,b-z\right\rangle =\norm{b-z}^{2}>0$$\then$
$b\in H^{+}$. By the projection lemma $\forall x\in C$ there holds
$\left\langle x-z,b-z\right\rangle \le0$ $\then$ $C\subseteq H^{-}$.%
\end{insight}

Final missing step for the Farkas lemma proof is: Why $H$ goes through
the origin? $H=\left\{ x\big|\left\langle x-z,b-z\right\rangle =0\right\} $.
We know that $b\in H^{+}$ (if $x=z\then0$ $x=0$ means $\left\langle x-z,b-z\right\rangle \le0$
\begin{insight}
Not sure why, todo%
\end{insight}


\subsubsection{The correct order}

Projection lemma $\then$ Separation lemma. Specialize $C=$ convex
cone $\then$ Farkas $\then$ Strong duality (The proof for strong
duality will be shown a bit later).

\section{The Ellipsoid algorithm}

There exists a bounded and convex body $P$ in $\R^{n}$. Problem:
decide whether $P=\emptyset$ or not. Assumptions about $P$: $P\subseteq B\left(0,R\right)$
(ball with radius $R$ around the origin ) if $P\ne\emptyset$$\then$
$P$ contains a ball of radius $r$.

We assume that these programs exist:
\begin{itemize}
\item Efficient Membership oracle: $x\in P$?
\item Efficient separation oracle: $H$ separating $x$ from $P$.
\end{itemize}
\begin{insight}
We assume that $R,r$ are known to us.%
\end{insight}

\begin{definition}
An ellipsoid is the image of the unit ball under affine transformation.
\end{definition}

$B\left(0,1\right)=\underbrace{\sum x_{i}^{2}}_{\left\langle x,x\right\rangle }\le1$

In other words an ellipsoid $E$ with \uline{center} at $c$ is
given by: $E=\left\{ x\big|\left(x-c\right)M\left(x-c\right)^{T}\le1\right\} $
where $M$ is a positive definite matrix.

We know by assumption $P$ is hiding somewhere inside a some enormous
ball. We can ask whether $0\in P$ if the answer is yes -> game over.
if the answer is no? we go to the separation oracle to separate $P$
from our origin.

step $k$ of the algorithm: We have an ellipsoid $E_{k}$ with center
$c_{k}$. we ask is $c_{k}\in P$ using our membership oracle, if
the answer is yes - game over $P\ne\emptyset$. Else we take a separating
hyperplane using our separation oracle. Previously we knew that $P\subseteq E_{k}$
now we know $P\subseteq E_{K}^{+}$ the half ellipsoid which we cut
using the hyperplane. Now we create a new ellipsoid which will contain
$E_{k}^{+}$.
\begin{definition}
a half-ellipsoid is the intersection of an ellipsoid $E$ with a half-space
$H^{+}$ through $c$ the center of $E$.
\end{definition}

\begin{proposition}
For every half-ellipsoid $E^{+}\subseteq\R^{n}$ there exist an ellipsoid
$E'$ such that $E^{+}\subset E'$ and $vol\left(E'\right)\le\left(1-\frac{1}{2n^{2}}\right)Vol\left(E\right)$
\begin{insight}
this isn't the exact constant $1-\frac{1}{2n^{2}}$ it might be different
but that's okay.%
\end{insight}
\end{proposition}

The algorithm keeps on going until we find membership or we reach
an ellipsoid which has a smaller volume then a ball of radius $r$
(because $P$ should contain a ball of radius $r$).

\begin{insight}
The idea is to keep on running the ellipsoid and we cut the volume
really fast. if we get to a volume which is really small - smaller
then $r$ we can conclude that $P$ can't be inside something which
it should hold (a ball of radius $r$.)%
\end{insight}
\begin{insight}
An affine transformation $x\to Ax+b$ multiplies all volumes by $\left|detA\right|$.
therefore it suffices to prove the proposition when $E=B\left(0,1\right)$
. $E^{+}=E\cap\left\{ x\big|x_{1}>0\right\} $. %
\begin{insight}
We apply a reverse transformation such that we move the ellipsoid
back to the unit ball and we can calculate it there and then put it
back.%
\end{insight}

\begin{insight}
Ayelet's 3 last pages%
\end{insight}

We want to use this ellipsoid as a tool to solve linear programming
- for feasibility. %
\begin{insight}
Complete from video, There is a missing segment here 15:37%
\end{insight}

\uline{Q}: is there a strongly polynomial algorithm for $LP$?
$LP$ is in $P$ (polynomial time) because ellipsoid can solve it
in polynomial time. we don't know if there is algorithm $LP$. This
is still an open question and we don't know this. in practice this
is okay. neat.

The Cone of PSD matrices:

Membership oracle: $M\in PSD$ ? we can do that easily.

Separation oracle: $\sum a_{ij}m_{ij}<0$ and $\sum a_{ij}q_{ij}\ge0$
$\forall Q\in PSD$.%
\end{insight}


\section{Last Lecture}
\begin{itemize}
\item We will talk about strong duality from Farkas.
\item We will talk about the ellipsoid algorithm once more.
\item We will show that determinant is volume of a matrix
\end{itemize}

\subsection{Farkas $\protect\then$ Strong LP duality}
\begin{thm}[Strong LP duality]
 Suppose that both LP and DLP are feasible, and define a non empty
polytope (=bounded set) $\then$ their optima are equal.
\begin{itemize}
\item LP:
\begin{align*}
\max\left\langle c,x\right\rangle \\
 & Ax=b\\
 & x\ge0
\end{align*}
\item DLP:
\begin{align*}
\min\left\langle b,y\right\rangle \\
 & yA\ge c
\end{align*}
\end{itemize}
\end{thm}

\begin{proof}
$ $

Easy part: under these assumptions $x,y$ feasible for LP/DLP respectively.
$\left\langle c,x\right\rangle \le\left\langle b,y\right\rangle $.
\ref{weak-LP-duality-proof}%
\begin{insight}
easy to show this. also we've seen it.%
\end{insight}

hard part: Let $y_{0}$ be an optimal solution for a DLP%
\begin{insight}
. if we can find an $x$ such it solve LP and the two solutions are
equal we are done%
\end{insight}
. Let us assume %
\begin{insight}
by contradiction?%
\end{insight}
{} that for every LP--feasible $x$, it holds that $\left\langle c,x\right\rangle <\left\langle b,y_{0}\right\rangle $.
This means that there is no $x\ge0$ such that $Ax=b$ and $\left\langle c,x\right\rangle =\left\langle b,y_{0}\right\rangle $,
in other words $\nexists x\ge0$ such that $\tilde{A}x=\tilde{b}$,
$\tilde{A}=\begin{pmatrix}A\\
c
\end{pmatrix},\tilde{b}=\begin{pmatrix}b\\
\left\langle y_{0},b\right\rangle 
\end{pmatrix}$. By Farkas, %
\begin{insight}
it isn't feasible if and only if there is such vector%
\end{insight}
{} there is a vector $w$ which shows that no such $x$ exists. i.e
$w$ such that $w\tilde{A}\ge0$ and $\left\langle w,\tilde{b}\right\rangle <0$.
Three cases to check: The last coordinate of $w$ is $\left\{ -1,0,1\right\} $
\begin{insight}
i assume we can normalize the vector to get such last coordinate.%
\end{insight}
\begin{insight}
(we normalize $w$ with a positive value such that the last entry
of $w$ is either $\left\{ 1,-1,0\right\} $)%
\end{insight}
:
\begin{itemize}
\item $w=\begin{pmatrix}z\\
-1
\end{pmatrix}$: $w\tilde{A}\ge0\iff zA\ge c$. This means that $\left\langle z,b\right\rangle <\left\langle y_{0},b\right\rangle $
(from $\left\langle w,\tilde{b}\right\rangle <0$ ). We found a $z$
that it is better then $y_{0}$.
\item If $w=\begin{pmatrix}z\\
1
\end{pmatrix}$ then $zA\ge-c$ (from $w\tilde{A}\ge0$) this means that $\left\langle z+y_{0},b\right\rangle =\left\langle z,b\right\rangle +\left\langle y_{0},b\right\rangle <0$.
($\left\langle \begin{pmatrix}z\\
1
\end{pmatrix},\begin{pmatrix}b\\
y_{0},b
\end{pmatrix}\right\rangle <0\iff\left\langle z,b\right\rangle +\left\langle y_{0},b\right\rangle <0$) 
\begin{claim*}
For some $\epsilon>0$ small enough: $\left(1+\epsilon\right)y_{0}+\epsilon z$
is a better solution than $y_{0}$ (for DLP).
\begin{proof}
Lets look at: 
\[
\left(\left(1+\epsilon\right)y_{0}+\epsilon z\right)A=y_{0}A+\underbrace{\epsilon\left(y_{0}+z\right)A}_{\ge0}\ge c
\]
 because $y_{0}A\ge c$ and $zA\ge-c$ so $\left(y_{0}+z\right)A\ge0$,
Now 
\[
\left\langle \left(1+\epsilon\right)y_{0}+\epsilon z,b\right\rangle <\left\langle y_{0},b\right\rangle 
\]
which means we can construct a vector which is a combination of $y_{0}$
and $z$ to get a better solution (smaller). We need to see that $\left\langle \epsilon y_{0}+\epsilon z,b\right\rangle <0$
to show that. \\

\begin{insight}
My take on trying to simplifying it: We know that $\left\langle z+y_{0},b\right\rangle =\left\langle z,b\right\rangle +\left\langle y_{0},b\right\rangle <0$
so $\left\langle z+y_{0},b\right\rangle <0\iff\epsilon\left\langle z+y_{0},b\right\rangle <0$
for $\epsilon>0$. Now 
\[
\left\langle \left(1+\epsilon\right)y_{0}+\epsilon z,b\right\rangle =\left\langle y_{0}+\epsilon z+\epsilon y_{0},b\right\rangle =\epsilon\underbrace{\left\langle z+y_{0},b\right\rangle }_{<0}+\underbrace{\left\langle y_{0},b\right\rangle }_{<0}
\]
 which means we improve the solution $y_{0}$ because $\left(1+\epsilon\right)y_{0}+\epsilon z$
is smaller then simply $y_{0},b$ .%
\end{insight}
\end{proof}
\end{claim*}
\item what happens when the last coordinate is $0$? $w=\begin{pmatrix}z\\
0
\end{pmatrix}$ $zA\ge0,\left\langle z,b\right\rangle <0$ this means that the DLP
is unbounded. Explanation: Assume $m$ is positive, $\left(y+mz\right)A\ge c$,
$\left\langle y+mz,b\right\rangle <\left\langle y_{0},b\right\rangle $
which is a better solution. %
\begin{insight}
We can make $m$ as big as we want.. 
\[
\left\langle \begin{pmatrix}mz\\
0
\end{pmatrix},\begin{pmatrix}b\\
\left\langle y_{0},b\right\rangle 
\end{pmatrix}\right\rangle =\left\langle mz,b\right\rangle <0
\]
\end{insight}
\end{itemize}
\begin{insight}
We showed a contradiction. %
\end{insight}

\end{proof}
%

\subsection{Back to the ellipsoid algorithm}
\begin{enumerate}
\item An ellipsoid is the image of a ball ($\ell_{2}$ ball) under an affine
transformation.
\item We consider %
\begin{insight}
David: ``should be bounded'', David later: ``Welp it is bounded,
it is surrounded by a ball''%
\end{insight}
{} a convex body (fully dimensional) $P\subseteq\R^{n}$. $P\subseteq B\left(0,R\right)$,
guarantee:
\begin{enumerate}
\item if $p\ne0$ $\then$ it contains a ball of radius $r$.
\item Oracles: We can answer efficiently the queries:
\begin{enumerate}
\item Membership: $x\overset{?}{\in}P$.
\item Separation: $H$ that separates $x$ from $P$.
\end{enumerate}
\end{enumerate}
\item We consider ellipsoid $E_{k}$ with centers $x_{k}$ $E_{0}=B\left(0,R\right)$
, $x_{0}=0$
\item Step: Query $x_{k}\overset{?}{\in}P$, In case the answer is yes $\then P\ne\emptyset$.
Else : Let $\tilde{H}$ be a separating hyperplane, $H||\tilde{H}$
goes through $x_{k}$. %
\begin{insight}
We took the parallel hyperplane that goes through $x_{k}$..%
\end{insight}
$E_{k}^{+}=E_{k}\cap H^{+}$.
\end{enumerate}
\begin{insight}
We never work with the ellipsoid we work with matrices over the ball..%
\end{insight}

\begin{lemma}
If $E\subseteq\R^{n}$ is an ellipsoid with center $x$, and $H$
as hyperplane through $x$ Let $E^{+}=E\cap H^{+}$ $\then$ $\exists$
Ellipsoid $E'\supset E^{+}$ where $vol\left(E'\right)<\left(1-\frac{1}{2n^{2}}\right)vol\left(E\right)$.
\end{lemma}

This means that the ellipsoids we find shrink exponentially.

\subsubsection{personal note}

The ellipsoid can be used in order to find a feasible solution to
the LP problem: Once we find feasibilty we know from strong duality
that the optimal solution is achievable when the DLP and LP solutions
intersect. Lets find the feasible solution for the following problem:
$Ax=b,x\ge0$ and $y^{T}A\ge c^{T}$, $c^{T}x=y^{T}b$. $x,y$ that
are feasible also uphold the best solution for the LP and DLP problem. 

Using the ellipsoid method we should be able to find a feasible solution
(a member of the group) to a linear programming problem using a separation
and membership oracle. 

With that knowledge about the ellipsoid method we advance to the next
subject.

\section{semi-definite programming}

Having a method of finding feasible solutions we face a new type of
optimization problems, lets dive in:

\begin{insight}
We are in $n^{2}$ dim space, in there there is a cone of the positive
semi definite matrices. What do we do here? What do we want to check?%
\end{insight}

$PSD_{n}:=$ Cone of (symmetric) $n\times n$ positive semi-definite
matrices.

This cone has good separation and membership oracles. 
\begin{itemize}
\item Membership, Given $n\times n$ matrix $A$ we can ask whether $A\overset{?}{\in}PSD_{n}$
and find a solution fast by finding all non-negative eigenvalues.
\item Separation, Suppose that $A\notin PSD_{n}$ $\then$ $\exists x$
such that $xAx^{T}=\sum\overbrace{x_{i}x_{j}}^{\text{coefficients}}\cdot\overbrace{a_{ij}}^{\text{variable}}<0$%
\begin{insight}
$xAx=\sum\limits _{i,j}x_{i}x_{j}a_{ij}<0$ so this is negative. The
variables are $a_{i,j}$ and $x_{i},x_{j}$ are the coefficients for
the separator.%
\end{insight}
{} where $\forall B\in PSD_{n}$, $xBx^{T}\ge0$. $x$ determine the
coefficients of $A$, and thus define a separator. For any point $C$
in the space, if $xCx=\sum\left(x_{i}x_{j}\right)\cdot c_{ij}\ge0$
(similar to $\left\langle x,c\right\rangle \ge0$) then $C\in PSD_{n}$
else $xCx<0$, $C\notin PSD$. Therefore $x$ defines a separator.
(How can we find such an $x$? We find a negative eigenvalue for $A$)
\end{itemize}
consequences of Ellipsoid Algorithm:
\begin{enumerate}
\item Efficiently decide whether a finite list of linear inequalities is
feasible. (find feasible solution for LP in polynomial time)
\item In addition require that a matrix made up of our variables in $PSD$.
\begin{insight}
$X$ is a matrix in $PSD_{n}$list of linear equations and inequalities
in $x_{ij}$ optimize some fin function in the $x_{ij}$.%
\end{insight}
{} In otherwords we can now solve questions as follows: given $X\in PSD$
and a list of linear equations and inequalities in $x_{i,j}$ , optimize
some linear function in the $x_{i,j}$.
\end{enumerate}
Lets see what we can do with such a tool


\subsection{approximation of Max cut}

Given a graph $G$ find the bi-partition $V\left(G\right)=A\cup B$
such that $e\left(A,B\right)$ is maximized (NP--Hard).

\subsubsection{Goemans--williamson algorithm}

Gives a $0.868$ approximation to max cut using semi-definite programming.
Max cut for $G=\left(V,E\right),V=\left\{ 1\dots n\right\} $ can
be defined as

\[
maxCut\left(G\right)=\max_{x:v\to\left\{ -1,1\right\} }\sum_{i,j\in E}\frac{1-x_{i}x_{j}}{2}
\]

Lets map the vertices into unit vectors. 
\[
\text{v-maxCut}\left(G\right)=\max_{\begin{array}{c}
x_{1}\dots x_{n}\\
\norm{x_{i}}=1
\end{array}}\sum_{i,j\in E}\frac{1-\left\langle x_{i},x_{j}\right\rangle }{2}
\]
 %
\begin{insight}
we are allowed to use any unit vectors that we want, So $v-maxCut\left(G\right)=\max_{\begin{array}{c}
x_{1}\dots x_{n}\\
\norm{x_{i}}=1
\end{array}}\sum_{i,j\in E}\frac{1-\left\langle x_{i},x_{j}\right\rangle }{2}$ is bigger then $\max_{x:v\to\left\{ -1,1\right\} }\sum_{i,j\in E}\frac{1-x_{i}x_{j}}{2}$%
\end{insight}
{} Obviously $\text{v-maxCut}\left(G\right)\ge\text{maxCut}\left(G\right)$
\begin{thm}
$ $
\begin{enumerate}
\item $1.13\cdot\text{maxCut}\left(G\right)\ge\text{v-maxCut}\left(G\right)$.
\item $\text{v-maxCut}\left(G\right)$ can be efficiently computed.(Using
semi-definite programming)
\end{enumerate}
\end{thm}

we construct a matrix $A=\begin{pmatrix}-x_{1}-\\
\vdots\\
-x_{n}-
\end{pmatrix},Y=AA^{T}$ so $Y\in PSD$. Each $x_{i}$ vector is a unit vectors. so $\forall i\quad y_{ii}=1$.
What we want is 
\[
\max\sum_{i,j\in E}\frac{1-y_{ij}}{2}
\]
 And we can do this using semi-definite programming. %
\begin{insight}
meaning we maximize $y_{ij}$ in a PSD matrix%
\end{insight}
. %
\begin{insight}
We want to maximize this because $y_{i,j}=\left\langle x_{i},x_{j}\right\rangle $%
\end{insight}

Two issues remain:
\begin{enumerate}
\item why is $\text{v-maxCut}\left(G\right)\le1.13\text{maxCut}\left(G\right)$
?
\item How do we actually partition $V\left(G\right)$?
\end{enumerate}
Given $Y$ we can recover the $x's$. Actually given optimal $x_{i}$'s
that solve this, we can efficiently find a partition that is $\ge0.868\cdot maxCut\left(G\right)$.

step two of the algorithm after we found $Y$. Pick a random hyperplane
through the origin, and partition $v$ accordingly. (We select a random
vector and take an inner product with all the vectors $x_{i}$'s)
People we're able to remove the randomness but we're not going to
get into that.

\begin{insight}
if $\left\langle x_{i},x\right\rangle \le0$ it's in partition $A$,
else $\left\langle x_{i},x\right\rangle >0$ its in partition $B$.%
\end{insight}

Suppose that $x_{i},x_{j}$ form an angle of $\theta\iff\left\langle x_{i},x_{j}\right\rangle =\cos\theta$.
What is the probability that a random hyperplane through the origin
separates $x_{i},x_{j}$, The answer is $\frac{\theta}{\pi}$ %
\begin{insight}
$\left(\frac{\theta}{180}\right)$%
\end{insight}
\begin{insight}
plot the two points and the angle between them and the origin, Now
we want to ask what is the chance that we get a line in between them%
\end{insight}
. $\then$ $\E{\text{size of the cut we create}}=\sum\limits _{i,j\in E}\frac{\theta_{i,j}}{\pi}$.
\begin{claim}
We create a cut that is $\ge0.868\text{maxCut}\left(G\right)$. Actually
we show that our cut (in expectation) is 
\[
\ge0.868\text{v-maxCut}\left(G\right)\ge0.868\text{maxCut}\left(G\right)
\]
 Our claim is that $\sum\frac{\theta_{i,j}}{\pi}\ge0.868\sum\frac{1-\left\langle x_{i},x_{j}\right\rangle }{2}$.
Actually: $\forall\pi\ge\theta\ge0\quad$ $\frac{\theta}{\pi}\ge0.868\left(\frac{1-\cos\theta}{2}\right)$.
\end{claim}


\subsection{Determinant is volume}
\begin{definition}
Let $x_{1}\dots x_{m}$ be vectors in $\R^{n}$ the parallelepiped
$\left(Makbilon\right)$ that they generate is $P=\left\{ \sum_{i=1}^{m}t_{i}x_{i}\big|\forall i\quad1\ge t\ge0\right\} $
\begin{insight}
Looks like diamond, or the sim's diamond in $\R^{3}$%
\end{insight}
\begin{insight}
in convex hull we demand that the $t_{i}$s sum into $1$.%
\end{insight}

Why do we consider it? This is the image of the unit cube under a
linear map given by the matrix whose rows are $x_{1}\dots x_{m}$
\end{definition}

\begin{thm}
The volume $vol\left(P\right)$ of the parallelepiped formed by $x_{1}\dots x_{m}\in\R^{n}$
is given by $vol^{2}\left(P\right)=\det AA^{T}$ %
\begin{insight}
$\det A=\det A^{T}$ and $\det AB=\det A\cdot\det B$%
\end{insight}
where $A$ is the matrix with rows $x_{1}\dots x_{m}$.
\end{thm}

\begin{proof}
By induction on $m$. For $m=1$ the matrix is just a vector, and
$\left\langle x,x\right\rangle =\norm x^{2}$.

$A=\begin{pmatrix}-x_{1}-\\
-x_{2}-\\
\vdots\\
-x_{m}-
\end{pmatrix}$ all we need to do is to express $x_{1}=a+b$ where $a\perp span\left\{ x_{2}\dots x_{m}\right\} $
and $b\in span\left\{ x_{2}\dots x_{m}\right\} $ as a result the
height is exactly $a$. $\begin{pmatrix}-a+b-\\
-x_{2}-\\
\vdots\\
-x_{m}-
\end{pmatrix}=\left(\begin{array}{c}
a+b\\
\hline M
\end{array}\right)$. 
\begin{align*}
\det AA^{T} & =\det\left(\begin{array}{cc}
\norm{a+b}^{2} & \left(a+b\right)M\\
\left(a+b\right)M^{T} & MM^{T}
\end{array}\right)=\\
 & =\det\left(\begin{array}{cc}
\norm a^{2} & 0\\
0 & MM^{T}
\end{array}\right)\overset{I.H}{=}\norm a^{2}\cdot vol^{2}\left(x_{2},\dots x_{m}\right)
\end{align*}

\end{proof}
%
Thanks for reading and a special thanks for everyone who contributed.
\end{document}
